\documentclass[11pt]{article}
\usepackage{amsmath}
\usepackage{amssymb}
% Change "review" to "final" to generate the final (sometimes called camera-ready) version.
% Change to "preprint" to generate a non-anonymous version with page numbers.

\usepackage[review]{acl} 
% \usepackage[final]{acl}
% \usepackage[preprint]{acl}

% Standard package includes
\usepackage{times}
\usepackage{latexsym}

% % ask 1
\usepackage{subcaption}

\usepackage[T1]{fontenc}
% For Vietnamese characters
% \usepackage[T5]{fontenc}
% See https://www.latex-project.org/help/documentation/encguide.pdf for other character sets

% This assumes your files are encoded as UTF8
\usepackage[utf8]{inputenc}

% This is not strictly necessary, and may be commented out,
% but it will improve the layout of the manuscript,
% and will typically save some space.
\usepackage{microtype}

% This is also not strictly necessary, and may be commented out.
% However, it will improve the aesthetics of text in
% the typewriter font.
\usepackage{inconsolata}

%Including images in your LaTeX document requires adding
%additional package(s)
\usepackage{graphicx}
\usepackage{booktabs}  % For professional table formatting (\toprule, \midrule, \bottomrule)
\usepackage{tabularx}  % For responsive table column widths
\usepackage{array}     % For advanced column formatting
\usepackage{multirow}  % For multi-row cells in tables
\usepackage{xcolor}    % For colored text in tables      

% \usepackage{comment}

% If the title and author information does not fit in the area allocated, uncomment the following
%
%\setlength\titlebox{<dim>}
%
% and set <dim> to something 5cm or larger.th

\title{Towards a Mechanistic Understanding of Propositional \\ Logical Reasoning in Large Language Models}

% \author{
%   \textbf{Danchun Chen\textsuperscript{1\thanks{Equal Contribution.}
%   % \thanks{Work done during research assistantship at PKU \textit{PILLAR}.}
%   },
%   \textbf{Qiyao Yan\textsuperscript{1\footnotemark[1] }},
%   % \textbf{Liangming Pan\textsuperscript{2}}}
%   \textbf{Liangming Pan\textsuperscript{1\thanks{Corresponding Author.}}}}
% \\
%   % \textsuperscript{1}Individual Researcher,
%   \textsuperscript{1}
%   MOE Key Lab of Computational Linguistics, Peking University
% \\
%   \texttt{dcchen.nexus@outlook.com},
%   \texttt{2300013233@stu.pku.edu.cn},
%   \texttt{liangmingpan@pku.edu.cn}
% }

\author{
  \textbf{Danchun Chen\thanks{Equal Contribution.}}, 
  \textbf{Qiyao Yan\footnotemark[1]},
  \textbf{Liangming Pan\thanks{Corresponding Author.}}
\\
  % 移除了 \textsuperscript{1}
  MOE Key Lab of Computational Linguistics, Peking University
\\
  \texttt{dcchen.nexus@outlook.com},
  \texttt{yqy.1@stu.pku.edu.cn},
  \texttt{liangmingpan@pku.edu.cn}
}


% 自定义符号
% \author{
%   \textbf{Danchun Chen}\textsuperscript{1}\footnotemark[1]\footnotemark[3],
%   \textbf{Qiyao Yan}\textsuperscript{2}\footnotemark[1],
%   \textbf{Liangming Pan}\textsuperscript{2}\footnotemark[2]
% \\
%   \textsuperscript{1}Independent Researcher,
%   \textsuperscript{2}Peking University
% \\
%   \texttt{dcchen.nexus@outlook.com},
%   \texttt{2300013233@stu.pku.edu.cn},
%   \texttt{liangmingpan@pku.edu.cn}
% }



\begin{document}
\maketitle  
% % Author footnotes using fnsymbol
% {
%   \renewcommand{\thefootnote}{\fnsymbol{footnote}}
%   \footnotetext[1]{Equal contribution.}
%   \footnotetext[2]{Corresponding author.}
%   \footnotetext[3]{Work done during research assistantship at PKU PILLAR.}
% }
\begin{abstract}
% While Large Language Models (LLMs) demonstrate impressive capabilities in logical reasoning, the internal computations driving these behaviors remain largely opaque. Do LLMs merely parrot statistical patterns, or do they implement algorithmic circuits akin to formal deduction? In this work, we investigate the mechanistic underpinnings of propositional logical reasoning within the \textit{Qwen3} model family \textit{(8B and 14B)}. By applying causal mediation analysis on a custom dataset covering diverse logical rules (\textit{e.g.}, De Morgan's Laws, Distributivity) across one-hop and two-hop reasoning tasks, we uncover a structured ``computational architecture" underlying reasoning. We identify three core phenomena: Information Transmission (semantic information aggregates at boundary tokens), Specialized Attention Heads (attention heads exhibit stable functional specialization), and Fact Retrospection (factual information is persistently re-accessed). Notably, these mechanisms generalize across model scales and logical rule types, suggesting that LLMs employ specific computational strategies to solve logical problems.
Understanding how Large Language Models (LLMs) perform logical reasoning internally remains a fundamental challenge. While prior mechanistic studies focus on identifying task-specific circuits, they leave open the question of what computational strategies LLMs employ for propositional reasoning. We address this gap through comprehensive analysis of 
% TODO: 加入更多模型的表述
\textit{Qwen3} (8B and 14B) on \textit{PropLogic-MI}, a controlled dataset spanning 11 propositional logic rule categories across one-hop and two-hop reasoning. Rather than asking ``which components are necessary,'' we ask ``how does the model organize computation?'' Our analysis reveals a coherent computational architecture comprising four interlocking mechanisms: Staged Computation (layer-wise processing phases), Information Transmission 
(information flow aggregation at boundary tokens),
% (information flow across stages),
Fact Retrospection (persistent re-access of source facts), and Specialized Attention Heads (functionally distinct head types). These mechanisms generalize across model scales, rule types, and reasoning depths, providing mechanistic evidence that LLMs employ structured computational strategies for logical reasoning. 
% Code and data are available at
\end{abstract}

% 议论，故事线（逻辑严密）
% 缺少说服力，1. CRGP. context(任务背景，关注的方向) ; 2. Review(related work, previous work 大概内容, existing works focus on ..., some works in logical reasoning) ; 3. Gap(先前的工作不足之处，什么是之前的工作没有做的，而某些东西是重要的；是否存在challenge？技术上的难度？要体现出工作的重要性，让别人觉得该工作的必要性；)；4. Proposal(讲自己要做的东西， to address this limitations...， 需要与前面的内容相联系， 与解决Gap相对应，不要过于深入技术细节，放到后面的篇幅中阐述，核心思路阐述）；5. Findings（核心的内容结论总结，最重要的发现）；6. Contributions(in summary).
% Introduction（通俗易懂，所有人）；related work（稍微深入，本科生）；Methods(正式、深入，研究生，某一两个部分特别深入细节，体现深度)；
\section{Introduction}

% The advent of transformer-based Large Language Models (LLMs) has revolutionized natural language processing, enabling models to perform complex tasks ranging from open-ended generation to multi-step problem solving \citep{openai2024openaio1card, deepseekai2025deepseekr1incentivizingreasoningcapability, comanici2025gemini25pushingfrontier}. Particularly noteable is their emergent ability to handle logical reasoning tasks, achieving over 85\% accuracy on rule-based benchmarks like PrOntoQA \citep{saparov2023languagemodelsgreedy} and ProofWriter \citep{tafjord2021proofwritergeneratingimplications}. Despite this empirical success, the effective mechanisms behind these capabilities remain a mystery. The field faces a central debate: are these models simply sophisticated pattern matchers relying on surface statistics \citep{bender2021stochasticparrots, dziri2023faithfatelimitstransformers, zhang2022paradoxlearningreasondata}, or have they internalized robust algorithms that mirror logical deduction \citep{li2023emergent, nanda2023progress}? To address this, the emerging field of Mechanistic Interpretability (MI) \citep{rai2024practicalreview, wang2022interpretabilitywildcircuitindirect, meng2023locatingeditingfactualassociations, olsson2022incontextlearninginductionheads, elhage2021mathematicalframeworktransformer} has begun to reverse-engineer transformers to causally uncover the mechanisms underlying model behaviors. 
% Recent efforts have applied MI to uncover the mechanisms behind LLM reasoning, including arithmetic \citep{kantamneni2025languagemodelsuse, stolfo2023mechanisticinterpretationarithmeticreasoning, yu2024interpretingarithmeticmechanismcomparative}, syllogistic inference \citep{kim2024reasoningcircuitslanguage} and combination tasks \citep{ye2025transformerslearnimplicit, wang2024languagemodelscan}.

%%% data: on LogicBench \citep{parmar2024logicbench}, models achieve 86\% on propositional logic but degrade to 74\% on non-monotonic reasoning, and continue to struggle with negations and complex inference chains \citep{saparov2023languagemodelsgreedy, dziri2023faithfatelimitstransformers}.

The capability of Large Language Models (LLMs) to execute complex logical reasoning has improved significantly \citep{cheng2025empoweringllmslogicalreasoning, openai2024openaio1card, deepseekai2025deepseekr1incentivizingreasoningcapability}. State-of-the-art evaluations demonstrated this trend: \textit{GPT-5.2} achieves 92\% accuracy on \textit{GPQA-Diamond}, while \textit{Gemini 3 Pro} reaches 76\% on \textit{SimpleBench}. 
% However, this empirical performance contrasts with the opacity of the underlying computational mechanisms. 
However, the underlying computational mechanisms behind such high empirical performance remain poorly understood. 
% Furthermore, systematic evaluations on propositional logic reveal persistent gaps: models achieve 86\% on LogicBench \citep{parmar2024logicbench} but degrade to 74\% on non-monotonic reasoning, and continue to struggle with negations and complex multi-step inference \citep{saparov2023languagemodelsgreedy, dziri2023faithfatelimitstransformers}. 
% The field thus faces a central debate: 
This gap in understanding raises a fundamental question: are these models simply sophisticated pattern matchers relying on surface statistics \citep{bender2021stochasticparrots, dziri2023faithfatelimitstransformers, zhang2022paradoxlearningreasondata}, or have they internalized robust algorithms that mirror logical deduction \citep{li2023emergent, nanda2023progress}? 
To address this ambiguity, the emerging field of \textit{Mechanistic Interpretability} (MI) 
% \citep{rai2024practicalreview, wang2022interpretabilitywildcircuitindirect, meng2023locatingeditingfactualassociations, olsson2022incontextlearninginductionheads, elhage2021mathematicalframeworktransformer} 
\citep{sharkey2025openproblemsmechanisticinterpretability}
has begun to reverse-engineer transformers to causally uncover the mechanisms underlying model behaviors. Recent efforts have applied MI to reasoning domains including arithmetic \citep{kantamneni2025languagemodelsuse, stolfo2023mechanisticinterpretationarithmeticreasoning, yu2024interpretingarithmeticmechanismcomparative}, syllogistic inference \citep{kim2024reasoningcircuitslanguage}, and multi-hop composition \citep{ye2025transformerslearnimplicit, wang2024languagemodelscan}. Their findings suggest that LLMs may employ structured, interpretable strategies when processing logical reasoning, rather than relying solely on opaque pattern matching.


% The capability of Large Language Models (LLMs) to execute complex logical reasoning has reached unprecedented levels, marking a fundamental shift from statistical pattern matching to what appears to be genuine deductive proficiency. Recent comprehensive surveys on reasoning capabilities demonstrate that state-of-the-art models now consistently achieve remarkable performance on general reasoning benchmarks—such as GSM8K, MATH, and Big-Bench Hard—with leading models often exceeding 90\% accuracy on tasks requiring multi-step deduction \citep{huang2023reasoningsurvey, yu2023natureofreasoning, openai2024openaio1card, deepseekai2025deepseekr1incentivizingreasoningcapability}. Yet, this empirical success stands in stark contrast to the opacity of the underlying mechanisms. The field faces a central epistemic debate: do these models merely exploit sophisticated surface-level correlations \citep{bender2021stochasticparrots, dziri2023faithfatelimitstransformers}, or have they spontaneously internalized robust algorithmic circuits that mirror formal logical operations \citep{li2023emergent, nanda2023progress}? To address this ambiguity, the emerging field of Mechanistic Interpretability (MI) \citep{rai2024practicalreview, wang2022interpretabilitywildcircuitindirect, meng2023locatingeditingfactualassociations, olsson2022incontextlearninginductionheads, elhage2021mathematicalframeworktransformer} has begun to reverse-engineer transformers to causally uncover the mechanisms underlying model behaviors. 
% Recent efforts have applied MI to uncover the mechanisms behind LLM reasoning, including arithmetic \citep{kantamneni2025languagemodelsuse, stolfo2023mechanisticinterpretationarithmeticreasoning, yu2024interpretingarithmeticmechanismcomparative}, syllogistic inference \citep{kim2024reasoningcircuitslanguage} and combination tasks \citep{ye2025transformerslearnimplicit, wang2024languagemodelscan}.

% In arithmetic, researchers have uncovered how transformers encode numerical representations \citep{kantamneni2025languagemodelsuse} and perform multi-step calculations \citep{stolfo2023mechanisticinterpretationarithmeticreasoning, yu2024interpretingarithmeticmechanismcomparative}. For syllogistic reasoning, \citet{kim2024reasoningcircuitslanguage} identified circuits implementing logical rules. In graph-based tasks, \citet{saparov2025transformerscanlearn} analyzed connectivity reasoning mechanisms.
% Closely related are studies on multi-hop reasoning, where models must chain multiple facts to derive conclusions \citep{ye2025transformerslearnimplicit, wang2024languagemodelscan}. Implicit reasoning research investigates how models reason without explicit intermediate steps—\citet{wang2024grokkedtransformersimplicit} show that grokking enables implicit multi-hop reasoning, while \citet{lin2025implicitreasoningtransformers} reveal that models often rely on shortcuts rather than systematic deduction. Chain-of-thought mechanistic studies \citep{dutta2024thinkstepbystep} further reveal how models implement sequential reasoning when explicitly prompted.

% However, \textbf{propositional logical reasoning remains largely unexplored} in mechanistic interpretability research. Only \citet{aimpliesb} provide circuit-level analysis, but their study is limited to a single logical operator structure (implication chains) with simplified length-2 problems. Broader logical rules—De Morgan's Laws, Distributivity, Commutativity—and their underlying mechanisms in production-scale LLMs remain unknown. There remains a significant gap in understanding how LLMs orchestrate internal components to solve comprehensive propositional logical reasoning tasks across diverse rule types and multi-hop scenarios. Furthermore, it is unclear whether the mechanisms identified in smaller ``toy" models scale to dense, production-grade LLMs like the Qwen series, and how these mechanisms adapt when the reasoning depth increases from single-step to multi-hop inference.
% However, \textbf{propositional logical reasoning}, a cornerstone of systematic reasoning, remains largely unexplored in MI. While \citet{aimpliesb} provides an initial circuit-level analysis on simple implication chains, the mechanisms governing broader logical rules (\textit{e.g.}, \textit{De Morgan's Laws}, \textit{Distributivity}) in production-scale LLMs remain obscure. This gap is further exacerbated by the lack of structural control in existing benchmarks, which hinders the fine-grained causal interventions necessary to distinguish true algorithmic execution from statistical shortcuts.
Propositional logical reasoning, \textit{i.e.}, the capacity to manipulate boolean operators and chain inference rules, represents a cornerstone of formal deduction that underpins more complex reasoning tasks. Despite its foundational role, this domain remains largely unexplored in MI. Behavioral studies reveal that model accuracy degrades significantly with increased proof depth \citep{saparov2023languagemodelsgreedy, dziri2023faithfatelimitstransformers}, yet the internal computational mechanisms responsible for these failures remain opaque. A prior mechanistic work, \citet{aimpliesb}, provides circuit-level analysis identifying four attention head families in \textit{Mistral-7B}, \textit{Gemma-2-9B}, and \textit{Gemma-2-27B}. While they demonstrate that similar circuits emerge across these models, their study is constrained to a single operator structure (\textit{i.e.}, implication chains) using simplified one-hop problems. Broader propositional rules, including \textit{De Morgan's Laws}, \textit{Distributivity}, and \textit{Commutativity}, as well as applying multiple different rules in a multi-step reasoning process, remain mechanistically unexplored. Moreover, their circuit-discovery paradigm focuses on identifying \textit{which specific components} are necessary for a given model-task pair, rather than characterizing \textit{how} the reasoning process is organized at a higher level of abstraction, namely the generalizable computational strategies that may transfer across model architectures and logical structures.

% In this paper, we address this gap with a more comprehensive analysis. First, 

% we substantially expand the scope of analysis to \textit{11 propositional logic rule categories} across \textit{one-hop and two-hop reasoning} in \textit{Qwen3 (8B, 14B)}, demonstrating that 

% identified mechanisms (Revise)

% generalize not only across rule types but also across reasoning depths. Second, and more fundamentally, we shift analytical focus from model-specific circuit identification to generalizable mechanistic principles. Rather than asking \textit{``which heads are necessary for this specific task,''} we ask \textit{``what computational strategies does the model employ to solve logical problems in general?''} This perspective reveals a \textit{Staged Computational Hierarchy} where early, middle, and late layers assume distinct functional roles, along with three interlocking phenomena: \textit{Information Transmission} (semantic information aggregates at boundary tokens), \textit{Fact Retrospection} (factual information is persistently re-accessed), and \textit{Specialized Attention Heads} (attention heads exhibit stable functional specialization). These phenomena persist across model scales, rule types, and reasoning depths. 

In this paper, we address this gap with a more comprehensive analysis. 
First, we shift analytical focus from model-specific circuit identification to generalizable mechanistic principles. Rather than asking \textit{``which heads are necessary for this specific task,''} we ask \textit{``what computational strategies does the model employ to solve logical problems in general?''} As shown in Figure~\ref{fig:overview}, 
% this perspective reveals a \textit{Staged Computational Hierarchy} where early, middle, and late layers assume distinct functional roles, along with another three interlocking phenomena: \textit{Information Transmission} (semantic information aggregates at boundary tokens), \textit{Fact Retrospection} (factual information is persistently re-accessed), and \textit{Specialized Attention Heads} (attention heads exhibit stable functional specialization). 
this perspective reveals four interlocking mechanisms: (1) Staged Computation, with layer-wise processing phases where early, middle, and late layers assume distinct functional roles; (2) Information Transmission, where semantic information aggregates at boundary tokens; (3) Fact Retrospection, in which factual information is persistently re-accessed; and (4) Specialized Attention Heads, where attention heads exhibit stable functional specialization. 
Second, we 
% substantially 
expand the scope of analysis to 11 propositional logic rule categories across one-hop and two-hop reasoning in \textit{Qwen3} (8B, 14B), demonstrating that these mechanisms persist across model scales, rule types, and reasoning depths. 
Third, we introduce a prompt-order intervention that moves the queried expression before the factual assignments (\texttt{expr\_first}) while preserving the underlying logical content, labels, and corruption semantics. This intervention lets us probe whether staged computation is primarily prompt-oriented or model-oriented. For readability, the main text focuses on \texttt{facts\_first}; the dedicated prompt-order analysis is reported in Appendix~\ref{sec:appendix_prompt_order}. 
Through the above analysis, we provide a more profound, complete and generalizable understanding of the mechanism of propositional logical reasoning. 

%%% PLM: Through the above analysis, we provide a more complete, deeper, and more generalizable picture / understanding of the mechanism of propositional logical reasoning. 


%% PLM: 再加强。我们的区别更显著一点。！分析角度！
% Limitations：1）shallow - deep 2）miss some points, we cover points 3）。。。more logical rules 
% 有什么证据表现出它们分析比较肤浅？有什么分析角度也很重要，但是它们没有做？


\begin{figure*}[t]
    \centering
    % Top Row
    \begin{subfigure}[b]{0.46\textwidth}
        \centering
        \includegraphics[width=\linewidth]{./figures/overview_1_stage.pdf}
        \vspace{-0.5em}
        \caption{\textbf{Staged Computation}}
        \label{fig:overview_pipeline}
    \end{subfigure}
    \hspace{0.5em}
    \begin{subfigure}[b]{0.46\textwidth}
        \centering
        \includegraphics[width=\linewidth]{./figures/overview_2_info.pdf}
        \vspace{-1.8em}
        \caption{\textbf{Information Transmission}}
        \label{fig:info_convergence}
    \end{subfigure}

    \vspace{0.9em}

    % Bottom Row
    \begin{subfigure}[b]{0.46\textwidth}
        \centering
        \includegraphics[width=\linewidth]{./figures/overview_4_facts.pdf}
        \vspace{-0.9em}
        \caption{\textbf{Fact Retrospection}}
        \label{fig:repetitive_looking}
    \end{subfigure}
    \hspace{0.5em}
    \begin{subfigure}[b]{0.46\textwidth}
        \centering
        \includegraphics[width=\linewidth]{./figures/overview_3_heads.pdf}
        \vspace{-0.9em}
        \caption{\textbf{Specialized Attention Heads}}
        \label{fig:attention_heads}
    \end{subfigure}

    \vspace{-0.3em}
    \caption{\textbf{Overview of the propositional logic reasoning mechanisms in \textit{Qwen3}.}
    (a) Staged Computation: MLP patching reveals layer-wise processing phases.
    (b) Information Transmission: Semantic content aggregates at segment-terminal tokens.
    (c) Fact Retrospection: Fact tokens maintain persistent causal influence across depth.
    (d) Specialized Attention Heads: Specialized heads implement the macroscopic mechanisms.}
    \vspace{-0.5em}
    \label{fig:overview}
\end{figure*}

% In this paper, we bridge this gap by conducting a rigorous mechanistic analysis of the \textit{Qwen3} family on a comprehensive suite of boolean logic tasks, constructing a causal graph of the model's inference process based on $11$ standard logical rules.
% %%% PLM: use what methods? 1-2 sentences

% In this paper, we bridge this gap by conducting a rigorous mechanistic analysis of the \textit{Qwen3} family on a comprehensive suite of propositional logic tasks. We employ causal \textit{activation patching} to isolate the necessary model components that mediates correct predictions, coupled with \textit{attention head analysis} to map the functional topology of information flow.
% Our experiments reveal that logical inference in \textit{Qwen3} is not monolithic, but rather manifests as a \textbf{Staged Computational Hierarchy}, where early layers encode facts, middle layers execute logical calculations, and late layers integrate conclusions. 
% Within this macroscopic framework, as shown in Figure~\ref{fig:overview}, three interlocking mechanisms that collectively orchestrate the flow of information:
% (1) \textbf{Information Transmission}: the model gathers distributed context into critical focal points at semantic boundaries;
% (2) \textbf{Specialized Attention Heads}: distinct groups of heads specialize in specific sub-processes, forming a coordinated assembly line akin to collaborative problem-solving, where some heads parse structure, others bind variables, and still others retrieve facts; and
% (3) \textbf{Fact Retrospection}: a self-verification strategy where late-layer heads persistently revisit source premises.
% % (1) \textbf{Information Transmission}, which facilitates efficient transmission by compressing scattered semantic content into discrete ``information hubs'' at boundary tokens; 
% % (2) \textbf{Specialized Attention Heads}, implying a specialized division of labor where distinct heads emerge to execute stage-specific tasks---from ``splitting'' syntax to ``binding'' variables; and 
% % (3) \textbf{Fact Retrospection}, a verification mechanism akin to working memory where late-stage heads persistently re-access original ground-truth tokens.
% % These findings provide structural evidence for the hypothesis that LLMs spontaneously instantiate a progressive computational pipeline that elevates raw token streams into verifiable logical conclusions.
% These findings suggest that LLMs may employ structured, interpretable strategies when processing logical rules, rather than relying solely on opaque pattern matching.
% %%% PLM: 去和前面的别人方法的Limitation或者你提到的Key RQ呼应一下；process/understand logical rule; 不一定要很确定的结论；讲的委婉一点；Positive Evidence

%%% PLM: Cognitive process before computational process



% connection to out work 分布到其他相关工作中；并非堆叠论文介绍。讲清楚和别人工作的区别。将RG展开讲。Review + Gap + Proposal。
\section{Related Work}

\paragraph{Component-Level Mechanistic Interpretability.}
\textit{Mechanistic interpretability} (MI) has identified specific components responsible for individual computational primitives: induction heads for in-context learning \citep{olsson2022incontextlearninginductionheads}, MLP layers that store factual associations \citep{meng2023locatingeditingfactualassociations, geva2023dissectingrecallfactual, ferrando2025informationflowlarge}, and binding ID vectors through which attention heads link entities to their attributes \citep{feng2024languagemodelsrepresent, wang2023labelwordsanchors}. Causal mediation analysis \citep{vig2020causalinterpretingneuralNLP, hase2024doeslocalizationinform, zhang2024towardsbestpractices, heimersheim2024useinterpretactivation} and circuit-based methods \citep{elhage2021mathematicalframeworktransformer, wang2022interpretabilitywildcircuitindirect} extend this by identifying minimal subgraphs causally necessary for specific tasks.
% , while sparse autoencoders \citep{marks2024sparsfeaturecircuits, engels2024decomposingdarkmatter} and function vectors \citep{todd2024functionvectorslarge} further explore distributed representations.
%
% For reasoning specifically, \citet{kim2024reasoningcircuitslanguage} decompose syllogistic inference (\textit{e.g.}, ``All A are B; All B are C; therefore All A are C'') into modular circuits, identifying premise-retrieval and conclusion-assembly components within transformer layers. Their work demonstrates that even deductive reasoning can be localized to identifiable sub-networks. Similarly, \citet{aimpliesb} identify four functional head types (locating, mover, processing, and decision heads) for propositional implication reasoning in \textit{Mistral-7B} and \textit{Gemma-2}.
%
% These studies share a component-level lens: they ask \textit{which} heads or layers implement a given sub-task, providing fine-grained causal evidence for individual mechanisms. However, they leave open the question of how these components coordinate into coherent end-to-end reasoning strategies, particularly across diverse logical rules and multi-hop inference chains.

For reasoning specifically, \citet{kim2024reasoningcircuitslanguage} decompose syllogistic inference into modular circuits with premise-retrieval and conclusion-assembly components, and \citet{aimpliesb} identify four functional head types (locating, mover, processing, and decision heads) for propositional implication in \textit{Mistral-7B} and \textit{Gemma-2}. These studies share a component-level lens, asking \textit{which} heads or layers implement a given sub-task, but leave open how these components coordinate into end-to-end reasoning strategies across diverse logical rules.

% \paragraph{System-Level Computational Strategies.}
% Less explored is how individual components organize into network-wide reasoning strategies. \citet{dutta2024thinkstepbystep} take a step in this direction by tracing how chain-of-thought (CoT) prompting restructures internal information flow, showing that explicit intermediate tokens serve as anchor points that enable staged computation across layers. Their analysis reveals that CoT provides external scaffolding for sequential reasoning steps.
% %
% Our work addresses the complementary question: how do models organize multi-step logical reasoning \textit{without} explicit intermediate tokens? We find that even in the implicit (no-CoT) setting, models exhibit temporal staging across layers, information convergence at boundary tokens, and repeated fact retrieval, which parallel but do not require the external scaffolding that CoT provides. This implicit reasoning perspective also connects to work on entity binding: \citet{feng2024languagemodelsrepresent} show that models route entity identity through binding ID vectors in attention heads, a single-primitive operation. Our Transmission Heads perform an analogous but higher-level function, routing \textit{composed logical content} (not just entity identity) across token positions toward semantic boundary tokens during multi-step reasoning.
% %
% Furthermore, where \citet{meng2023locatingeditingfactualassociations} demonstrate that single-fact retrieval localizes to specific MLP layers that store subject-attribute associations, our Fact Retrospection mechanism reveals that compositional reasoning demands \textit{distributed, recurring} access to premises across network depth. Specifically, models repeatedly re-access premise information at multiple layers rather than retrieving it once. This suggests that the transition from factual recall to logical reasoning involves a qualitative shift in how stored knowledge is accessed.
% %
% Compared with the component-level studies above, our system-level perspective, combined with broader rule coverage (11 categories) and multi-hop scenarios, reveals organizational patterns (\textit{e.g.}, staged computation, boundary-token convergence) that are invisible to analyses focused on individual heads or circuits alone.

\paragraph{System-Level Computational Strategies.}
Less explored is how individual components organize into network-wide reasoning strategies. \citet{dutta2024thinkstepbystep} trace how chain-of-thought (CoT) prompting restructures internal information flow, showing that explicit intermediate tokens serve as anchor points for staged computation. Our work addresses the complementary question: how do models organize multi-step logical reasoning \textit{without} such explicit tokens? We find that even in the implicit (no-CoT) setting, models exhibit temporal staging across layers, information convergence at boundary tokens, and repeated fact retrieval, which parallel but do not require the external scaffolding that CoT provides.
%
Prior component-level studies have shown that models route entity identity through binding ID vectors \citep{feng2024languagemodelsrepresent} and localize single-fact retrieval to specific MLP layers \citep{meng2023locatingeditingfactualassociations}. Our findings extend beyond these single-primitive operations: Transmission Heads route \textit{composed logical content} across token positions, and Fact Retrospection reveals that compositional reasoning demands \textit{distributed, recurring} access to premises across network depth rather than one-time retrieval. These patterns suggest a qualitative shift from factual recall to logical reasoning in how stored knowledge is accessed and coordinated.




% \paragraph{Mechanistic Interpretability of Transformers.}
% This area reverse-engineers transformer computations to understand capabilities like language modeling \citep{olsson2022incontextlearninginductionheads}, factual recall \citep{meng2023locatingeditingfactualassociations, geva2023dissectingrecallfactual, ferrando2025informationflowlarge}, and entity binding \citep{feng2024languagemodelsrepresent, wang2023labelwordsanchors}. 
% The most popular methods used in MI is causal mediation analysis, representive by patching methods
% %%% PLM: transition ??? the most popular methods used in MI is causal mediation analysis, representive by patching methods. 
% % At their core, patching methods are variants of causal mediation analysis 
% \citep{vig2020causalinterpretingneuralNLP, hase2024doeslocalizationinform, zhang2024towardsbestpractices, heimersheim2024useinterpretactivation}. Circuit-based interpretability \citep{elhage2021mathematicalframeworktransformer, wang2022interpretabilitywildcircuitindirect} identifies minimal subgraphs for specific tasks, while more recent work explores distributed representations through sparse autoencoders \citep{marks2024sparsfeaturecircuits, engels2024decomposingdarkmatter} and function vectors \citep{todd2024functionvectorslarge}.
% However, these studies primarily focus on language modeling primitives (\textit{e.g.}, copying, entity tracking) or isolated task components, without examining end-to-end reasoning pathways in multi-step logical tasks. Our work extends mechanistic interpretability to propositional logical reasoning, tracing information flow across residual streams, attention heads, and MLPs for comprehensive logical rules across one-hop and two-hop scenarios, representing a more systematic investigation than prior component-level analyses.

% \paragraph{Coarse-Grained Evaluation of Reasoning in LLMs.}
% Recent work has extensively evaluated LLM reasoning across benchmarks \citep{tafjord2021proofwritergeneratingimplications, hendrycks2021measuringmathematicalproblem, saparov2023languagemodelsgreedy, berglund2024reversalcursellms, dziri2023faithfatelimitstransformers}, including emerging thinking models that extend test-time computation \citep{muennighoff2025s1simpletesttime, ye2025transformerslearnimplicit, wang2025systemthinkingdiagrambased}. These studies primarily benchmark performance on sophisticated tasks, assessing \emph{whether} models can reason. While valuable for understanding capabilities, they do not provide mechanistic insights into \emph{how} models reason internally. Our work complements this by focusing on interpretable, fine-grained analysis of the computational mechanisms underlying logical reasoning.

% \paragraph{Fine-Grained Analysis of How LLMs Reason.}
% Fewer studies provide mechanistic analysis of reasoning. Recent work has examined arithmetic \citep{stolfo2023mechanisticinterpretationarithmeticreasoning, yu2024interpretingarithmeticmechanismcomparative, kantamneni2025languagemodelsuse}, syllogistic reasoning \citep{kim2024reasoningcircuitslanguage}, graph connectivity \citep{saparov2025transformerscanlearn}, and chain-of-thought mechanisms \citep{dutta2024thinkstepbystep}. Some work studies small models trained on symbolic tasks \citep{brinkmann2024mechanisticanalysistransformertrained, wang2024modularadditiontransformer, zhu2024learningcircuittransformers}. These studies primarily investigate other reasoning types (arithmetic, graphs) or use toy models trained specifically for the task. In contrast, we study production-scale pretrained LLMs (Qwen3-8B, 14B) on comprehensive propositional logic rules, providing insights more directly applicable to real-world deployment.

% \paragraph{Multi-hop Reasoning and Propositional Logic.}
% Studies on multi-hop reasoning reveal that transformers can learn to chain facts \citep{ye2025transformerslearnimplicit, wang2024languagemodelscan}, though data requirements grow exponentially with reasoning depth. Implicit reasoning research \citep{wang2024grokkedtransformersimplicit, lin2025implicitreasoningtransformers} investigates how models reason without explicit intermediate steps—\citet{wang2024grokkedtransformersimplicit} show that grokking enables implicit reasoning, while \citet{lin2025implicitreasoningtransformers} reveal reliance on shortcuts.
% For propositional logic specifically, \citet{aimpliesb} provide the only prior mechanistic analysis, identifying four modular sub-circuits for implication reasoning (\textit{e.g.}, ``A or B implies C''). However, their study is limited to a single logical operator structure and simplifies problems to length-2 to ensure tractable analysis. Broader logical rules—De Morgan's Laws, Distributivity, Commutativity—and their underlying mechanisms remain unexplored. We address this gap by investigating a comprehensive suite of propositional logic rules across one-hop and two-hop reasoning in production-scale LLMs, identifying higher-level mechanisms (Information Transmission, Specialized Attention Heads, Fact Retrospection) that generalize across diverse logical structures.
% \subsection{Mechanistic Analysis of Propositional Logic Reasoning}
% % Propositional logic reasoning—requiring Boolean operators and rule chaining—represents fundamental systematic reasoning.

% \paragraph{Mechanistic Analysis and Our Contribution.}
% \citet{aimpliesb} provide the only prior mechanistic analysis, identifying sub-circuits for implication reasoning in Mistral-7B and Gemma-2. However, they examine a \textbf{single operator structure} with \textbf{length-2} simplification. Broader rules—De Morgan's Laws, Distributivity, Commutativity, Associativity, Identity, Domination, and others—remain unexplored. \textbf{We address this gap} by analyzing Qwen3 (8B, 14B) across \textbf{11 propositional logic rule types}, uncovering \textbf{generalizable mechanisms} (Information Transmission, Specialized Attention Heads, Fact Retrospection) that persist across diverse logical structures. Unlike behavioral approaches testing \emph{whether} models reason, and \citet{aimpliesb}'s narrow analysis, we reveal \emph{how} models implement comprehensive propositional reasoning, providing evidence for structured algorithmic computation over pattern matching.


\paragraph{Propositional Logic Reasoning in LLMs}
Benchmarks such as PrOntoQA \citep{saparov2023languagemodelsgreedy}, ProofWriter \citep{tafjord2021proofwritergeneratingimplications}, and LogicBench \citep{parmar2024logicbench} have been developed to evaluate LLM propositional reasoning. Behavioral studies reveal systematic limitations specific to logical inference: accuracy degrades significantly with proof depth \citep{saparov2023languagemodelsgreedy, dziri2023faithfatelimitstransformers}, models struggle with proof planning (selecting correct inference steps when multiple valid options exist \citep{saparov2023languagemodelsgreedy}) and exhibit particular difficulty with negated premises \citep{parmar2024logicbench}. These findings fuel the ongoing debate: do LLMs implement genuine logical algorithms \citep{nanda2023progress}, or merely exploit statistical shortcuts \citep{dziri2023faithfatelimitstransformers}? However, behavioral evaluation alone only answers \textit{whether} models succeed, not \textit{how} models complete the reasoning internally. 
% internal computation 
% the internal computation analysis. 
% To the best of our knowledge, the only mechanistic study on propositional logic is \citet{aimpliesb},  
% Mechanistic studies on propositional logic are scarce. (revise) 
% \citet{aimpliesb} provided an initial circuit-level analysis, identifying retrieval, processing, and decision sub-circuits for implication reasoning in Mistral-7B and Gemma-2. However, their work examines only a \textit{single operator structure} (implication chains) with simplified length-2 problems. Broader propositional rules, including De Morgan's Laws, Distributivity, and Commutativity, remain mechanistically unexplored.
To the best of our knowledge, the only mechanistic study on propositional logic is \citet{aimpliesb}, which decomposes implication reasoning into four functional head types (locating, mover, processing, and decision heads) in \textit{Mistral-7B} 
% TODO: 有点重复
and \textit{Gemma-2} models. Their component-level analysis successfully pinpoints specific heads responsible for rule retrieval, information movement, and answer selection.
Different from this head-level functional decomposition, we characterize network-wide computational patterns: how processing is temporally organized across layers, how information flows across token positions
%toward semantic boundaries 
, and how facts are repeatedly retrieved throughout network depth. This pattern-level perspective, combined with broader rule coverage and multi-hop scenarios, reveals mechanisms that generalize across diverse logical structures, inference chain length and model scales.

% \textit{We address this gap} by analyzing \textit{Qwen3} (8B, 14B) across \textit{11 propositional logic rule categories} with one-hop and two-hop reasoning. We uncover \textbf{three generalizable mechanisms}, namely Information Transmission, Specialized Attention Heads, and Fact Retrospection, that persist across diverse logical structures and model scales, providing mechanistic evidence for how LLMs implement propositional reasoning.

\section{Methodology}
% In this section, we formalize the propositional logic reasoning task, describe the construction of our dataset, and introduce the mechanistic interpretability techniques used to dissect the model's internal computations.

Our approach departs from the circuit-discovery paradigm \citep{elhage2021mathematicalframeworktransformer, wang2022interpretabilitywildcircuitindirect} 
% that asks 
which primarily asks 
% TODO: 侧重circuit
``which heads are necessary?'' 
We instead focus on ``what computational strategies does the model employ?'' This shift is motivated by two considerations. First, circuit faithfulness is difficult to establish, as multiple circuits may implement the same function and ablating one may simply shift computation to another. Second, circuits often fail to generalize across different models, limiting their scientific value for understanding \emph{how transformers reason} rather than \emph{how a particular model is wired}. Accordingly, our analysis characterizes network-wide computational patterns rather than minimal circuits.


\subsection{Task Formulation}
% We formulate the propositional logical reasoning task as a next-token prediction problem. Let $\mathcal{M}$ denote a causal language model. Given a prompt sequence $x = [F_1, F_2, ...R_1,...F_k,...R_m, Q]$, where $\{F_i\}$ represents a set of premises (facts), $\{R_i\}$ are logic expression, and $Q$ is the query, the model acts as a function that maps the sequence to a probability distribution over the vocabulary: $P(y|x) = \mathcal{M}(x)$. For a Boolean logic task, the target output $y$ is constrained to binary truth values: $y \in \{\text{True}, \text{False}\}$. To strictly evaluate implicit reasoning, we constrain the model to output the answer immediately following the query, prohibiting explicit chain-of-thought generation. This forces the model to perform the reasoning steps internally within its hidden states.

We formulate the propositional logical reasoning task as a next-token prediction problem. Given a prompt $x = [F_1, \ldots, F_k, R_1, \ldots, R_m, Q]$, where $\{F_i\}$ are premises (\textit{i.e.}, facts), $\{R_j\}$ are logical expressions, and $Q$ is the query, the model $\mathcal{M}$ predicts the answer $y \in \{\texttt{True}, \texttt{False}\}$. To strictly evaluate implicit reasoning, we require the model to output the answer immediately following the query without \textit{Chain-of-Thought} (CoT) generation, 
% forcing all reasoning to occur internally within hidden states. 
confining all reasoning to internal hidden states.


\subsection{Dataset Construction}
\label{sec:dataset}
\paragraph{PropLogic-MI}
While benchmarks like \textit{PrOntoQA} and \textit{ProofWriter} have established a foundation for evaluating reasoning accuracy in LLMs, they remain insufficient for fine-grained MI. These datasets often introduce uncontrolled linguistic variability and entangled reasoning steps, which act as confounding factors for causal analysis methods like activation patching. 
To address this limitation, we introduce \textit{PropLogic-MI}, a highly controlled dataset designed specifically for circuit discovery. \textit{PropLogic-MI} focuses on atomic logical operations across 11 propositional logic rule categories (Table~\ref{tab:rule_coverage}), providing strictly aligned clean/corrupted pairs for one-hop and two-hop reasoning. This structural precision enables researchers to isolate the exact neural components responsible for logical transitions, shifting the focus from measuring performance to dissecting the underlying mechanism.

% \paragraph{Rule Coverage.}
% Our dataset encompasses 11 fundamental propositional logic rule categories, spanning core operations of Boolean algebra.
% Table~\ref{tab:rule_coverage} summarizes each rule type with its formal definition and illustrative prompt templates.

% \begin{table*}[t]
% \centering
% \small
% % 适当减小一点列间距，让内容更紧凑
% \setlength{\tabcolsep}{6pt}
% % 布局调整：
% % l: Category 左对齐
% % c: Formula 居中对齐 (数学公式居中更美观)
% % X: Prompt 自动填充，且用左对齐
% \begin{tabularx}{\textwidth}{@{}lcX@{}}
% \toprule
% \textbf{Rule Category} & \textbf{Formal Definition} & \textbf{Example Prompt Template} \\
% \midrule
% Identity 
%     & $P \land \top \equiv P$ 
%     & \texttt{"A is True, A and True is"} \\
%     & $P \lor \bot \equiv P$ 
%     & \texttt{"A is False, A or False is"} \\
% \addlinespace[0.4em] % 稍微增加组间距，阅读更轻松
% Domination 
%     & $P \land \bot \equiv \bot$ 
%     & \texttt{"A is True, A and False is"} \\
%     & $P \lor \top \equiv \top$ 
%     & \texttt{"A is False, A or True is"} \\
% \addlinespace[0.4em]
% Idempotent 
%     & $P \land P \equiv P$ 
%     & \texttt{"A is True, A and A is"} \\
%     & $P \lor P \equiv P$ 
%     & \texttt{"A is False, A or A is"} \\
% \addlinespace[0.4em]
% Double Negation 
%     & $\lnot(\lnot P) \equiv P$ 
%     & \texttt{"A is True, ($\lnot$($\lnot$A)) is"} \\
% \addlinespace[0.4em]
% Excluded Middle 
%     & $P \lor \lnot P \equiv \top$ 
%     & \texttt{"A is True, A or $\lnot$A is"} \\
% \addlinespace[0.4em]
% Contradiction 
%     & $P \land \lnot P \equiv \bot$ 
%     & \texttt{"A is False, A and $\lnot$A is"} \\
% \addlinespace[0.4em]
% Commutative 
%     & $P \land Q \equiv Q \land P$ 
%     & \texttt{"A is T, B is F, A and B is"} \\ % 稍微缩写 T/F 如果行太挤
% \addlinespace[0.4em]
% Associative 
%     & $(P \land Q) \land R \equiv P \land (Q \land R)$ 
%     & \texttt{"A is T, B is T, C is F, (A and B) and C is"} \\
% \addlinespace[0.4em]
% Distributive 
%     & $P \land (Q \lor R) \equiv (P \land Q) \lor (P \land R)$ 
%     & \texttt{"A is T, B is F, C is T, A and (B or C) is"} \\
% \addlinespace[0.4em]
% De Morgan 
%     & $\lnot(P \land Q) \equiv \lnot P \lor \lnot Q$ 
%     & \texttt{"A is T, B is F, ($\lnot$A or $\lnot$B) is"} \\
% \addlinespace[0.4em]
% Absorption 
%     & $P \land (P \lor Q) \equiv P$ 
%     & \texttt{"A is T, B is F, A and (A or B) is"} \\
% \bottomrule
% \end{tabularx}
% \caption{\textbf{Propositional logic rule coverage.} Our dataset systematically covers 11 fundamental rule categories. For each rule covers its formal Boolean algebra definition and a corresponding template. Variables \{A, B, C\} are instantiated with truth values (True/False). Note that for longer prompts, we abbreviate True/False as T/F in this table.}
% \label{tab:rule_coverage}
% \end{table*}





\paragraph{Prompt Template}
% For each rule, we instantiate prompts by randomly assigning truth values to atomic propositions in all possibility. For example, for \textit{De Morgan's Law}:
For each rule, we instantiate prompts by enumerating all possible truth value assignments to atomic propositions. For example, for \textit{De Morgan's Law}:
\begin{center}
\texttt{"A is True, B is False, (¬A or ¬B) is"}
\end{center}
where the model is expected to predict the correct truth value (\texttt{True} in this case). The ground-truth answer is computed via symbolic evaluation.

\paragraph{Reasoning Depth}
We construct prompts under two complexity settings: (1) One-Hop Reasoning, where the query is directly derived from facts via a single rule, and (2) Two-Hop Reasoning, where an intermediate conclusion must first be derived before predicting the final answer. Representative examples are provided in Table~\ref{tab:reasoning_depth}. 
% (Appendix~\ref{sec:appendix_dataset}).


% \paragraph{Evaluation Protocol.}
% We analyze pre-trained models in a zero-shot setting without fine-tuning. For each prompt, we measure the model's logit difference (Eq.~\ref{eq:logit_diff}) to assess reasoning performance and conduct activation patching experiments to identify causally important components.



\subsection{Analytical Methods}
\label{sec:analytical_methods}

\paragraph{Causal Framework}
% To rigorously trace the information flow underlying logical reasoning, we move beyond correlational analysis (\textit{e.g.}, attention maps) to establish causal sufficiency. To this end, we employ \textit{Activation Patching}~\citep{vig2020causalinterpretingneuralNLP, meng2023locatingeditingfactualassociations}, a method grounded in counterfactual interventions. The core idea follows a specific protocol: construct a \emph{clean} input that yields the correct answer and a \emph{corrupted} input that yields a different (counterfactual) prediction. We then selectively replace (\textit{i.e.} patch) activations in the corrupted run with their clean counterparts. If patching a specific component causes the model's output to recover the correct answer, we identify this component as critical in the logical circuit.
To rigorously trace information flow underlying logical reasoning, we move beyond correlational analysis (\textit{e.g.}, attention maps) to establish causal sufficiency. We employ \textit{Activation Patching}, a method grounded in counterfactual interventions, and the procedure is as follows: construct a \emph{clean} input that yields the correct answer and a \emph{corrupted} input that yields a different prediction; then selectively patch activations in the corrupted run with their clean counterparts. If patching a specific component recovers the correct answer, we identify it as causally critical.

% \paragraph{Causal Framework.}
To rigorously trace the information flow underlying logical reasoning, we move beyond correlational analysis (\textit{e.g.}, attention maps) to establish \textit{causal necessity}. 
% Mere high activation or attention weights do not guarantee a component's contribution to the final output. 
% Therefore, 
To this end, we employ \textbf{Activation Patching}~\citep{vig2020causalinterpretingneuralNLP, meng2023locatingeditingfactualassociations}, 
%%% PLM: AP 怎么做的 A -> Patch -> B...
%%% PLM: judge the causal relation between a cause A and an effect B; 怎么实现？
a method grounded in counterfactual interventions. The core premise is to isolate specific model components and ask: \emph{Does restoring the internal state of this component from a correct execution implies the recovery of the correct output in a corrupted execution?} By systematically intervening on activations, we can distinguish between components that merely process features and those that act as critical circuit nodes for deductive inference.


\paragraph{Implementation and Metrics}
We construct paired inputs for this counterfactual setup: a \emph{clean} input $x_{\text{clean}}$ yielding target answer $Y_{\text{clean}}$, and a \emph{corrupted} input $X_{\text{corrupt}}$ minimally edited (\textit{e.g.}, flipping a truth value) to induce a contrasting prediction $Y_{\text{corrupt}}$.
% We operate paired inputs to construct this counterfactual setup. We define a clean input $x_{\text{clean}}$ designed to predict a specific target answer $y_{\text{clean}}$, and a corrupted input $x_{\text{corrupt}}$, which is minimally edited (\textit{e.g.}, naturally flipping a truth value) to induce a contrasting prediction $y_{\text{corrupt}}$.
We first perform a forward pass on $X_{\text{clean}}$ to cache activations $h_{l,t}$ (or $h_{l,h,t}$ for attention head $h$) at layer $l$ and token position $t$. And then we perform a forward pass on $X_{corrupt}$ and get $Y_{corrupt}$. We again run a forward pass on $X_{\text{corrupt}}$, at this time we intervene by replacing the target activation with its cached clean counterpart and get $Y_{patched}$.
% We first perform a forward pass on $x_{\text{clean}}$ to cache the activations $h_{l,i}$ at layer $l$ and position $i$. During a forward pass on $x_{\text{corrupt}}$, we intervene by replacing the activation at $(l, i)$ with the cached $h_{l,i}$.
In some experiments, we also apply mean ablation (i.e., replacing the output with its mean activation) to selected model components to evaluate their functional roles, such as MLP outputs and attention head outputs.

\noindent We quantify the causal effect of this intervention using the \textit{Probability Difference} (PD) metric. Let $w_{\texttt{True}}$ and $w_{\texttt{False}}$ denote the tokens for the two truth values(assume that \texttt{True} is the correct answer on $X_{clean}$, while \texttt{False} is the correct answer on $X_{corrupt}$, normally the two are mutually exclusive). The PD on the patched output is:
% $LogitDiff = Logits(w_{true})-Logits(w_{false}$
\begin{equation}
    \text{PD} = \text{P}(w_{\texttt{True}}) - \text{P}(w_{\texttt{False}})
    \label{eq:logit_diff}
\end{equation}
% A positive $\Delta \text{Logits}$ indicates that the component $(l,i)$ retains critical information required to steer the prediction back to $y_{\text{clean}}$.
Here $P(w)$ denotes the probability assigned by the model to token $w$ in the output vocabulary. A positive \text{PD} indicates that the component $(l,t)$ (or $(l,h,t)$) retains critical information required to steer the prediction back to $y_{\text{clean}}$.
To quantify the causal effect of patching, we define the \textit{Probability Difference Shift} ($\text{dPD}$):
\begin{equation}
    % \text{dLD}_{l,t} = \text{LD}_{\text{patched}}^{(l,t)} - \text{LD}_{\text{corrupt}}
    \text{dPD} = \text{PD}_{\text{patched}} - \text{PD}_{\text{baseline}}
    \label{eq:dld}
\end{equation}
where $\text{PD}_{\text{baseline}}$ is the logit difference on the unpatched corrupted run. A positive $\text{dPD}$ indicates successful restoration toward the correct answer.


% We quantify the causal effect of this intervention using the \textit{Logit Difference} (LD) metric on the patched output:
% % \begin{equation}
% %     \label{eq:logit_diff}
% %     % \Delta \text{Logits} = \text{logit}_{\text{patched}}(y_{\text{clean}}) - \text{logit}_{\text{patched}}(y_{\text{flipped}})
% %     \Delta \text{Logits} = \text{logit}(y_{\text{clean}}) - \text{logit}(y_{\text{flipped}})
% % \end{equation}
% \begin{equation}
%     \text{LogitDiff} = \text{Logits}(w_{\texttt{True}}) - \text{Logits}(w_{\texttt{False}})
%     \label{eq:logit_diff}
% \end{equation}
% A positive $\Delta \text{Logits}$ indicates that the component $(l,i)$ retains critical information required to steer the prediction back to $y_{\text{clean}}$.

% To trace the information flow, we primarily employ \textit{Activation Patching} (also known as \textit{Causal Mediation Analysis}).
% We define a clean input $x_{clean}$ (yielding the correct answer) and a corrupted input $x_{corrupt}$ (yielding an incorrect answer or a flipped truth value). We perform a forward pass on $x_{clean}$ and cache the activations $h_{l,i}$ at layer $l$ and position $i$. During a forward pass on $x_{corrupt}$, we intervene by replacing the activation at a specific $(l, i)$ with the cached $h_{l,i}$ from the clean run.
% We verify the causal effect of this intervention using the \textbf{Logit Difference} metric:
% \begin{equation}
%     \label{eq:logit_diff}
%     \Delta \text{Logits} = \text{logit}(y_{correct}) - \text{logit}(y_{wrong})
% \end{equation}
% A high positive shift in $\Delta \text{Logits}$ after patching indicates that the patched component stores critical information for the correct deduction.

\subsection{Experimental Setup}
\label{sec:experiments-setup}
% We analyze pretrained \textit{Qwen3} models on \textit{PropLogic-MI} using causal intervention analysis.

\paragraph{Models}
Our primary analysis centers on \textit{Qwen3-8B} (36 layers, 32 attention heads per layer), with same-family replication on \textit{Qwen3-14B}. The current codebase and appendix also include cross-model extensions to \textit{Llama-3.1-8B-Instruct} and \textit{Mistral-7B-Instruct-v0.1}. All analyses use pretrained checkpoints in zero-shot settings without fine-tuning.

% \paragraph{Dataset.}
% From the 370 examples generated via template instantiation (74 one-hop, 296 two-hop), we retain 224 examples (60.5\%) where Qwen3-8B correctly predicts both clean and corrupted prompts—42 one-hop (56.8\%) and 182 two-hop (61.5\%). This filtering ensures activation patching measures genuine mechanism disruption rather than pre-existing failures, while maintaining balanced coverage across reasoning depths and rule categories.


% \paragraph{Dataset}
\paragraph{PropLogic-MI}
We constructed a total of $20000$ samples ($10000$ one-hop, $10000$ two-hop) by populating the templates in Table~\ref{tab:rule_coverage} with randomized truth values, thereby generating strictly aligned clean and corrupted pairs. Under the default \texttt{facts\_first} prompt order used in the main text, \textit{Qwen3-8B} achieves $86.98\%$ clean accuracy with a $75.44\%$ dual-correct rate ($15087/20000$), while \textit{Qwen3-14B} reaches $87.42\%$ clean accuracy with a $76.55\%$ dual-correct rate ($15310/20000$). The alternative \texttt{expr\_first} prompt order is substantially harder: clean accuracy drops to $74.22\%$ / $74.39\%$ and the dual-correct rate drops to $53.89\%$ / $52.80\%$ for \textit{Qwen3-8B} / \textit{Qwen3-14B}. We use this harder setting as a prompt-order intervention: because it preserves the logical task while reordering the surface presentation, it lets us test whether staged computation follows the prompt layout or reflects a more model-internal computational preference. We therefore keep \texttt{facts\_first} in the main figures for readability, and report the prompt-order comparisons in Appendix~\ref{sec:appendix_performance} and Appendix~\ref{sec:appendix_prompt_order}. Unless otherwise specified, our mechanistic analysis is conducted on samples that the model answers correctly on both the clean and corrupted prompts.
% To ensure the validity of our mechanistic interventions, we constraint that the model must accurately predict the clean and corrupt pair to differentiate genuine circuit disruption from pre-existing performance failures. 
% Following evaluation with \textit{Qwen3-8B}, we retain $224$ validated examples ($60.5\%$), comprising $42$ one-hop ($56.8\%$) and $182$ two-hop ($61.5\%$) cases. 
% This filtered subset preserves sufficient statistical power while maintaining balanced coverage across reasoning depths and rule categories.


\paragraph{Evaluation}
We apply activation patching and mean-ablation analyses across all 11 rule categories (§\ref{sec:dataset}) using $\text{PD}$ (Eq.~\ref{eq:logit_diff}) and $\text{dPD}$ (Eq.~\ref{eq:dld}) to quantify component importance. The full experimental inventory now includes: (1) residual-stream patching for token-level information flow; (2) region-level mean ablation for MLP outputs; (3) matched region-level mean ablation for attention outputs; (4) a joint attention+MLP region ablation; and (5) a three-step attention-head discovery / validation pipeline. The main text focuses on the Qwen3 MLP, token, and head results, while the additional region-level variants and cross-model extensions are summarized in the appendix. All analyses run on NVIDIA H100 GPUs using TransformerLens~\citep{nanda2022transformerlens}. Detailed implementation is provided in Appendix~\ref{sec:appendix_setup}.




% We conduct activation patching experiments on Qwen3-8B (36 layers, 32 heads per layer) across all 11 logical rule categories described in §\ref{sec:dataset}. To verify cross-model consistency, key experiments are replicated on Qwen3-14B. All analyses use pre-trained models without fine-tuning, executed on NVIDIA H100 GPUs using TransformerLens. Detailed experimental setup is provided in Appendix~\ref{sec:appendix_setup}.




% \subsection{Models and Evaluation}
% \label{sec:models}

% We primarily use Qwen3-8B (36 layers, 32 attention heads per layer) from Hugging Face. To verify cross-model consistency, key experiments are replicated on Qwen3-14B. All experiments use pre-trained checkpoints without fine-tuning, analyzing zero-shot logical reasoning capabilities. Evaluation uses the logit difference metric (Eq.~\ref{eq:logit_diff}) to quantify the causal importance of model components. Detailed experimental setup, including hardware specifications and implementation details, is provided in Appendix~\ref{sec:appendix_setup}.

%\begin{figure*}[t]
%  \includegraphics[width=\textwidth]{./figures/1.pdf}
%  \caption{A figure describing the \textbf{Information Transmission} mechanism. At first few layers, the model splits the logical reasoning problem into different semantic regions. In this example, they are clause 1, clause 2 and logical expression. And then, in each region, the information converges into the last token position, when we call it \textbf{Convergence Point}. Specifically, the fact regions(clause 1 and clause 2) converge earlier than the logical expression region. Notably, in some other examples, the fact regions may converge at a very late moment. And at a specific layer, the all information converges into the last token position to form the final answer.}
%  \label{fig:information_convergence}
%\end{figure*}
%



% \begin{figure*}[t]
%   \includegraphics[width=\textwidth]{./figures/fig1.pdf}
% 	\caption{\textbf{Overview of the mechanistic landscape in Qwen3 during logical reasoning.} The schematic illustrates three governing behaviors identified in our analysis: (1) \textbf{Information Transmission}, where contextual information within semantic intervals progressively aggregates at the segment's terminal token; (2) \textbf{Specialized Attention Heads}, which exhibit stable, specialized functional roles (\textit{e.g.}, splitting, binding, and retrieving) across layers; and (3) \textbf{Fact Retrospection}, characterized by the sustained re-accessing of factual token representations throughout the network's depth. These mechanisms collectively orchestrate the propagation and maintenance of logical information.}
% 	\label{fig:overview}
% \end{figure*}

\begin{figure*}[t]
    \centering
    \begin{subfigure}[b]{0.31\textwidth}
      \centering
      \includegraphics[width=\linewidth]{./figures_experiments/reports/mlp_analysis/mlp_8b_one_hop/facts_region/mlp_facts_region_scores.png}
      \caption{Facts (One-hop)}
      \label{fig:stage_facts}
    \end{subfigure}
    \hfill
    \begin{subfigure}[b]{0.31\textwidth}
      \centering
      \includegraphics[width=\linewidth]{./figures_experiments/reports/mlp_analysis/mlp_8b_one_hop/expression_region/mlp_expression_region_scores.png}
      \caption{Expression (One-hop)}
      \label{fig:stage_expression}
    \end{subfigure}
    \hfill
    \begin{subfigure}[b]{0.31\textwidth}
      \centering
      \includegraphics[width=\linewidth]{./figures_experiments/reports/mlp_analysis/mlp_8b_one_hop/query_region/mlp_query_region_scores.png}
      \caption{Query (One-hop)}
      \label{fig:stage_query}
    \end{subfigure}

    \vspace{0.8em}

    \begin{subfigure}[b]{0.31\textwidth}
      \centering
      \includegraphics[width=\linewidth]{./figures_experiments/reports/mlp_analysis/mlp_8b_two_hop/facts_region/mlp_facts_region_scores.png}
      \caption{Facts (Two-hop)}
      \label{fig:stage_facts_twohop}
    \end{subfigure}
    \hfill
    \begin{subfigure}[b]{0.31\textwidth}
      \centering
      \includegraphics[width=\linewidth]{./figures_experiments/reports/mlp_analysis/mlp_8b_two_hop/expression_region/mlp_expression_region_scores.png}
      \caption{Expression (Two-hop)}
      \label{fig:stage_expression_twohop}
    \end{subfigure}
    \hfill
    \begin{subfigure}[b]{0.31\textwidth}
      \centering
      \includegraphics[width=\linewidth]{./figures_experiments/reports/mlp_analysis/mlp_8b_two_hop/query_region/mlp_query_region_scores.png}
      \caption{Query (Two-hop)}
      \label{fig:stage_query_twohop}
    \end{subfigure}

    \caption{Region-wise MLP mean-ablation scores in \textit{Qwen3-8B} under \texttt{facts\_first}, rendered directly from the saved analysis outputs in
  \texttt{figures\_experiments/reports/mlp\_analysis}. Top row: one-hop. Bottom row: two-hop. Across both depths, the Facts Region remains strongly early-dominant, while the
  Expression Region occupies the early-to-middle part of the network. Two-hop reasoning preserves this coarse staged template but broadens the later computation: the Facts Region
  early-band BCR decreases from $0.856$ to $0.791$, the Expression Region middle-band BCR rises from $0.304$ to $0.380$, and the Query Region shifts from a broader early-middle
  profile in one-hop to a more clearly middle-to-late bridge in two-hop (middle/late BCR $0.329/0.093 \rightarrow 0.487/0.140$). Figure~\ref{fig:staged_computation_14B} extends the
  same MLP analysis to \textit{Qwen3-14B}, and Figure~\ref{fig:attn_region_triples} shows the matched attention-region results.}
    \label{fig:stage}
    \vspace{-0.5em}
  \end{figure*}

% 说明文，总分的方式（较为清晰）。
% 先总述一下，再分点撰写。
% 一张总图将三个机制都概括进去。
% Method： 先讲目的，再讲方法，不要直接平铺直叙讲方法。 为了....， 我们....。（Motivation + method)
% Experiments （找一个模板，A implies B），3~4篇相关论文参考，查看实验部分的组织。层次递进的效果。先讲重要的发现。
% 展现实验结果时：精确（表格，数字）和生动（搭配具体例子，搭配小图）。
\section{Experiments}
\label{sec:results}

% We adopt a top-down analytical strategy, progressing from macroscopic computational patterns to microscopic attention mechanisms. Our investigation begins with \textit{Staged Computation}, which reveals that \textit{Qwen3} processes logical problems in distinct layer-wise phases. This raises a natural question: \textit{how is information organized and transmitted across these stages?} We address this through \textit{Information Transmission} analysis, discovering that semantic content aggregates at segment-terminal tokens. Notably, this analysis reveals an unexpected phenomenon: factual information remains causally influential far deeper into the network than a simple encode-then-forget model would predict. This observation motivates our \textit{Fact Retrospection} hypothesis, suggesting that the model actively revisits source facts throughout processing. Finally, we validate these macroscopic findings at the attention head level, identifying \textit{Specialized Attention Heads} with specialized, stable functional roles (Splitting, Transmission, and Fact-Retrieval) that mechanistically implement the observed information flow patterns. Figure~\ref{fig:overview} provides a schematic overview.


% We present findings in four subsections. We begin with Staged Computation (§\ref{sec:staged}), establishing that layer-wise processing phases exist. This motivates examining Information Transmission (§\ref{sec:info})—how information flows between stages. The analysis reveals an unexpected finding: fact tokens maintain causal influence deep into the network, prompting our Fact Retrospection investigation (§\ref{sec:retro}). Finally, Specialized Attention Heads (§\ref{sec:heads}) identifies specific heads implementing these mechanisms. Figure~\ref{fig:overview} provides a schematic overview.

We present our findings in four subsections, progressing from macroscopic to microscopic analysis. Section~\ref{sec:staged} establishes layer-wise processing phases via MLP patching. Sections~\ref{sec:info} and~\ref{sec:retro} trace information flow through residual stream analysis, revealing both transmission patterns and persistent premises (\textit{i.e.}, truth values) access. Section~\ref{sec:heads} validates these patterns at the attention head level. Figure~\ref{fig:overview} provides a schematic overview.

\subsection{Staged Computation}
\label{sec:staged}

We begin with a fundamental question: \textit{when does computation occur
during logical reasoning, and what is processed at each stage?} 
% TODO: 为什么MLP？动机不够
Following
prior interpretability work that identifies MLP layers as primary sites for knowledge processing \citep{geva-etal-2021-transformer}, 
% TODO: 或许可以再加几篇工作证明出发点
% we employ MLP patching to reveal when different computational steps occur. 
% we employ MLP patching to reveal when different parts of the input are processed.
% Specifically, we hypothesize that MLP layers handle binding entities to truth values,
% retrieving logical rules, and applying them to derive answers.
% We begin with a fundamental question: \textit{when does computation occur during logical reasoning, and what is processed at each stage?} In transformer architectures, attention mechanisms primarily perform linear operations (aggregating and routing information across tokens), while MLPs execute the nonlinear transformations that constitute actual computation. We hypothesize that for logical reasoning, the core computational work should be concentrated in MLP layers, including binding entities to their truth values, retrieving relevant logical rules, and applying them to derive the answer. Based on this, we employ MLP patching to reveal \textit{when} different parts of the input are processed.
we employ MLP patching to reveal when different parts of the input are processed via partitioning each prompt
into three semantically distinct regions: the \textit{Facts Region} (\textit{e.g.}, \texttt{A is True, B is False}), the \textit{Expression Region} (\textit{e.g.}, \texttt{($\neg$A or $\neg$B)}), and the \textit{Query Region}, which contains the query anchor together with the answer-format tail. For each region, we perform mean ablation on the MLP layer outputs across all layers in clean-prompt runs, quantifying the causal impact via 
$\text{dPD}$.  
The results are illustrated in Figure~\ref{fig:stage}.

The results reveal a clear staged pattern. In early layers (L0-8), the Facts Region exhibits the highest patching effect, indicating that factual information (entity-value bindings) is primarily processed at this stage. In middle layers (L10-16), the Expression Region becomes dominant, reflecting the resolution of logical operators. The Query Region is primarily important in late layers where it aggregates preceding computations into the final prediction, though it also shows notable influence in the middle of the network. Note that precise layer boundaries vary slightly across different analytical methods (MLP patching vs. residual stream patching), but the qualitative three-stage pattern remains consistent.

This staged organization (\textit{i.e.}, facts first, expression second, integration last) mirrors the compositional structure of the input itself. The model respects the semantic hierarchy inherent in logical statements, handling premises before conclusions. This suggests that \textit{Qwen3} implements a systematic, compositional approach to logical reasoning, where the temporal sequence of computation aligns with the logical dependencies in the task. Appendix~\ref{sec:appendix_prompt_order} shows that this coarse early/middle/late decomposition survives the \texttt{expr\_first} intervention, although the relative weight of each region shifts with prompt order.

The current codebase extends this region-level analysis beyond MLPs. Figure~\ref{fig:qwen3_region_band_overview} summarizes the three completed Qwen3-8B \texttt{facts\_first} band-level overviews under a matched layout: MLP region ablation, attention-region ablation, and joint attention+MLP ablation. The overall stage ordering is preserved across all three branches, but the attention branch assigns relatively more mass to middle-layer Expression processing and late Query integration, whereas the joint branch remains closest to the original MLP decomposition while sharpening early Facts dominance. We therefore use the MLP branch as the main staged-computation probe, and treat the other two as supporting region-level interventions on the same computation.

\begin{figure}[t]
    \centering

    \begin{subfigure}[b]{0.95\linewidth}
        \centering
        \includegraphics[width=\linewidth]{./figures_experiments/reports/mlp_analysis/comparison/band_metrics_panel.png}
        \caption{MLP}
        \label{fig:qwen3_band_mlp}
    \end{subfigure}

    \vspace{0.5em}

    \begin{subfigure}[b]{0.95\linewidth}
        \centering
        \includegraphics[width=\linewidth]{./figures_experiments/reports/attn_analysis/comparison/band_metrics_panel.png}
        \caption{Attention}
        \label{fig:qwen3_band_attn}
    \end{subfigure}

    \vspace{0.5em}

    \begin{subfigure}[b]{0.95\linewidth}
        \centering
        \includegraphics[width=\linewidth]{./figures_experiments/reports/attn_mlp_analysis/comparison/band_metrics_panel.png}
        \caption{Attention+MLP}
        \label{fig:qwen3_band_attn_mlp}
    \end{subfigure}

    \caption{\textbf{Qwen3-8B region-level band summaries across the three patching analyses.} All three panels correspond to the same \texttt{facts\_first} setting and are rendered from the saved \texttt{band\_metrics\_panel.png} outputs. Each panel compares one-hop and two-hop settings across the Facts, Expression, and Query regions using the shared BMI/BCR/SBI summary. The MLP and joint attention+MLP branches show the clearest early-facts / middle-expression / late-query decomposition, while the attention branch shifts comparatively more mass toward middle-layer expression routing and late query integration.}
    \label{fig:qwen3_region_band_overview}
\end{figure}


% In this section, we empirically investigate the internal mechanisms of the Qwen3 model during logical reasoning tasks. To ensure a tractable analysis of valid reasoning pathways, we focus our experiments primarily on one-hop and two-hop reasoning scenarios, as the model demonstrates consistent performance in these regimes compared to more complex logical problems. Our analysis unveils three primary structural behaviors governing the model's inference process: (1) \textbf{Information Transmission}, where semantic information within a specific text segment progressively aggregates at the final token of that interval; (2) \textbf{Specialized Attention Heads}, a phenomenon where attention heads assume specialized, persistent functional roles—such as semantic splitting, entity binding, or fact retrieval—across different layers and tasks; and (3) \textbf{Fact Retrospection}, a mechanism wherein critical factual information is actively maintained and revisited throughout the network's depth rather than being processed solely in early layers. Collectively, these findings provide a granular view of how Qwen3 organizes, propagates, and sustains logical information, which is schematically summarized in Figure~\ref{fig:overview}.

%During our experiments, we find that the Qwen3 model isn't capable to solve some very complex logical reasoning problems. For mechanistic interpretibility, we should carefully choose some good questions, which can't be to complex or to easy. So, we mainly focus on one-hop and two-hop logical reasoning followed by section 3.2, which the model can handle most of the time.
%
%In this section, we will delve into the logical reasoning process of the model with thoroughness and present the mechanisms we have uncovered. And we will begin by examining a simple case of single-step reasoning and then progress to two-step reasoning.

% TODO: MLP patching score figure

% \begin{figure}[t]
%   \includegraphics[width=\columnwidth]{./figures/converge_1.png}
%   \includegraphics[width=\columnwidth]{./figures/converge_2.png}
%   \caption{Importance ratio at each}
%   \label{fig:logit difference}
% \end{figure}

\subsection{Information Transmission}
\label{sec:info}

\paragraph{Tracing Information Flow.} 
% Having established that computation occurs in distinct stages, a natural question arises: \textit{How does information move from one stage to the next?} 
% If facts are encoded in early layers and predictions emerge at the query token in late layers, there must be a transmission mechanism.
% We investigate this 
We next investigate \textit{how information flows between computational stages}
through residual stream patching, 
% given that the residual stream 
as it serves as the primary information pathway and accumulates representations across layers. By patching residual stream activations at specific token positions, we can trace \textit{where} stage-relevant information resides as it propagates through the network.

We conduct a controlled activation patching study: given a clean prompt $x_{\text{clean}}$ (\textit{e.g.}, \texttt{A is True, B is False, ($\neg$A or $\neg$B) is}) and a corrupt prompt $x_{\text{corrupt}}$ (\textit{e.g.}, A is True, B is True...), we patch token-wise residual stream activations from the clean run into the corrupt run and measure the resulting shift in logit difference. This reveals which token positions carry the causal information that determines the correct answer.
The resulting causal traces (Figure~\ref{fig:logit_difference}) reveal a clear pattern of information transmission. In Figure~\ref{fig:logit_difference}a, the patching effect is initially concentrated at token position 6 (\textit{i.e.}, the truth value \texttt{True}) in early layers, then shifts to token position 14 (\textit{i.e.}, the closing parenthesis \texttt{)}), and finally converges at the query token in late layers. Intra-segment tokens show negligible effects throughout. This trajectory indicates that information progressively aggregates at segment-terminal tokens before flowing to the final prediction position.

\paragraph{Pattern Quantification.} We corroborate this finding with an alternative patching configuration (Figure~\ref{fig:logit_difference}b), which identifies positions 2, 6, 14, and 15 as dominant convergence points. To quantify this pattern, we analyze the mean \text{|dPD|} (\textit{i.e.}, absolute probability-difference shift) across token categories and layer groups (Figure~\ref{fig:info_converge_token}). The results confirm that \texttt{facts\_value} tokens dominate in early layers (L0--13), while \texttt{query\_token} shows a substantial surge in late layers (L24--35), with \text{Mean |dPD|} values increasing over 2.5$\times$. Extended experiments for \textit{Qwen3-14B}, two-hop reasoning, and the \texttt{expr\_first} prompt-order intervention are provided in Appendix~\ref{sec:appendix_info_converge}, Appendix~\ref{sec:appendix_prompt_order}, and Figure~\ref{fig:info_converge_appendix}.
These findings establish that information transmission in \textit{Qwen3} follows a structured pattern: semantic content aggregates at the final token of each semantic segment (\textit{e.g.}, \texttt{True} for the fact \texttt{A is True}, and \texttt{)} for the expression region) and ultimately converges at the query token (\texttt{is}). These segment-terminal tokens serve as information hubs for downstream propagation, providing the infrastructure for the staged computation observed in Section~\ref{sec:staged}. 
% Notably, \texttt{facts\_value} tokens maintain non-trivial causal importance even in middle and late layers (Figure~\ref{fig:info_converge_token}), contradicting a simple encode-then-forget model. This persistent influence suggests active re-access of source facts throughout processing.


\begin{figure}[t]
  \centering
  \begin{subfigure}[b]{0.48\linewidth}
    \includegraphics[width=\linewidth]{./figures/info_trans_1.png}
    \caption{B altered (True$\to$False)}
    \label{fig:patch_b}
  \end{subfigure}
  \hfill
  \begin{subfigure}[b]{0.48\linewidth}
    \includegraphics[width=\linewidth]{./figures/info_trans_2.png}
    \caption{A altered (True$\to$False)}
    \label{fig:patch_a}

  \end{subfigure}

  % \caption{Residual stream patching reveals information convergence.
  % Normalized logit difference after patching clean activations into the corrupt prompt \texttt{A is True, B is True, ($\neg$A or $\neg$B) is}.
  % Blue regions indicate successful restoration of the correct answer.
  % (a) Effect concentrates at token 6 (truth value), shifts to token 14 (closing parenthesis), then converges at the terminal query token.
  % (b) Causal effects emerge at segment-terminal positions: 2, 6 (truth values), 14 (expression end), and 15 (query token).}
  \caption{Residual stream patching reveals information convergence.
  Logit difference after patching clean activations into the corrupt prompt \texttt{A is True, B is True, ($\neg$A or $\neg$B) is}.
  Values are \textbf{normalized} per layer to highlight relative token importance.
  Blue regions indicate successful restoration of the correct answer.
  (a) Effect concentrates at token 6 (truth value), shifts to token 14 (closing parenthesis), then converges at the terminal query token.
  (b) Causal effects emerge at segment-terminal positions: 2, 6 (truth values), 14 (expression end), and 15 (query token).}
  \label{fig:logit_difference}
  \vspace{-0.5em}
\end{figure}

\begin{figure}[t]
    \centering
    \includegraphics[width=\columnwidth]{./figures/info-converge-onehop-8b.png}
    \caption{Token-wise information convergence in \textit{Qwen3-8B} (One-hop).
    Mean |dPD| (probability-difference shift caused by patching; see Appendix~\ref{sec:appendix_info_converge}) per token category across layer groups.
    Early layers (L0-13): \texttt{facts\_value} tokens exhibit the highest causal importance, reflecting factual encoding.
    Late layers (L24-35): \texttt{query\_token} shows a dramatic surge, indicating information convergence toward the final prediction position.
    Error bars denote standard error of the mean (SEM).}
    \vspace{-0.5em}
    \label{fig:info_converge_token}
\end{figure}

% \subsection{Information Transmission}
% In our investigation into the internal mechanisms of the Qwen3-8B model, we have identified a recurrent structural phenomenon in its forward propagation, which we term \textbf{Information Transmission}. Specifically, we observe that in responses to distinct input prompts, when the token sequence is segmented into coherent semantic intervals, the final token within each such interval often exhibits a distinctive role in information aggregation.

% This phenomenon manifests as follows: beyond a certain layer in the transformer stack, if the representation at this final-token position is altered or patched—\textit{i.e.}, replaced with the representation from another forward pass or a modified vector—the model's output undergoes significant and measurable change. Notably, this sensitivity is not uniformly distributed across all tokens in the same semantic interval; other tokens within the same segment do not demonstrate a comparable impact when similarly patched.

% Furthermore, the influence of patching this specific token persists robustly across several subsequent layers. This suggests that information pertinent to the semantic segment—or information necessary for downstream processing—tends to converge and remain resident at this positional locus over a defined depth of the network. We interpret this as evidence that the model retains and propagates interval-summary or segment-closing information at these token positions across multiple layers of computation.

% This behavior underscores a form of structural inductive bias within the transformer architecture, where boundary tokens of semantic chunks may act as information hubs, aggregating contextual meanings from preceding tokens in the interval and sustaining that aggregated signal through later layers. Our findings contribute to the growing understanding of how transformers organize and propagate linguistic and conceptual information in a layer-wise manner.

% To illustrate this phenomenon concretely, we conduct a controlled case study. We define the clean prompt as "A is True, B is False, (¬A or ¬B) is", and the corrupt prompt as "A is True, B is True, (¬A or ¬B) is". Both prompts are fed into the model, and a single forward pass is performed to obtain the corresponding activations and outputs.

% Since the residual stream contains the majority of the information propagated within each transformer block, we use the token-wise residual activations from the clean prompt to patch the corresponding positions in the corrupt prompt and then compute the resulting output logits. After each patching intervention, we measure the logit difference between the correct answer (“False”) and the incorrect answer (“True”) to quantify the causal effect of each token-level intervention.

% \begin{figure}[t]
%   \includegraphics[width=0.48\linewidth]{./figures/02.png} \hfill
%   \includegraphics[width=0.48\linewidth]{./figures/03.png}
%   \caption{Logit differences at each token position and layer were measured after patching the corrupt prompt with the \textbf{residual activations} from the clean prompt (specifically, we patch the residual stream before each layer). Higher positive values (red) indicate a stronger bias toward the incorrect answer, whereas higher negative values (green) indicate a stronger bias toward the correct answer. Note that the logit differences are normalized to the range [-1,1] for consistent visualization. Both patching experiments share the same original prompt, ``A is True, B is True, ($\neg$A or $\neg$B) is''. For the left panel, the prompt used for patching is ``A is True, B is False, ($\neg$A or $\neg$B) is'', whereas for the right panel, the patching prompt is ``A is False, B is True, ($\neg$A or $\neg$B) is''. From the resulting heatmaps, we can clearly observe both the convergence of task-relevant information and the corresponding pathways along which this information is transported across the network.}
%   \label{fig:logit difference}
% \end{figure}

% As shown in Figure~\ref{fig:logit difference} left, the patching effect at token position 6 — corresponding to the “True” token in the corrupt prompt — is particularly pronounced and persists up to approximately layer 23. Following this, token position 14, corresponding to “)” respectively, exhibit strong effects. Beyond about layer 25, the final token shows a significant impact that extends to the penultimate layer, while the remaining tokens display negligible influence. We therefore assume that the final token in each semantic region serves as a salient representation that aggregates the information of the entire region. Information from each region is progressively propagated to the subsequent regions and ultimately aggregated at the final token of the prompt.

% When we replace the clean prompt with “A is False, B is True, (¬A or ¬B) is”, we use the corresponding residual stream activations to patch the same corrupted prompt. The resulting effects are shown in the right panel of Figure~\ref{fig:logit difference}. Notably, at convergence positions 2, 6, 14, and 15, the patching induces a pronounced change in the logit difference, indicating that these positions play a dominant role in mediating the propagation and integration of task-relevant information. To further substantiate this result, we performed additional experiments using varied rule sets and two-hop reasoning — full experimental details are provided in Appendix 1.

% \subsection{Stage Calculating}
% We observed that Qwen3-8B spontaneously decomposes logical-reasoning prompts into a sequence of staged, hierarchical computations. Across prompt variants the precise layer boundaries of these stages shift slightly, but the same qualitative pattern recurs: early layers primarily process fact tokens, middle layers operate on expression (calculation / reasoning) tokens, and late layers integrate final-token information to produce the model's output. Importantly, interventions applied to the MLP outputs reveal that MLPs are a major driver of this staged computation.

% \paragraph*{Layered/staged computation} When analyzing information flow through Qwen3-8B on a range of logically structured prompts, we find that the model performs reasoning in a sequence of stages localized to different depth ranges. Early layers (approximately layers 0–10) are causally important for processing tokens that carry facts; intervening on MLP outputs at these token positions substantially alters final outputs. Middle layers (about 10–18) are similarly critical for tokens that encode intermediate expressions or transformations, and perturbations to these MLP outputs lead to large downstream effects. Finally, layers beyond about 20 are primarily involved in integrating the final token(s); perturbing the MLP output at the last-token position in these late layers also strongly affects model outputs. Across alternative prompts the numeric boundaries shift slightly, but the qualitative staging is robust.

\begin{figure*}[t]
    \centering
    % --- 第一个子图 (Figure a) ---
    \begin{subfigure}[b]{0.32\textwidth}
        \centering
        % 替换为你实际的文件名
        \includegraphics[width=\linewidth]{./figures/split_heads.pdf} 
        \caption{Splitting heads}
        \label{fig:splitting_head}
    \end{subfigure}
    \hfill % 这里的 \hfill 负责在图片之间产生弹性间距
    % --- 第二个子图 (Figure b) ---
    \begin{subfigure}[b]{0.32\textwidth}
        \centering
        % 替换为你实际的文件名
        \includegraphics[width=\linewidth]{./figures/info_heads.pdf}
        \caption{Information Transmission heads}
        \label{fig:convergence_head}
    \end{subfigure}
    \hfill
    % --- 第三个子图 (Figure c) ---
    \begin{subfigure}[b]{0.32\textwidth}
        \centering
        % 替换为你实际的文件名
        \includegraphics[width=\linewidth]{./figures/facts_heads.pdf}
        \caption{Fact-Retrieval heads}
        \label{fig:facts_head}

    \end{subfigure}
    
    % --- 总标题 ---
    \caption{Visualization of attention matrices for three distinct types of heads. The prompt used is \texttt{A is True, B is False, ($\neg$A or $\neg$B) is}. The notation $(x, y)$ denotes the $y$-th head in the $x$-th layer.
    (a) The Splitting heads (early layers) attend predominantly to the comma delimiter.
    (b) The Information Transmission heads (intermediate layers) reveal lower-triangular structures within Fact and Expression regions.
    (c) The Fact-Retrieval heads (deep layers) show tokens in the Expression Region or Query Token (\texttt{is}) attending to \texttt{True}/\texttt{False} indicators.}
    \label{fig:combined_heads}
    \vspace{-0.5em}
\end{figure*}

\begin{figure}[t]
    \centering
    \includegraphics[width=\columnwidth]{./figures_experiments/reports/heads_analysis_qwen3_8b_facts_first/taxonomy/head_taxonomy_line_chart.png}
    \caption{Layer-wise distribution of the three head roles in the saved \textit{Qwen3-8B} \texttt{facts\_first} taxonomy results, rendered from \texttt{figures\_experiments/reports/heads\_analysis\_qwen3\_8b\_facts\_first/taxonomy/head\_taxonomy\_line\_chart.png}. Splitting Heads concentrate in early layers, Transmission Heads span the middle of the network, and Fact-Retrieval Heads peak in middle-to-late layers. Appendix~\ref{appendix_heads} extends this analysis to both model scales and both prompt orders, and further shows that ablating the selected heads causes large accuracy drops while random controls do not.}
\label{fig:counts}
\vspace{-0.5em}
\end{figure}

\subsection{Fact Retrospection}
\label{sec:retro}
An unexpected finding from the preceding analysis is that fact tokens maintain
causal importance even in late layers (Figure~\ref{fig:info_converge_token} shows that \texttt{facts\_value} tokens exhibit sustained comparatively high Mean $|\text{dPD}|$ values across all layer groups).
% , contradicting a simple encode-then-forget model. 
% The Information Transmission analysis reveals an unexpected phenomenon: \textit{why do fact tokens maintain causal importance in deep layers?} Once information is transmitted to segment-terminal tokens, we would expect source positions to become causally irrelevant. Yet \texttt{facts\_value} tokens exhibit sustained comparatively high $|\text{dLD}|$ values across all layer groups (Figure~\ref{fig:info_converge_token}), contradicting a simple encode-then-forget model.
This pattern holds across one-hop and two-hop reasoning (Appendix~\ref{sec:appendix_info_converge}).
We propose the Fact Retrospection hypothesis to explain this observation: rather than processing facts once and discarding the original representations, the model actively revisits source fact tokens throughout its depth. This mechanism of active information maintenance may serve several functions: (1) Error Mitigation, where retaining access to original facts allows potential correction of misinterpretations in intermediate steps; (2) Integration Facilitation, since later reasoning steps may require renewed access to earlier facts for final synthesis; and (3) Robustness Enhancement, as redundant access to source information provides resilience against representation degradation.

% If this hypothesis is correct, we should observe attention heads that persistently attend back to fact positions in later layers.

\subsection{Specialized Attention Heads}
\label{sec:heads}
With the macroscopic computational patterns characterized in the preceding sections, a fundamental question remains: \textit{what specific attention heads implement these mechanisms?} 
We uncover \textit{Specialized Attention Heads}, \textit{i.e.}, a stable division of labor in which different heads repeatedly realize different sub-computations of logical reasoning. To test this claim, we use a three-step pipeline. First, we screen heads by clean-prompt mean ablation, ranking them by how strongly ablating that head decreases the correct-minus-incorrect probability difference. Second, we classify the retained heads by their attention patterns using boundary focus, intra-region aggregation, and fact-focus statistics, which yields three recurring roles. Third, we causally validate each role on held-out clean prompts by ablating the top-$k$ heads from that role and comparing against a \textit{Random-Outside} control sampled from heads that were not selected in the screening stage.

We group the discovered heads into three main categories, each corresponding to one of the macroscopic mechanisms identified earlier.

\paragraph{Splitting Heads} In early layers, these heads concentrate on comma tokens (Figure~\ref{fig:splitting_head}), which act as boundary markers between semantic regions. They therefore operationalize Staged Computation: by marking clause boundaries, they separate the factual premises from the expression to be solved.
\paragraph{Transmission Heads} These heads aggregate information within a semantic region and route it toward the region's terminal token. Their attention maps form lower-triangular blocks inside the Facts and Expression regions (Figure~\ref{fig:convergence_head}), matching the Information Transmission pattern identified from residual-stream patching. A subclass additionally binds variables to their truth values.
\paragraph{Fact-Retrieval Heads} These heads repeatedly attend from the query token or expression tokens back to truth-value tokens (\texttt{True}/\texttt{False}) (Figure~\ref{fig:facts_head}). Their concentration in middle-to-late layers (Figure~\ref{fig:counts}) directly supports the Fact Retrospection hypothesis: the model does not simply encode facts once, but re-accesses them during later computation.

The layer-wise organization is robust. Figure~\ref{fig:counts} shows the one-hop \textit{Qwen3-8B} case, where Splitting Heads cluster in early layers, Transmission Heads occupy the middle of the network, and Fact-Retrieval Heads remain active into later layers. Appendix~\ref{appendix_heads} shows the same coarse layout for both \textit{Qwen3-8B} and \textit{Qwen3-14B}, and for both \texttt{facts\_first} and \texttt{expr\_first} prompts, although the relative frequency of each role changes with prompt format.

More importantly, this taxonomy is causally predictive rather than merely descriptive. Across all four validation settings (\textit{Qwen3-8B}/\textit{Qwen3-14B} $\times$ \texttt{facts\_first}/\texttt{expr\_first}), ablating the top-32 heads of a single role lowers accuracy by 20.4--39.6 points for Fact-Retrieval, 19.6--27.4 points for Transmission, and 10.0--19.6 points for Splitting. Ablating top-32 heads pooled across all three roles lowers accuracy by 25.4--35.6 points. In contrast, ablating 32 \textit{Random-Outside} heads leaves accuracy at or above baseline in three of the four settings, and causes only a 0.3-point drop in the remaining case. The strongest effect appears in \textit{Qwen3-8B} with \texttt{facts\_first} prompts: ablating the top-64 Fact-Retrieval heads reduces accuracy from 96.6\% to 50.6\%, whereas ablating 64 random outside heads still leaves 91.1\% accuracy. These interventions show that the extracted heads are functionally important for solving the task, and that the three roles capture complementary parts of the reasoning mechanism rather than superficial attention motifs. Appendix~\ref{appendix_heads} reports the complete Qwen3 Step 3 accuracy tables for all strategies and all $k$ values.

\begin{table}[t]
\centering
\small
\renewcommand{\arraystretch}{0.96}
\setlength{\tabcolsep}{3.5pt}
\resizebox{\columnwidth}{!}{%
\begin{tabular}{@{} l c c c c @{}}
\toprule
\textbf{Role ($k=32$)} & \textbf{8B FF} & \textbf{8B EF} & \textbf{14B FF} & \textbf{14B EF} \\
\midrule
Base accuracy & 96.6 & 92.8 & 94.8 & 97.6 \\
Fact-Retrieval & 57.0 & 63.2 & 63.8 & 77.2 \\
Transmission & 69.2 & 73.2 & 69.2 & 76.6 \\
Splitting & 80.4 & 73.2 & 84.8 & 86.8 \\
Mixed-All & 62.0 & 67.4 & 61.4 & 62.0 \\
Random-Outside & 97.2 $\pm$ 0.4 & 92.5 $\pm$ 0.8 & 98.3 $\pm$ 0.2 & 98.2 $\pm$ 0.2 \\
\bottomrule
\end{tabular}%
}
\caption{\textbf{Qwen3 head-role accuracy validation on held-out clean prompts.} Accuracy is measured after ablating the top-$32$ heads for each role, using balanced validation samples drawn from the filtered dual-correct pool. FF = \texttt{facts\_first}; EF = \texttt{expr\_first}. Role-specific ablations consistently reduce accuracy, whereas the matched \textit{Random-Outside} control remains near baseline.}
\label{tab:qwen3_head_accuracy_main}
\end{table}

% \begin{figure}[t]
%     \centering
%     \includegraphics[width=0.8\columnwidth]{./figures/attn_heads}
%     \caption{Visualization of attention matrices for three distinct types of heads. The prompt used across all figures is ``A is True, B is False, ($\neg$A or $\neg$B) is''. The notation $(x, y)$ denotes the $y$-th head in the $x$-th layer. 
%     \textbf{first row:} The \textbf{Splitting heads} (early layers) demonstrate a pattern where tokens predominantly attend to the comma (``,'') delimiter. 
%     \textbf{second row:} The \textbf{Information Transmission heads} (intermediate layers) reveal two distinct, relatively independent lower-triangular structures within the ``Fact'' and ``Expression'' regions. 
%     \textbf{third row:} The \textbf{Fact-Retrieval heads} (deep layers) show that tokens within the ``Expression'' region (notably the final token) allocate significant attention to the ``True''/``False'' indicators.}
% \label{fig:combined_heads}
% \end{figure}


%  \begin{figure}[t]
%   \includegraphics[width=\columnwidth]{./figures/converge.pdf}
%   \caption{The attention matrix of the \textbf{Information Transmission heads} $(x, y)$ located in the intermediate layers. The visualization reveals two distinct, relatively independent lower-triangular structures within the "Fact" and "Expression" regions, respectively. Note that we use the notation $(x, y)$ to denote the $y$-th attention head situated in the $x$-th layer and the prompt used is "A is True, B is False, (¬A or ¬B) is". The later figures use the same prompt.}
%   \label{fig:converge}
% \end{figure}

%  \begin{figure}[t]
%   \includegraphics[width=\columnwidth]{./figures/split.pdf}
%   \caption{Visualization of the attention matrix for the \textbf{Splitting heads} located in the model's early layers. The heatmap demonstrates a distinct pattern wherein all input tokens predominantly attend to the comma (",") delimiter. }
%   \label{fig:split}
% \end{figure}

%  \begin{figure}[t]
%   \includegraphics[width=\columnwidth]{./figures/look.pdf}
%   \caption{The attention matrix for the \textbf{Fact-Retrieval heads} $(x, y)$ in the deep layers. The heatmap reveals that tokens within the subsequent "Expression" region allocate significant attention weights to the "True"/"False" token positions. Notably, the final token in the sequence also exhibits a pronounced attentional bias towards these truth-value indicators.}
%   \label{fig:look}
% \end{figure}

\section{Discussion}
\subsection{Implications of the Findings}

Our results suggest that propositional logical reasoning in transformers is not implemented by a monolithic circuit but rather organized as a multi-phase computational strategy: early-layer encoding, middle-layer information routing, and late-layer decision making. This staged organization persists across 11 rule categories, two reasoning depths, and two model scales, which points to a general architectural property rather than a task-specific shortcut.

A notable finding is Fact Retrospection, where models repeatedly re-access premise information across multiple layers rather than encoding it once. This recurring retrieval pattern is qualitatively different from the single-shot factual recall observed in knowledge localization studies \citep{meng2023locatingeditingfactualassociations}, and may reflect the additional computational demands of composing multiple premises under a logical rule. Whether this distributed access pattern extends to other forms of compositional reasoning (e.g., arithmetic, spatial) remains an open question.

The staged computation we observe in the no-CoT setting also raises an interesting question about chain-of-thought prompting: if models already implement internal staging, CoT may succeed not by introducing staged reasoning \textit{de novo}, but by aligning the model's internal computation with an explicit external structure. Testing whether CoT reshapes or merely amplifies the patterns we identify could clarify the mechanistic role of intermediate tokens.

\subsection{Methodological Considerations}

Several methodological choices in our analysis warrant discussion. First, our head taxonomy (Splitting, Transmission, Fact-Retrieval) is derived from attention pattern clustering and validated through group ablation (Table~\ref{tab:qwen3_head_accuracy_main}), but we have not yet performed targeted ablation of individual heads within each type or matched random controls of equal size. Such experiments would strengthen the causal claims by distinguishing necessary heads from redundant ones within each category.

Second, our intervention suite mixes two families of operations: clean-to-corrupt activation patching for residual-stream analyses, and mean-ablation replacements for region- and head-level analyses. The latter are more conservative than zeroing activations, but they can still introduce distributional shifts because the replacement vectors need not lie on the model's natural trajectory for the intervened sample. Further work should compare within-sequence means, dataset means, and cleaner causal baselines across all branches.

Third, our primary metric is the logit difference between the correct and incorrect tokens. While this captures the model's relative preference, it may miss cases where the correct token's absolute probability changes substantially without altering the logit difference. Supplementing with the probability of the correct token would provide a more complete picture.

% \subsection{From Circuit Discovery to Mechanism Characterization}

% Our approach departs from the circuit-discovery paradigm \citep{elhage2021mathematicalframeworktransformer, wang2022interpretabilitywildcircuitindirect} that asks ``which heads are necessary?'' We instead ask ``what computational strategies does the model employ?''

% This shift is motivated by two observations. First, circuit faithfulness is difficult to establish: multiple circuits may implement the same function, and ablating one may simply shift computation to another. Second, circuits often fail to generalize across models, limiting their scientific value for understanding \emph{how transformers reason} rather than \emph{how a particular model is wired}.

% Our findings (Staged Computation, Information Transmission, Fact Retrospection, Specialized Attention Heads) describe patterns at an abstraction level that may generalize across architectures, aligning with cognitive science's distinction between algorithmic and implementational analysis.


% \subsection{Lazy Reasoning: Information Flow vs. Logical Computation}

% Our analysis primarily traces information flow patterns, but an important distinction exists between information flow and logical computation. Our directional patching analysis (Appendix~\ref{sec:lazy_reasoning}, Figure~\ref{fig:lazy_reasoning}) suggests that the actual logical computation may be deferred to late layers at the \textit{Query Token}, consistent with a lazy reasoning strategy where facts are encoded early, routed through middle layers, and combined only at the final decision point.



% \paragraph{Other Attention Head Patterns.}
\subsection{Other Attention Head Patterns.}
% \label{sec:other_heads}
Beyond the above three primary head types, we observe additional patterns including: (1) \textit{Idle Heads} that consistently attend to the first token regardless of content; (2) \textit{Information Binding Heads} in early layers (specifically, L0-10) that associate entities with truth values; (3) \textit{Self-Processing Heads} exhibiting diagonal attention where tokens primarily attend to themselves; and (4) \textit{Expression Processing Heads} in L18-22 that concentrate on expression tokens before attention convergence. These patterns, while not causally validated, suggest additional functional diversity. More Details are provided in Appendix~\ref{sec:appendix_other_heads}.



% \paragraph{Functional Coupling in Specialized Heads.}
\subsection{Functional Coupling in Specialized Heads.}
We observe that individual attention heads often exhibit multiple functional roles simultaneously. For example, a head may show Splitting behavior at delimiter positions while also contributing to Transmission within semantic regions; similarly, Fact-Retrieval patterns frequently co-occur with Transmission patterns in middle-layer heads. This functional coupling suggests that our categorization of Specialized Attention Heads represents idealized abstractions rather than mutually exclusive classes. Such multi-functionality likely reflects parameter efficiency: the model implements diverse reasoning sub-processes by reusing the same heads across different computational roles. This observation aligns with the superposition hypothesis in neural networks, where individual components participate in multiple overlapping circuits. Details are provided in Appendix~\ref{sec:appendix_func_couple}.
% Future work should develop metrics to quantify the degree of functional overlap and investigate whether multi-functional heads differ in causal importance from specialized ones.

% \subsection{Lazy Reasoning: Information Flow vs. Logical Computation}

% \paragraph{Lazy Reasoning.}

% Our analysis primarily traces information flow patterns, but an important distinction exists between information flow and logical computation. To investigate where actual computation occurs, we analyze \text{dLD} rather than $|\text{dLD}|$ to differentiate restoration and interferance: positive \text{dLD} indicates that patching restores the clean prediction, while negative \text{dLD} indicates interference with already-correct reasoning. This analysis (Appendix~\ref{sec:lazy_reasoning}, Figure~\ref{fig:lazy_reasoning}) reveals a notable asymmetry at the Query Token: when the model predicts differently on corrupted vs. clean inputs, patching in late layers produces strong positive effects; when predictions already match, patching shows opposite signs across layers, \textit{i.e.}, positive in middle layers but negative in late layers. This sign reversal suggests a \textit{lazy reasoning} hypothesis: middle layers primarily transmit information to the Query Token, while actual logical computation occurs in late layers. 
% % Further experiments are needed to validate this hypothesis.



% Our analysis primarily traces information flow patterns, but an important distinction exists between information flow and logical computation. To investigate where actual computation occurs, we employ signed \text{dLD} rather than absolute values: positive \text{dLD} indicates that patching restores the correct answer, while negative \text{dLD} indicates interference with already-correct reasoning. This signed analysis (Appendix~\ref{sec:lazy_reasoning}, Figure~\ref{fig:lazy_reasoning}) reveals a notable asymmetry: when the model predicts differently on corrupted vs. clean inputs, patching the Query Token in late layers produces strong positive effects (\textit{i.e.}, restoration); when predictions already match, the same patching produces negative effects (\textit{i.e.}, interference). This differential pattern suggests that actual logical computation is deferred to late layers at the \textit{Query Token}, consistent with a lazy reasoning strategy where facts and expressions are transferred to the \textit{Query Token} through layers, with logical computation deferred to the final decision point.



\section{Conclusion and Future Work}
In this paper, we present a mechanistic analysis of propositional logical reasoning in \textit{Qwen3} models, adopting a top-down analytical strategy 
from
% that effectively bridges the gap between 
macroscopic patterns 
to
% and
microscopic implementations.
Our investigation uncovers a coherent computational architecture: the model organizes reasoning into temporally distinct phases,
% (\textit{Staged Computation}), 
routes information through semantic boundary tokens,
% (\textit{Information Transmission}), 
maintains active re-access to source facts throughout processing, 
% (\textit{Fact Retrospection}), 
and implements these strategies through specialized attention heads.
% (\textit{Specialized Attention Heads}). 
These four phenomena form an interlocking system where each mechanism supports and validates the others. Furthermore, these patterns generalize across model scales, rule types and reasoning depths, suggesting fundamental computational strategies rather than task-specific shortcuts. By characterizing \textit{how} models reason rather than \textit{which components} are necessary, 
% we contribute mechanistic evidence to the debate on whether LLMs implement genuine algorithmic processes or rely solely on pattern matching.
we contribute mechanistic evidence to the hypothesis that LLMs implement algorithmic processes rather than rely solely on pattern matching.


\paragraph{Future Directions}
Our findings open several promising avenues for future research. First, extending the analysis to deeper reasoning chains (\textit{i.e.}, three-hop and beyond) would test whether the identified mechanisms scale gracefully under increased complexity or encounter qualitative bottlenecks. Second, extending the cross-model comparison beyond the currently analyzed \textit{Qwen}, \textit{Llama}, and \textit{Mistral} families would establish how broadly these mechanisms transfer across different pre-training regimes and architectural choices. 
% Third, applying similar mechanistic techniques to other logic systems such as first-order logic or modal logic could reveal whether propositional-level mechanisms serve as building blocks for more complex reasoning. 
% Third, the mechanisms we identify (particularly Fact Retrospection and Specialized Attention Heads) provide potential targets for interventional approaches: selectively enhancing or suppressing specific heads may improve reasoning accuracy or interpretability. 
% Finally, the lazy reasoning pattern observed in our directional patching analysis (Appendix~\ref{sec:lazy_reasoning}) warrants deeper investigation into computational efficiency tradeoffs, as deferring computation to late layers may reflect an optimization strategy with implications for reasoning speed and robustness.
Finally, analyzing training dynamics would reveal how and when these mechanisms emerge during pretraining, \textit{i.e.}, whether they develop gradually, appear suddenly via phase transitions (\textit{e.g.}, grokking), or require specific training data distributions. 



\section*{Limitations}
Our analysis has several methodological limitations. 
First, while we characterize attention head specialization in detail, the internal computations within MLP layers remain largely unexplored. Our MLP patching experiments reveal when MLPs are causally important, but not what specific transformations they perform. Understanding how MLPs encode logical rules, bind variables to truth values, or compute Boolean operations represents a significant open challenge.
% Second, activation patching establishes causal sufficiency but not necessity??: alternative circuits may implement the same function, and ablating one pathway may simply shift computation to another. Our identified mechanisms represent sufficient components for reasoning, not necessarily the only components involved.
Second, our analysis characterizes the presence of mechanisms in pre-trained models but does not explain how or when these mechanisms emerge during training. Whether these patterns arise from specific training data, emerge gradually through pretraining, or appear suddenly via phase transitions remains unknown. Finally, while we interpret our findings as evidence for structured computation, we cannot definitively rule out that these patterns reflect sophisticated statistical shortcuts that merely approximate genuine logical algorithms.



% Bibliography entries for the entire Anthology, followed by custom entries
%\bibliography{anthology,custom}
% Custom bibliography entries only
\bibliography{custom}

% \clearpage

\appendix
\section*{Contents of Appendix}
\begin{itemize}
    \item[\ref{sec:appendix_setup}] \hyperref[sec:appendix_setup]{Experimental Setup Details}
    \begin{itemize}
        \item[\ref{sec:appendix_hardware}] \hyperref[sec:appendix_hardware]{Hardware and Software}
        \item[\ref{sec:appendix_implementation}] \hyperref[sec:appendix_implementation]{Implementation Details}
        \item[\ref{sec:appendix_reproducibility}] \hyperref[sec:appendix_reproducibility]{Reproducibility}
        \item[\ref{sec:appendix_dataset}] \hyperref[sec:appendix_dataset]{Dataset Details}
    \end{itemize}
    \item[\ref{sec:appendix_results}] \hyperref[sec:appendix_results]{More Results}
	    \begin{itemize}
	        \item[\ref{sec:appendix_performance}] \hyperref[sec:appendix_performance]{Behavioral Performance on PropLogic-MI}
	        \item[\ref{sec:appendix_prompt_order}] \hyperref[sec:appendix_prompt_order]{Prompt-Order Intervention: \texttt{expr\_first}}
	        \item[\ref{sec:appendix_stage}] \hyperref[sec:appendix_stage]{Staged Computation}
	        \item[\ref{sec:appendix_region_variants}] \hyperref[sec:appendix_region_variants]{Additional Region-Level Ablations}
	        \item[\ref{sec:appendix_info_converge}] \hyperref[sec:appendix_info_converge]{Information Transmission}
	        \item[\ref{appendix_heads}] \hyperref[appendix_heads]{Specialized Attention Heads: Robustness and Validation}
	        \item[\ref{sec:appendix_llama}] \hyperref[sec:appendix_llama]{Cross-Model Extensions: Llama-3.1-8B-Instruct and Mistral-7B-Instruct-v0.1}
	        \item[\ref{sec:appendix_full_accuracy}] \hyperref[sec:appendix_full_accuracy]{Complete Accuracy Tables}
	        \item[\ref{sec:failure_analysis}] \hyperref[sec:failure_analysis]{Failure Case Analysis}
	    \end{itemize}
        \item[\ref{sec:appendix_discussion}] \hyperref[sec:appendix_discussion]{Discussion}
    \begin{itemize}
        \item[\ref{sec:appendix_other_heads}] \hyperref[sec:appendix_other_heads]{Other Attention Head Patterns}
        \item[\ref{sec:appendix_func_couple}] \hyperref[sec:appendix_func_couple]{Functional Coupling in Specialized Heads}
        % \item[\ref{sec:lazy_reasoning}] \hyperref[sec:lazy_reasoning]{Lazy Reasoning}
    \end{itemize}
\end{itemize}


\section{Experimental Setup Details}
\label{sec:appendix_setup}

\subsection{Hardware and Software}
\label{sec:appendix_hardware}            

\paragraph{Hardware.}
All experiments are conducted on NVIDIA H100 GPUs with 96GB memory. The large memory capacity allows us to process the entire \textit{Qwen3-14B} model in float16 precision without requiring model parallelism.

\paragraph{Software Environment.}
We use PyTorch 2.0 with CUDA 11.8. Activation patching experiments are implemented using the TransformerLens library \citep{nanda2022transformerlens}, which provides convenient hooks for accessing and modifying intermediate activations in transformer models.
\begin{figure}[!h]
\centering
\includegraphics[width=\columnwidth]{./figures/method_metric.pdf}
% \caption{\textbf{Activation patching procedure.} We perform a forward pass on the clean input $x_{\text{clean}}$ to cache activations at component $(l,i)$. During execution of the corrupted input $x_{\text{corrupt}}$, we replace the activation at the targeted component with the cached clean activation. The resulting shift in logit difference $\Delta \text{Logits}$ quantifies the causal importance of that component for steering the model toward the correct prediction.}
\caption{Activation patching procedure. We focus on the logits difference between \texttt{True} token and \texttt{False} token in $Y_{patched}$ output logits.}
\vspace{-0.5em}
\label{fig:activate_patching}
\end{figure}
\subsection{Implementation Details}
\label{sec:appendix_implementation} 

\paragraph{General intervention protocol.}
All three analysis branches operate on paired prompts, namely a clean prompt $x_{\mathrm{clean}}$ and a minimally perturbed corrupted prompt $x_{\mathrm{corr}}$. Unless noted otherwise, we first run the model on $x_{\mathrm{clean}}$ to cache the relevant activations, then run the model on $x_{\mathrm{corr}}$ to obtain a baseline prediction, and finally re-run the model with a targeted intervention. For the final answer position, we define a two-way probability difference over the True/False vocabulary:
\begin{equation}
\mathrm{PD}(x) = p_{\theta}(\texttt{True}\mid x) - p_{\theta}(\texttt{False}\mid x),
\end{equation}
where the probabilities are computed by a softmax restricted to the two answer tokens. We then use the label-aligned version
\begin{equation}
\mathrm{PD}_{\mathrm{label}}(x, y) =
\begin{cases}
\mathrm{PD}(x), & y = \texttt{True},\\
-\mathrm{PD}(x), & y = \texttt{False},
\end{cases}
\end{equation}
so larger values always indicate stronger support for the correct answer. For an intervention producing logits $\tilde{x}$, the main shift metric is
\begin{equation}
\mathrm{dPD} = \mathrm{PD}_{\mathrm{label}}(\tilde{x}, y) - \mathrm{PD}_{\mathrm{label}}(x, y).
\end{equation}
All reported values are computed from this probability-difference form rather than raw logit subtraction.

\paragraph{Sampling and filtering.}
For mechanistic analyses we operate on the dual-correct subset, \textit{i.e.}, samples for which the model answers both the clean and corrupted prompts correctly. Within each experiment, rows are further balanced by rule type to avoid domination by high-frequency rules. The main Qwen3 runs use balanced per-rule budgets of up to 3000 rows for MLP region ablation, 1000 rows for token patching, 2000 rows for head discovery/refinement, 1000 rows for head-validation PD curves, and 500 rows for the Step 3 accuracy tables.

\paragraph{Layer partition.}
For token- and region-level summaries we partition layers into three contiguous stages. In the token analysis, we use fixed boundaries Early = $[0,13]$, Middle = $[14,23]$, and Late = $[24, L-1]$, where $L$ is the number of transformer layers. In the MLP analysis, we instead split the full depth into three approximately equal contiguous bands using \texttt{numpy.array\_split}, again denoted Early, Middle, and Late.

\paragraph{MLP region ablation.}
This protocol underlies the Staged Computation results in §\ref{sec:staged} and the main-text visualization in Figure~\ref{fig:stage}.
The MLP analysis does not patch clean activations into a corrupted run. Instead, for a chosen semantic region $R$ and layer $l$, we run the clean prompt and replace the MLP outputs for all positions $i \in R$ by the within-sequence mean MLP vector at that layer:
\begin{equation}
\tilde{m}^{(l)}_i = \frac{1}{T}\sum_{j=1}^{T} m^{(l)}_j, \qquad i \in R,
\end{equation}
where $T$ is the sequence length and $m^{(l)}_j$ is the original \texttt{mlp\_out} vector. The effect is measured by
\begin{equation}
R_{\mathrm{norm}}^{(l)}(R) = \frac{|\mathrm{dPD}^{(l)}(R)|}{\max\left(p_{\mathrm{sel}}(x), \varepsilon\right)},
\end{equation}
where $p_{\mathrm{sel}}(x)=\max\{p(\texttt{True}\mid x), p(\texttt{False}\mid x)\}$ on the unablated prompt and $\varepsilon=10^{-6}$. We report both this absolute score and its signed counterpart. Regions are defined by character-span alignment and then mapped back to token indices: \textbf{Facts Region}, \textbf{Expression Region}, and \textbf{Query Region}. The current implementation folds the older CLI aliases \texttt{constrain\_region} and \texttt{terminal\_token} into \texttt{query\_region} for backward compatibility. To summarize stage structure, we additionally compute: 
\begin{equation}
\mathrm{BMI}_b = \frac{1}{|b|}\sum_{l \in b} R_{\mathrm{norm}}^{(l)}, \qquad
\mathrm{BCR}_b = \frac{\sum_{l \in b} R_{\mathrm{norm}}^{(l)}}{\sum_{l=1}^{L} R_{\mathrm{norm}}^{(l)}},
\end{equation}
and the signed band index
\begin{equation}
\mathrm{SBI}_b = \frac{1}{|b|}\sum_{l \in b} \widetilde{R}_{\mathrm{norm}}^{(l)},
\end{equation}
where $b \in \{\text{Early},\text{Middle},\text{Late}\}$ and $\widetilde{R}_{\mathrm{norm}}^{(l)}$ denotes the signed score.

\paragraph{Attention-region and joint attention+MLP ablations.}
The current codebase adds two region-level companions to the main MLP branch. In \textbf{attention-region ablation}, we replace \texttt{attn\_out} for all positions in a selected region by the within-sequence mean attention-output vector at that layer. In \textbf{joint attention+MLP ablation}, we apply the same mean replacement to both \texttt{attn\_out} and \texttt{mlp\_out} inside the same region during one forward pass. Both branches reuse the same region parser, the same $\mathrm{dPD}$-based scoring, and the same BMI/BCR/SBI summaries as the MLP branch. Their results are summarized in Appendix~\ref{sec:appendix_region_variants}.

\paragraph{Token-wise residual patching.}
This procedure supports the Information Transmission analysis in §\ref{sec:info}, including the aggregated token-category result in Figure~\ref{fig:info_converge_token} and its appendix-scale extensions in Appendix~\ref{sec:appendix_info_converge}.
The current token analysis is a two-stage pipeline. In the raw stage, we patch \texttt{resid\_pre} one layer and one token at a time from the clean prompt into the corrupted prompt, and save the resulting layer-by-token patching tensor for downstream analysis. In the released Qwen3 runs, this raw stage uses the dual-correct subset, \texttt{strict\_length\_match=True}, and a balanced sample budget of up to $1000$ examples per setting. Because chat formatting may prepend template tokens, raw prompt positions are first aligned to model-input positions using tokenizer offset mappings; if offset-based alignment fails, the implementation falls back to a contiguous body-span mapping inside the model input. For a sample with $T$ aligned raw-token positions, we obtain a matrix $\mathrm{dPD}_{l,i}$ with $l \in [0,L-1]$ and $i \in [1,T]$.

The raw stage also produces a simple token grouping used for quick diagnostics and plotting. These simple categories are \texttt{facts\_value}, \texttt{query\_token}, \texttt{expr\_last}, \texttt{expression\_token}, and \texttt{others}. The refined stage then reclassifies each token into a more detailed taxonomy using prompt-order-aware string heuristics:
\texttt{facts\_value}, \texttt{variable\_in\_facts}, \texttt{query\_token}, \texttt{variable\_in\_expr}, \texttt{operator}, \texttt{expr\_last}, optional \texttt{derived\_assignment}, and \texttt{others}.
Here, the query token is the first \texttt{is} in \texttt{expr\_first} prompts and the last \texttt{is} in \texttt{facts\_first} prompts; the expression span is resolved relative to that query position; and for two-hop prompts the intermediate assignment span is separately identified so that it can be marked as \texttt{derived\_assignment}. This prompt-order-sensitive classifier is important because the same lexical tokens occupy different structural roles under \texttt{facts\_first} and \texttt{expr\_first}.

For both the simple and refined stages, we partition layers into Early = $[0,13]$, Middle = $[14,23]$, and Late = $[24,L-1]$, where $L$ is the model depth ($36$ for Qwen3-8B and $40$ for Qwen3-14B). The raw stage writes sample-level patching results to \texttt{patching\_results.pkl}, and the refined stage re-reads these tensors to compute category-level summaries. This separation lets us revise token taxonomies and aggregation rules without rerunning the expensive layer-by-token patching pass.

For each category $c$ and stage $s$, we report two statistics. The first measures total category contribution within a sample:
\begin{equation}
S_{c}^{(s)} = \sum_{i \in c} \frac{1}{|s|}\sum_{l \in s} |\mathrm{dPD}_{l,i}|.
\end{equation}
The second normalizes by the number of tokens in the category:
\begin{equation}
M_{c}^{(s)} = \frac{1}{|c|}\sum_{i \in c} \frac{1}{|s|}\sum_{l \in s} |\mathrm{dPD}_{l,i}|.
\end{equation}
We average these sample-level quantities across the dataset and report means and standard errors. Intuitively, $S_{c}^{(s)}$ captures aggregate influence mass, while $M_{c}^{(s)}$ measures average per-token importance. In addition, the raw stage reports a sample-averaged heatmap over categories and layers, making it possible to inspect how the same token class changes its causal profile across depth before any stage-wise aggregation.

\paragraph{Attention-head discovery and validation.}
This is the methodological backbone of the Specialized Attention Heads analysis in §\ref{sec:heads}, the role-layout figure in Figure~\ref{fig:counts}, and the robustness/validation results in Appendix~\ref{appendix_heads}.
The head analysis consists of three steps. \textbf{Step 1} is a necessity screen on clean prompts: we ablate one head at a time by replacing its output $z_{l,h}$ with the mean output of that head computed over a balanced clean sample set. Heads are ranked by
\begin{equation}
\mathrm{Sel}(l,h) = -\mathbb{E}[\mathrm{dPD}_{\mathrm{label}}(l,h)],
\end{equation}
so larger values correspond to heads whose ablation most strongly reduces correct-answer support. We first run a small probe set over all heads, keep the top candidates globally together with the top heads from each layer, and then re-evaluate only this candidate pool on a larger balanced set.
\textbf{Step 2} classifies the retained heads using attention-pattern statistics measured on clean prompts: fact-focus ratio, boundary-focus ratio, intra-clause ratio, attention entropy, top-1 mass, and a split-half stability score. After feature standardization, each head receives rule-based scores for three roles,
\begin{equation}
\begin{split}
\mathrm{Score}_{\mathrm{FR}}
&= f - 0.20\,i - 0.15\,b \\
&\quad + 0.20\,t - 0.10\,e + 0.08\,s,
\end{split}
\end{equation}
\begin{equation}
\begin{split}
\mathrm{Score}_{\mathrm{SP}}
&= b - 0.15\,i - 0.10\,f \\
&\quad + 0.15\,t - 0.08\,e + 0.10\,s,
\end{split}
\end{equation}
\begin{equation}
\begin{split}
\mathrm{Score}_{\mathrm{TR}}
&= i - 0.20\,b - 0.10\,f \\
&\quad + 0.18\,t - 0.08\,e + 0.12\,s,
\end{split}
\end{equation}
where $f,b,i,e,t,s$ denote the standardized feature values for fact-focus, boundary-focus, intra-clause concentration, entropy, top-1 mass, and stability, respectively. The winning label is assigned as Fact-Retrieval, Splitting, or Transmission.
\textbf{Step 3} validates these roles on held-out clean prompts by ablating the top-$k$ heads from each role or role combination, with $k \in \{1,2,4,8,16,32,64\}$. We report both accuracy and relative PD degradation:
\begin{equation}
\rho_{\mathrm{abs}} = \frac{|\mathrm{dPD}|}{|\mathrm{PD}_{\mathrm{clean}}|}, \qquad
\rho_{\mathrm{signed}} = \frac{\mathrm{dPD}}{\mathrm{PD}_{\mathrm{clean}}},
\end{equation}
aggregated over balanced validation samples. The matched \textit{Random-Outside} control is constructed by sampling heads from layers used in the validation stage but outside the Step 1 selected pool.

\paragraph{Patching granularities.}
Across the full project we therefore intervene at three levels of granularity:
\begin{itemize}
    \item \textbf{Residual stream}: patch $\mathrm{resid\_pre}_{l,i}$ at layer $l$ and token position $i$.
    \item \textbf{Attention heads}: replace the output $z_{l,h}$ of head $h$ at layer $l$ by its dataset mean, either at the query token only or at all token positions.
    \item \textbf{Region-level mean ablation}: replace \texttt{mlp\_out}, \texttt{attn\_out}, or both inside a semantic region by the corresponding within-sequence mean vector.
\end{itemize}
\begin{figure*}[t]
    \centering
    \begin{subfigure}[b]{0.31\textwidth}
        \centering
        \includegraphics[width=\linewidth]{./figures_experiments/reports/mlp_analysis/mlp_14b_one_hop/facts_region/mlp_facts_region_scores.png}
        \caption{\textbf{Facts (One-hop)}}
        \label{fig:facts_14B_onehop}
    \end{subfigure}
    \hfill
    \begin{subfigure}[b]{0.31\textwidth}
        \centering
        \includegraphics[width=\linewidth]{./figures_experiments/reports/mlp_analysis/mlp_14b_one_hop/expression_region/mlp_expression_region_scores.png}
        \caption{\textbf{Expression (One-hop)}}
        \label{fig:expression_14B_onehop}
    \end{subfigure}
    \hfill
    \begin{subfigure}[b]{0.31\textwidth}
        \centering
        \includegraphics[width=\linewidth]{./figures_experiments/reports/mlp_analysis/mlp_14b_one_hop/query_region/mlp_query_region_scores.png}
        \caption{\textbf{Query (One-hop)}}
        \label{fig:query_14B_onehop}
    \end{subfigure}
    
    \vspace{0.8em}
    
    \begin{subfigure}[b]{0.31\textwidth}
        \centering
        \includegraphics[width=\linewidth]{./figures_experiments/reports/mlp_analysis/mlp_14b_two_hop/facts_region/mlp_facts_region_scores.png}
        \caption{\textbf{Facts (Two-hop)}}
        \label{fig:facts_14B_twohop}
    \end{subfigure}
    \hfill
    \begin{subfigure}[b]{0.31\textwidth}
        \centering
        \includegraphics[width=\linewidth]{./figures_experiments/reports/mlp_analysis/mlp_14b_two_hop/expression_region/mlp_expression_region_scores.png}
        \caption{\textbf{Expression (Two-hop)}}
        \label{fig:expression_14B_twohop}
    \end{subfigure}
    \hfill
    \begin{subfigure}[b]{0.31\textwidth}
        \centering
        \includegraphics[width=\linewidth]{./figures_experiments/reports/mlp_analysis/mlp_14b_two_hop/query_region/mlp_query_region_scores.png}
        \caption{\textbf{Query (Two-hop)}}
        \label{fig:query_14B_twohop}
    \end{subfigure}
    
    \caption{Region-wise MLP mean-ablation scores in \textit{Qwen3-14B}, rendered directly from the saved outputs in \texttt{figures\_experiments/reports/mlp\_analysis}. Top row: one-hop \texttt{facts\_first}. Bottom row: two-hop \texttt{facts\_first}. Across both depths, the Facts Region is most influential in early layers, the Expression Region peaks in the middle of the network, and the Query Region dominates middle-to-late integration. This figure therefore extends the staged computation analysis in the main text to a larger model while confirming that the overall region ordering remains stable across reasoning depth.}
    \label{fig:staged_computation_14B}
\end{figure*}


All forward passes are performed with \texttt{torch.no\_grad()} to reduce memory consumption and disable gradient computation. We use \texttt{float16} precision for all activations and model weights.

\paragraph{Computational Cost.}
A single residual stream patching experiment (36 layers $\times$ approximately 16 token positions) takes approximately 30 seconds on an H100 GPU for \textit{Qwen3-8B}. Attention head patching (36 layers $\times$ 32 heads) requires approximately 2 minutes. The total computation time for all experiments presented in this paper is approximately 50 GPU-hours.

\subsection{Reproducibility}
\label{sec:appendix_reproducibility}       

Code and data will be made available upon publication. Our implementation is based on publicly available models from Hugging Face and the open-source TransformerLens library, ensuring that all experiments can be reproduced by the community.

\begin{table*}[t]
\centering
\small
\renewcommand{\arraystretch}{0.9}
\setlength{\tabcolsep}{6pt}
\begin{tabularx}{\textwidth}{@{} l c X @{}}
\toprule
% \textbf{Category} & \textbf{Formal Definitions (Duals)} & \textbf{Example Prompt Templates} \\
\textbf{Category} & \textbf{Formal Definitions} & \textbf{Example Prompt Templates (One-Hop)} \\
\midrule
% --- Basic Laws ---
Identity
    & $P \land \top \equiv P, \quad P \lor \bot \equiv P$
    & \texttt{"A is T, A and T is"}, \quad \texttt{"A is F, A or F is"} \\
Domination
    & $P \land \bot \equiv \bot, \quad P \lor \top \equiv \top$
    & \texttt{"A is T, A and F is"}, \quad \texttt{"A is F, A or T is"} \\
Idempotent
    & $P \land P \equiv P, \quad P \lor P \equiv P$
    & \texttt{"A is T, A and A is"}, \quad \texttt{"A is F, A or A is"} \\
\midrule
% --- Negation Laws ---
Dbl. Negation
    & $\lnot(\lnot P) \equiv P$
    & \texttt{"A is T, ($\lnot$($\lnot$A)) is"} \\
Excluded Mid.
    & $P \lor \lnot P \equiv \top$
    & \texttt{"A is T, A or $\lnot$A is"} \\
Contradiction
    & $P \land \lnot P \equiv \bot$
    & \texttt{"A is F, A and $\lnot$A is"} \\
\midrule
% --- Multi-Variable Laws ---
Commutative
    & $P \land Q \equiv Q \land P$
    & \texttt{"A is T, B is F, A and B is"} \\
Associative
    & $(P \land Q) \land R \equiv P \land (Q \land R)$
    & \texttt{"A is T, B is T, C is F, (A and B) and C is"} \\
Distributive
    & $P \land (Q \lor R) \equiv (P \land Q) \lor (P \land R)$
    & \texttt{"A is T, B is F, C is T, A and (B or C) is"} \\
De Morgan
    & $\lnot(P \land Q) \equiv \lnot P \lor \lnot Q$
    & \texttt{"A is T, B is F, ($\lnot$A or $\lnot$B) is"} \\
Absorption
    & $P \land (P \lor Q) \equiv P$
    & \texttt{"A is T, B is F, A and (A or B) is"} \\
\bottomrule
\end{tabularx}
% \caption{\textbf{Rule coverage.} We cover 11 fundamental categories. For rules with dual forms (Identity, Domination, Idempotent), both formulas and example templates are shown with horizontal spacing. For each rule covers its formal Boolean algebra definition and a corresponding template. Variables \{A, B, C\} are instantiated with truth values (True/False). T/F denotes True/False.}
\caption{\textbf{Rule coverage.} \textit{PropLogic-MI} systematically covers 11 fundamental rule categories. Each rule includes its formal Boolean algebra definition and a corresponding template. Variables \{\texttt{A, B, ..., Z}\} are instantiated with truth values (\texttt{True}/\texttt{False}). T/F denotes \texttt{True}/\texttt{False}.}
\vspace{-1.0em}
\label{tab:rule_coverage}
\end{table*}

\begin{table}[t]
\centering
\small
\renewcommand{\arraystretch}{0.9}
\setlength{\tabcolsep}{4pt}
\begin{tabularx}{\columnwidth}{@{}lX@{}}
\toprule
\textbf{Depth} & \textbf{Prompt Template \& Reasoning Process} \\
\midrule
\multirow{3.5}{*}{\textbf{One-Hop}}
& \texttt{"A is True, B is False, (A and B) is"} \newline
  \color{gray}\textit{\footnotesize $\hookrightarrow$ Logic: $A \land B$} \\
\addlinespace[0.4em]
& \texttt{"A is True, ($\lnot$($\lnot$A)) is"} \newline
  \color{gray}\textit{\footnotesize $\hookrightarrow$ Logic: $\lnot(\lnot A) \equiv A$} \\
\addlinespace[0.4em]
& \texttt{"A is True, B is False, ($\lnot$A or $\lnot$B) is"} \newline
  \color{gray}\textit{\footnotesize $\hookrightarrow$ Logic: De Morgan ($\lnot A \lor \lnot B$)} \\
\midrule
\multirow{5}{*}{\textbf{Two-Hop}}
& \texttt{"A is T, B is A and T, C is F, B and C is"} \newline
  \color{gray}\textit{ \footnotesize $\hookrightarrow$ Chain: $(A \land \top) \to B; \quad (B \land C) \to \text{Ans}$} \\
\addlinespace[0.4em]
& \texttt{"A is T, B is F, C is $\lnot$(A and B), D is T, C or D is"} \newline
  \color{gray}\textit{\footnotesize $\hookrightarrow$ Chain: $\text{NAND}(A,B) \to C; \quad (C \lor D) \to \text{Ans}$} \\
\addlinespace[0.4em]
& \texttt{"A is T, B is F, C is A and B, D is T, C or D is"} \newline
  \color{gray}\textit{\footnotesize $\hookrightarrow$ Chain: $(A \land B) \to C; \quad (C \lor D) \to \text{Ans}$} \\
\bottomrule
\end{tabularx}
\caption{\textbf{Probing prompts across reasoning depths.} We probe the model with queries requiring direct rule application (One-Hop) versus those necessitating intermediate latent variable derivation (Two-Hop).}
\label{tab:reasoning_depth}
\end{table}

\subsection{Dataset Details}
\label{sec:appendix_dataset}

\textit{PropLogic-MI} is released in two prompt-order variants, namely \path{dataset/proplogic_mi.jsonl} for \texttt{facts\_first} and \path{dataset/proplogic_mi_expr_first.jsonl} for \texttt{expr\_first}. Both corpora are generated with the same seed and the same underlying logical rows, differing only in prompt order. Each file contains exactly $20000$ unique clean/corrupted pairs with no duplicate symbolic prompts. All released rows use the \texttt{answer\_suffix} ending, \textit{i.e.}, they terminate with ``Answer with one word only: True or False.'' The symbolic compact prompt is the one used throughout the main experiments, while each row also stores symbolic verbose variants and semi-natural expression strings for auxiliary analyses.

\paragraph{Rule inventory and global balance.}
The dataset spans the 11 rule templates listed in Table~\ref{tab:rule_coverage}: identity, domination, idempotent, double negation, excluded middle, contradiction, commutative, associative, distributive, De Morgan, and absorption. These rules are grouped into three categories already stored in the JSONL rows: \texttt{basic\_laws}, \texttt{negation\_laws}, and \texttt{multi\_variable\_laws}. The final $20000$ released rows are nearly rule-balanced overall: each rule contributes either $1818$ or $1819$ examples, so the rule-balance span is only one example.

\paragraph{Depth allocation in the released corpora.}
Generation is requested with a nominal $50/50$ one-hop/two-hop split, but the final released files are constrained by uniqueness under the corruption semantics. As recorded in the generation logs, the maximum feasible one-hop allocation under these constraints is $9523$ examples, so the final corpora contain $9523$ one-hop rows and $10477$ two-hop rows. This adjustment is deterministic and identical for the \texttt{facts\_first} and \texttt{expr\_first} files. The overall dataset still covers all 11 rules, but the hop split is not fully crossed in the released $20$k version: six rules appear in both depths (\textit{identity}, \textit{domination}, \textit{idempotent}, \textit{double\_negation}, \textit{excluded\_middle}, \textit{contradiction}), whereas the five multi-variable rules (\textit{commutative}, \textit{associative}, \textit{distributive}, \textit{demorgan}, \textit{absorption}) are allocated to one-hop only in the released unique corpora.

\paragraph{Row structure.}
Each JSONL row contains the identifiers and semantics needed for both behavioral evaluation and mechanistic analysis: \texttt{id}, \texttt{hop}, \texttt{rule}, \texttt{category}, \texttt{formal\_definition}, the clean fact assignment \texttt{facts}, the corrupted assignment \texttt{corrupted\_facts}, the clean and corrupted labels, the queried expression in symbolic and semi-natural form, compact symbolic prompts, verbose symbolic prompts, and a \texttt{meta} object describing the corruption and variable remapping. Two-hop rows additionally store \texttt{intermediate\_var}, \texttt{intermediate\_expr\_symbolic}, and \texttt{intermediate\_expr\_semi\_natural}, making the latent derived assignment explicit.

\paragraph{One-hop construction.}
For one-hop rows, we first instantiate a symbolic rule template, then randomly remap the template variables (\textit{e.g.}, $A,B,C$) to fresh uppercase symbols drawn from a reserved pool that excludes \texttt{F} and \texttt{T}. We next sample a Boolean assignment over the variables appearing in the instantiated expression and evaluate the resulting expression to obtain the clean label. Under \texttt{facts\_first}, the compact prompt takes the form
\texttt{X is True, Y is False, [expr] is. Answer with one word only: True or False.}
Under \texttt{expr\_first}, the same logical row is linearized as
\texttt{[expr] is? X is True, Y is False. Answer with one word only: True or False.}

\paragraph{Two-hop construction.}
For two-hop rows, we first sample a one-hop expression $e_{1}$ for a rule and remap its variables to fresh symbols. We then introduce two additional symbols: an intermediate variable $M$ and a gate variable $G$. The intermediate assignment is defined as $M = e_{1}$, and the final query is either $M \land G$ or $M \lor G$. The choice of conjunction versus disjunction is randomized unless the step-1 expression is a constant under all environments, in which case the query operator is chosen deterministically to avoid degeneracy. In the compact \texttt{facts\_first} template, a two-hop row therefore has the form
\texttt{[facts], M is [step-1 expr], [final expr] is. Answer with one word only: True or False.}
and in \texttt{expr\_first} the queried final expression is moved to the front while the factual assignments and derived assignment remain in the tail.

\paragraph{Corruption mechanism.}
The default corruption procedure performs an exact label-flipping search over the factual variables only. Given a clean row, it searches for a corrupted fact assignment whose label differs from the clean label, and stores the changed variables together with the Hamming distance inside \texttt{meta.corruption}. In the released \texttt{facts\_first} corpus, $19808/20000$ rows use this standard \texttt{label\_flipping\_fact\_search} mechanism. A small residual set of $192$ one-hop rows corresponds to expressions that are constant with respect to the available facts; for these rows, the generator falls back to a label-flipping expression substitution, recorded as \texttt{label\_flipping\_expr\_fallback}. The \texttt{expr\_first} corpus mirrors the same logical rows and therefore inherits the same corruption metadata pattern.

\paragraph{Filtering for mechanistic analysis.}
The raw dataset itself is not filtered by model behavior. Model-specific filtering happens only after inference: for mechanistic experiments we retain the dual-correct subset where a model answers both the clean and corrupted prompts correctly. This post-hoc filtering is what produces files such as \path{artifacts/filtered_dual_correct_Qwen3-8b.jsonl}, which are then used by the MLP, token, and head analyses described above.

% \begin{table}[t]
% \centering
% \small
% \setlength{\tabcolsep}{4pt}
% \begin{tabularx}{\columnwidth}{@{}lX@{}}
% \toprule
% \textbf{Depth} & \textbf{Prompt Template \& Reasoning Process} \\
% \midrule
% \multirow{3.5}{*}{\textbf{One-Hop}}
% & \texttt{"A is True, B is False, (A and B) is"} \newline
%   \color{gray}\textit{\footnotesize $\hookrightarrow$ Logic: $A \land B$} \\
% \addlinespace[0.4em]
% & \texttt{"A is True, ($\lnot$($\lnot$A)) is"} \newline
%   \color{gray}\textit{\footnotesize $\hookrightarrow$ Logic: $\lnot(\lnot A) \equiv A$} \\
% \addlinespace[0.4em]
% & \texttt{"A is True, B is False, ($\lnot$A or $\lnot$B) is"} \newline
%   \color{gray}\textit{\footnotesize $\hookrightarrow$ Logic: De Morgan ($\lnot A \lor \lnot B$)} \\
% \midrule
% \multirow{5}{*}{\textbf{Two-Hop}}
% & \texttt{"A is T, B is A and T, C is F, B and C is"} \newline
%   \color{gray}\textit{ \footnotesize $\hookrightarrow$ Chain: $(A \land \top) \to B; \quad (B \land C) \to \text{Ans}$} \\
% \addlinespace[0.4em]
% & \texttt{"A is T, B is F, C is $\lnot$(A and B), D is T, C or D is"} \newline
%   \color{gray}\textit{\footnotesize $\hookrightarrow$ Chain: $\text{NAND}(A,B) \to C; \quad (C \lor D) \to \text{Ans}$} \\
% \addlinespace[0.4em]
% & \texttt{"A is T, B is F, C is A and B, D is T, C or D is"} \newline
%   \color{gray}\textit{\footnotesize $\hookrightarrow$ Chain: $(A \land B) \to C; \quad (C \lor D) \to \text{Ans}$} \\
% \bottomrule
% \end{tabularx}
% \caption{\textbf{Probing prompts across reasoning depths.} We probe the model with queries requiring direct rule application (One-Hop) versus those necessitating intermediate latent variable derivation (Two-Hop).}
% \label{tab:reasoning_depth}
% \end{table}

\section{More Results}
\label{sec:appendix_results}

\subsection{Behavioral Performance on PropLogic-MI}
\label{sec:appendix_performance}
Before turning to the mechanistic analyses, we summarize the zero-shot behavioral performance of both models on the full \textit{PropLogic-MI} benchmark. Table~\ref{tab:appendix_performance_summary} reports the overall clean accuracy, corrupted accuracy, dual-correct rate, and hop-wise clean accuracy for all four evaluation settings. Table~\ref{tab:appendix_rule_performance} further breaks down clean accuracy by rule type.

\begin{table*}[t]
\centering
\small
\renewcommand{\arraystretch}{0.95}
\setlength{\tabcolsep}{6pt}
\begin{tabularx}{\textwidth}{@{} l l c c X c c @{} }
\toprule
\textbf{Model} & \textbf{Prompt Order} & \textbf{Clean} & \textbf{Corrupted} & \textbf{Dual-Correct Rate} & \textbf{One-Hop} & \textbf{Two-Hop} \\
\midrule
\textit{Qwen3-8B} & \texttt{facts\_first} & 86.98 & 88.17 & 75.44\% (15087/20000) & 86.35 & 87.55 \\
\textit{Qwen3-14B} & \texttt{facts\_first} & 87.42 & 89.11 & 76.55\% (15310/20000) & 92.64 & 82.68 \\
\textit{Qwen3-8B} & \texttt{expr\_first} & 74.22 & 79.25 & 53.89\% (10777/20000) & 79.09 & 69.79 \\
\textit{Qwen3-14B} & \texttt{expr\_first} & 74.39 & 77.91 & 52.80\% (10561/20000) & 88.34 & 61.72 \\
\bottomrule
\end{tabularx}
\caption{\textbf{Overall zero-shot performance on \textit{PropLogic-MI}.} All values except the dual-correct column are accuracies (\%). Dual-Correct Rate denotes the fraction of examples for which the model answers both the clean and corrupted prompts correctly; these examples form the candidate pool for the mechanistic analyses. One-Hop and Two-Hop denote clean-prompt accuracy at each reasoning depth.}
\label{tab:appendix_performance_summary}
\end{table*}

\begin{table}[t]
\centering
\small
\renewcommand{\arraystretch}{0.95}
\setlength{\tabcolsep}{5pt}
\begin{tabularx}{\columnwidth}{@{} X c c c c @{} }
\toprule
\textbf{Rule} & \textbf{8B FF} & \textbf{14B FF} & \textbf{8B EF} & \textbf{14B EF} \\
\midrule
Absorption & 84.8 & 95.0 & 84.4 & 85.4 \\
Associative & 92.1 & 99.0 & 61.8 & 76.0 \\
Commutative & 99.8 & 98.1 & 92.0 & 94.4 \\
Contradiction & 64.0 & 54.1 & 49.4 & 48.1 \\
De Morgan & 80.4 & 100.0 & 80.9 & 100.0 \\
Distributive & 72.8 & 71.5 & 76.8 & 88.8 \\
Domination & 95.9 & 88.6 & 62.5 & 50.3 \\
Double Negation & 84.8 & 87.1 & 72.7 & 70.3 \\
Excluded Middle & 98.7 & 86.1 & 72.9 & 51.4 \\
Idempotent & 92.1 & 85.6 & 83.5 & 76.5 \\
Identity & 91.4 & 96.7 & 79.4 & 77.2 \\
\bottomrule
\end{tabularx}
\caption{\textbf{Rule-wise clean accuracy (\%) on \textit{PropLogic-MI}.} FF = \texttt{facts\_first}; EF = \texttt{expr\_first}. Prompt order has a large impact on several rules, especially Associative, Domination, and Excluded Middle, while Commutative remains consistently easy and Contradiction remains consistently difficult.}
\label{tab:appendix_rule_performance}
\end{table}

Several trends are worth highlighting. First, prompt order has a large and consistent effect: moving from \texttt{facts\_first} to \texttt{expr\_first} lowers clean accuracy by 12.76 points in \textit{Qwen3-8B} and 13.03 points in \textit{Qwen3-14B}, while reducing the dual-correct rate by 21.55 and 23.75 points, respectively. Second, corrupted prompts are slightly easier in all four settings, with corrupted accuracy exceeding clean accuracy by 1.69--5.04 points. Third, hop depth interacts with prompt order: under \texttt{expr\_first}, both models show a substantial two-hop penalty, especially \textit{Qwen3-14B} (88.34\% one-hop vs. 61.72\% two-hop), whereas \texttt{facts\_first} remains considerably more stable. Finally, the rule-wise breakdown shows a persistent split between easy algebraic symmetries and harder conflict-style rules: Commutative is consistently above 92\%, De Morgan reaches 100\% in \textit{Qwen3-14B}, while Contradiction is the weakest rule in all four settings.

\subsection{Prompt-Order Intervention: \texttt{expr\_first}}
\label{sec:appendix_prompt_order}
We use \texttt{expr\_first} as a controlled prompt-order intervention. Relative to \texttt{facts\_first}, the logical graph, labels, variable assignments, and corruption semantics are unchanged; only the order of presentation is altered so that the queried expression appears before the factual assignments. This lets us ask whether the staged phenomena reported in the main text are primarily driven by the prompt layout itself or by a more persistent computational organization in the model. A purely prompt-oriented account would predict a large reordering of the stage profile when the queried expression is moved to the front. A purely model-oriented account would predict that the same coarse mechanisms persist with only minor quantitative shifts.

The observed pattern supports a mixed interpretation. Prompt order clearly matters behaviorally (Table~\ref{tab:appendix_performance_summary}), and it also redistributes where computation is concentrated. However, the coarse mechanisms identified in the main text do not disappear under \texttt{expr\_first}: staged MLP sensitivity remains visible, token-level convergence is still expected to center on the final query position, and the same three head roles remain recoverable with strong causal validation. We therefore interpret the \texttt{expr\_first} results as evidence that the existence of staged computation is largely model-driven, while the exact allocation of computation across regions is prompt-sensitive.

\paragraph{Staged computation under prompt reordering.}
The MLP results already show a systematic prompt-order shift rather than a collapse of the stage structure. In one-hop \textit{Qwen3-8B}, the early-band contribution ratio (BCR) of the Expression Region increases from $0.69$ under \texttt{facts\_first} to $0.77$ under \texttt{expr\_first}, while the Facts Region early-band BCR falls from $0.86$ to $0.69$. The same pattern appears in one-hop \textit{Qwen3-14B}, where the Expression Region early-band BCR rises from $0.68$ to $0.84$ and the Facts Region early-band BCR falls from $0.82$ to $0.60$. Meanwhile, the Query Region remains a late-stage aggregation site, with its late-band BCR increasing from $0.13$ to $0.19$ in one-hop \textit{Qwen3-8B}. This is consistent with a prompt-sensitive redistribution of computational load: placing the queried expression first pulls expression processing earlier, but it does not remove the late-stage answer integration step.

The same tendency persists in more difficult settings. In two-hop \textit{Qwen3-8B}, the late-band BCR of the Query Region rises from $0.17$ under \texttt{facts\_first} to $0.22$ under \texttt{expr\_first}, while the middle-band contribution remains dominant in both cases ($0.49$ vs. $0.43$). In \textit{Qwen3-14B}, one-hop \texttt{expr\_first} again shifts the Facts Region away from an overwhelmingly early-stage pattern (early-band BCR $0.82 \rightarrow 0.60$) while strengthening the expression-first computation profile (Expression Region early-band BCR $0.68 \rightarrow 0.84$). Thus, prompt reordering changes \emph{where} the model spends computation, but not the basic fact-processing $\rightarrow$ expression-processing $\rightarrow$ answer-integration template.

\begin{figure*}[t]
    \centering

    \begin{subfigure}[b]{0.9\textwidth}
        \centering
        \includegraphics[width=\linewidth]{./figures_experiments/reports/mlp_analysis_expr_first/comparison_expr_first/band_metrics_panel.png}
        \caption{\textit{Qwen3-8B}}
        \label{fig:expr_first_stage_8b}
    \end{subfigure}

    \vspace{0.8em}

    \begin{subfigure}[b]{0.9\textwidth}
        \centering
        \includegraphics[width=\linewidth]{./figures_experiments/reports/mlp_analysis_expr_first/comparison_expr_first_14B/band_metrics_panel.png}
        \caption{\textit{Qwen3-14B}}
        \label{fig:expr_first_stage_14b}
    \end{subfigure}

    \caption{\textbf{Staged computation under the \texttt{expr\_first} prompt-order intervention.} Band-level MLP ablation summaries for the released \texttt{expr\_first} runs, shown separately for \textit{Qwen3-8B} and \textit{Qwen3-14B}. Each panel compares one-hop and two-hop settings across the main semantic regions. The key pattern is a redistribution rather than a disappearance of stages: when the queried expression is placed first, the Expression Region acquires more early-band mass, the Facts Region becomes less exclusively early-layer dominated, and the Query Region retains a strong middle-to-late aggregation profile. This supports the claim that prompt order modulates the allocation of computation across regions, while the existence of a staged computational hierarchy remains stable.}
    \label{fig:expr_first_stage_panel}
\end{figure*}

\begin{figure*}[t]
    \centering

    \begin{subfigure}[b]{0.48\textwidth}
        \centering
        \includegraphics[width=\linewidth]{./figures_experiments/reports/mlp_analysis/mlp_8b_one_hop/query_region/mlp_query_region_scores.png}
        \caption{\textit{Qwen3-8B}, \texttt{facts\_first}}
        \label{fig:query_region_8b_ff}
    \end{subfigure}
    \hfill
    \begin{subfigure}[b]{0.48\textwidth}
        \centering
        \includegraphics[width=\linewidth]{./figures_experiments/reports/mlp_analysis_expr_first/mlp_8b_one_hop_expr_first/query_region/mlp_query_region_scores.png}
        \caption{\textit{Qwen3-8B}, \texttt{expr\_first}}
        \label{fig:query_region_8b_ef}
    \end{subfigure}

    \vspace{0.8em}

    \begin{subfigure}[b]{0.48\textwidth}
        \centering
        \includegraphics[width=\linewidth]{./figures_experiments/reports/mlp_analysis/mlp_14b_one_hop/query_region/mlp_query_region_scores.png}
        \caption{\textit{Qwen3-14B}, \texttt{facts\_first}}
        \label{fig:query_region_14b_ff}
    \end{subfigure}
    \hfill
    \begin{subfigure}[b]{0.48\textwidth}
        \centering
        \includegraphics[width=\linewidth]{./figures_experiments/reports/mlp_analysis_expr_first/mlp_14b_one_hop_expr_first/query_region/mlp_query_region_scores.png}
        \caption{\textit{Qwen3-14B}, \texttt{expr\_first}}
        \label{fig:query_region_14b_ef}
    \end{subfigure}

    \caption{\textbf{Query-region MLP sensitivity across prompt orders.} One-hop query-region ablation curves for \textit{Qwen3-8B} and \textit{Qwen3-14B} under both \texttt{facts\_first} and \texttt{expr\_first}. The Query Region subsumes the query anchor together with the answer-format suffix. It remains a stable middle-to-late integration site in all four settings, indicating that final answer formation depends on a broader suffix-aware region rather than on a single terminal token. Prompt reordering changes the amount of late-band mass, but does not eliminate this query-centric integration step.}
    \label{fig:query_region_prompt_order}
\end{figure*}

\paragraph{Specialized heads under prompt reordering.}
The head analysis shows the same pattern of prompt-sensitive reallocation atop a stable coarse mechanism. Under \texttt{expr\_first}, Splitting Heads become more numerous while Fact-Retrieval Heads become slightly less dominant: in \textit{Qwen3-8B}, the role totals shift from $66/129/175$ (Splitting / Transmission / Fact-Retrieval) under \texttt{facts\_first} to $116/114/167$ under \texttt{expr\_first}; in \textit{Qwen3-14B}, they shift from $87/169/217$ to $170/142/200$. Yet the causal validation remains strong. In \textit{Qwen3-8B}, the top-32 ablations under \texttt{expr\_first} still reduce accuracy by 29.6 points for Fact-Retrieval, 19.6 for Transmission, and 19.6 for Splitting, while the matched Random-Outside control changes accuracy by only 0.3 points. In \textit{Qwen3-14B}, the corresponding drops are 20.4, 21.0, and 10.8 points, while the Random-Outside control remains effectively at baseline. Thus, prompt reordering changes which heads are emphasized, but not the existence of a stable role taxonomy with strong causal relevance.

\begin{figure*}[t]
    \centering
    \begin{subfigure}[b]{0.48\textwidth}
        \centering
        \includegraphics[width=\linewidth]{./figures_experiments/reports/heads_analysis_qwen3_8b_expr_first/taxonomy/head_taxonomy_line_chart.png}
        \caption{\textit{Qwen3-8B}: role distribution}
        \label{fig:expr_first_heads_8b_taxonomy}
    \end{subfigure}
    \hfill
    \begin{subfigure}[b]{0.48\textwidth}
        \centering
        \includegraphics[width=\linewidth]{./figures_experiments/reports/heads_analysis_qwen3_14b_expr_first/taxonomy/head_taxonomy_line_chart.png}
        \caption{\textit{Qwen3-14B}: role distribution}
        \label{fig:expr_first_heads_14b_taxonomy}
    \end{subfigure}

    \vspace{0.8em}

    \begin{subfigure}[b]{0.48\textwidth}
        \centering
        \includegraphics[width=\linewidth]{./figures_experiments/reports/heads_analysis_qwen3_8b_expr_first/validation/plots/pd_abs_ratio_curve.png}
        \caption{\textit{Qwen3-8B}: validation curve}
        \label{fig:expr_first_heads_8b_curve}
    \end{subfigure}
    \hfill
    \begin{subfigure}[b]{0.48\textwidth}
        \centering
        \includegraphics[width=\linewidth]{./figures_experiments/reports/heads_analysis_qwen3_14b_expr_first/validation/plots/pd_abs_ratio_curve.png}
        \caption{\textit{Qwen3-14B}: validation curve}
        \label{fig:expr_first_heads_14b_curve}
    \end{subfigure}
    \caption{\textbf{Specialized head roles under the \texttt{expr\_first} prompt-order intervention.} Top row: layer-wise role distributions for \textit{Qwen3-8B} and \textit{Qwen3-14B} under \texttt{expr\_first}. Relative to \texttt{facts\_first}, Splitting Heads become more prevalent, consistent with the fact that the queried expression and the factual block now meet at a different structural boundary. Bottom row: Step 3 validation curves plotting $|\mathrm{dPD}|/|\mathrm{PD}_{\mathrm{clean}}|$ against the number of ablated heads. In both models, role-targeted ablations remain substantially more damaging than the Random-Outside control, confirming that prompt reordering changes the allocation of head roles without eliminating their causal importance. Full accuracy tables for the \texttt{expr\_first} runs are given in Tables~\ref{tab:qwen3_8b_ef_full_accuracy} and~\ref{tab:qwen3_14b_ef_full_accuracy}.}
    \label{fig:expr_first_heads_panel}
\end{figure*}

\subsection{Staged Computation}
\label{sec:appendix_stage}
We replicated the same MLP region mean-ablation analysis on \textit{Qwen3-14B} to verify cross-model consistency. Following the current implementation, we partition prompts into Facts, Expression, and Query regions, intervene on \texttt{mlp\_out} with within-sequence mean replacement, and summarize the resulting $\mathrm{dPD}$ traces using BMI/BCR/SBI. Figure~\ref{fig:staged_computation_14B} confirms that the same early-facts / middle-expression / late-query organization persists at the larger scale.

\subsection{Additional Region-Level Ablations}
\label{sec:appendix_region_variants}
The current codebase adds two region-level companions to the main MLP branch: \textbf{attention-region ablation} and \textbf{joint attention+MLP ablation}. Both reuse the same three-region parser and the same BMI/BCR/SBI summaries. To keep these branches visually comparable to the main-text MLP figure, Figures~\ref{fig:attn_region_triples} and~\ref{fig:attn_mlp_region_triples} group the three region curves for one model under the same one-hop \texttt{facts\_first} template in a single row.

Two patterns are especially consistent. First, attention-region ablation places relatively more BCR mass in the middle Expression Region and the late Query Region than pure MLP ablation. In one-hop \textit{Qwen3-8B}, the Expression Region middle-band BCR rises from $0.30$ in the MLP branch to $0.58$ in the attention branch, and the Query Region late-band BCR rises from $0.09$ to $0.12$. \textit{Qwen3-14B} shows the same shift ($0.32 \rightarrow 0.55$ for middle-band Expression BCR and $0.04 \rightarrow 0.09$ for late-band Query BCR). This is consistent with attention contributing more directly to cross-token routing and late integration.

Second, the joint attention+MLP intervention preserves the same coarse staged template while tightening early Facts Region dominance. In one-hop, the early-band Facts Region BCR becomes $0.86$ / $0.88$ in \textit{Qwen3-8B} / \textit{Qwen3-14B}, compared with $0.86$ / $0.82$ in the pure MLP branch. Thus, the joint branch does not overturn the staged picture; it reinforces the view that routing and nonlinear transformation contribute complementary pieces of the same three-stage computation.

\begin{figure*}[t]
    \centering
    \begin{subfigure}[b]{0.31\textwidth}
        \centering
        \includegraphics[width=\linewidth]{./figures_experiments/reports/attn_analysis/attn_8b_one_hop/facts_region/attn_facts_region_scores.png}
        \caption{\textit{Qwen3-8B}: Facts}
        \label{fig:attn_region_8b_facts}
    \end{subfigure}
    \hfill
    \begin{subfigure}[b]{0.31\textwidth}
        \centering
        \includegraphics[width=\linewidth]{./figures_experiments/reports/attn_analysis/attn_8b_one_hop/expression_region/attn_expression_region_scores.png}
        \caption{\textit{Qwen3-8B}: Expression}
        \label{fig:attn_region_8b_expr}
    \end{subfigure}
    \hfill
    \begin{subfigure}[b]{0.31\textwidth}
        \centering
        \includegraphics[width=\linewidth]{./figures_experiments/reports/attn_analysis/attn_8b_one_hop/query_region/attn_query_region_scores.png}
        \caption{\textit{Qwen3-8B}: Query}
        \label{fig:attn_region_8b_query}
    \end{subfigure}

    \vspace{0.8em}

    \begin{subfigure}[b]{0.31\textwidth}
        \centering
        \includegraphics[width=\linewidth]{./figures_experiments/reports/attn_analysis/attn_14b_one_hop/facts_region/attn_facts_region_scores.png}
        \caption{\textit{Qwen3-14B}: Facts}
        \label{fig:attn_region_14b_facts}
    \end{subfigure}
    \hfill
    \begin{subfigure}[b]{0.31\textwidth}
        \centering
        \includegraphics[width=\linewidth]{./figures_experiments/reports/attn_analysis/attn_14b_one_hop/expression_region/attn_expression_region_scores.png}
        \caption{\textit{Qwen3-14B}: Expression}
        \label{fig:attn_region_14b_expr}
    \end{subfigure}
    \hfill
    \begin{subfigure}[b]{0.31\textwidth}
        \centering
        \includegraphics[width=\linewidth]{./figures_experiments/reports/attn_analysis/attn_14b_one_hop/query_region/attn_query_region_scores.png}
        \caption{\textit{Qwen3-14B}: Query}
        \label{fig:attn_region_14b_query}
    \end{subfigure}
    \caption{\textbf{Attention-region mean-ablation results in the main-text layout.} Each row groups the three region curves for one model under the same one-hop \texttt{facts\_first} template, matching the presentation style of Figure~\ref{fig:stage}. Compared with the MLP branch, the attention branch puts relatively more weight on middle-layer expression processing and late query integration.}
    \label{fig:attn_region_triples}
\end{figure*}

\begin{figure*}[t]
    \centering
    \begin{subfigure}[b]{0.31\textwidth}
        \centering
        \includegraphics[width=\linewidth]{./figures_experiments/reports/attn_mlp_analysis/attn_mlp_8b_one_hop/facts_region/attn_mlp_facts_region_scores.png}
        \caption{\textit{Qwen3-8B}: Facts}
        \label{fig:attn_mlp_region_8b_facts}
    \end{subfigure}
    \hfill
    \begin{subfigure}[b]{0.31\textwidth}
        \centering
        \includegraphics[width=\linewidth]{./figures_experiments/reports/attn_mlp_analysis/attn_mlp_8b_one_hop/expression_region/attn_mlp_expression_region_scores.png}
        \caption{\textit{Qwen3-8B}: Expression}
        \label{fig:attn_mlp_region_8b_expr}
    \end{subfigure}
    \hfill
    \begin{subfigure}[b]{0.31\textwidth}
        \centering
        \includegraphics[width=\linewidth]{./figures_experiments/reports/attn_mlp_analysis/attn_mlp_8b_one_hop/query_region/attn_mlp_query_region_scores.png}
        \caption{\textit{Qwen3-8B}: Query}
        \label{fig:attn_mlp_region_8b_query}
    \end{subfigure}

    \vspace{0.8em}

    \begin{subfigure}[b]{0.31\textwidth}
        \centering
        \includegraphics[width=\linewidth]{./figures_experiments/reports/attn_mlp_analysis/attn_mlp_14b_one_hop/facts_region/attn_mlp_facts_region_scores.png}
        \caption{\textit{Qwen3-14B}: Facts}
        \label{fig:attn_mlp_region_14b_facts}
    \end{subfigure}
    \hfill
    \begin{subfigure}[b]{0.31\textwidth}
        \centering
        \includegraphics[width=\linewidth]{./figures_experiments/reports/attn_mlp_analysis/attn_mlp_14b_one_hop/expression_region/attn_mlp_expression_region_scores.png}
        \caption{\textit{Qwen3-14B}: Expression}
        \label{fig:attn_mlp_region_14b_expr}
    \end{subfigure}
    \hfill
    \begin{subfigure}[b]{0.31\textwidth}
        \centering
        \includegraphics[width=\linewidth]{./figures_experiments/reports/attn_mlp_analysis/attn_mlp_14b_one_hop/query_region/attn_mlp_query_region_scores.png}
        \caption{\textit{Qwen3-14B}: Query}
        \label{fig:attn_mlp_region_14b_query}
    \end{subfigure}
    \caption{\textbf{Joint attention+MLP mean-ablation results in the main-text layout.} Each row again groups one model and one prompt template across the three semantic regions. The resulting curves preserve the same coarse staged ordering as the MLP branch while sharpening early Facts dominance, consistent with attention routing and MLP transformation contributing complementary parts of the same mechanism.}
    \label{fig:attn_mlp_region_triples}
\end{figure*}

\subsection{Information Transmission}
\label{sec:appendix_info_converge}

\paragraph{Metric Definitions.}
The current token pipeline measures causal effects using the label-aligned True/False probability difference rather than the raw logit gap. Concretely, for the final answer position we define
$$\mathrm{PD}(x) = p_{\theta}(\texttt{True}\mid x) - p_{\theta}(\texttt{False}\mid x),$$
where the probabilities are computed by a two-way softmax restricted to the True/False answer tokens. For a patched run at layer $l$ and token position $i$, the token-level shift is
$$\mathrm{dPD}_{l,i} = \mathrm{PD}_{\mathrm{patched}}^{(l,i)} - \mathrm{PD}_{\mathrm{baseline}}.$$
Throughout this appendix, we therefore describe the token-level intervention effect using the updated \textbf{probability-difference shift} notation.

We employ two aggregation strategies:

\textbf{(1) Mean $|\mathrm{dPD}|$ across samples} (Figure~\ref{fig:info_converge_token} and Figure~\ref{fig:info_converge_appendix}): We aggregate $|\mathrm{dPD}|$ by: (i) averaging over layers within each stage (Early, Middle, Late), (ii) averaging over tokens belonging to the same category within each sample, and (iii) averaging across all samples in the dataset.

\textbf{(2) Mean $|\mathrm{dPD}|$ per token} ($\mathrm{M}|\mathrm{dPD}|$, Figure~\ref{fig:infor_convergence_multi_corrupted}): For single-sample case studies, we compute the mean absolute shift at each token position $i$ within a layer stage $S$:
$$\mathrm{M}|\mathrm{dPD}|_i^{(S)} = \frac{1}{|S|}\sum_{l \in S} |\mathrm{dPD}_{l,i}|.$$
This metric captures how much patching each token position shifts the model's answer preference within each layer group.

\paragraph{Single-Sample Case Studies.}
Figure~\ref{fig:infor_convergence_multi_corrupted} extends the single-sample analysis from Figure~\ref{fig:logit_difference} (main text) by systematically varying which fact tokens are corrupted. This provides a controlled validation of the Information Transmission phenomenon across different corruption configurations.

\begin{figure}[t]
  \includegraphics[width=\columnwidth]{./figures/info-convergence-multi-corrupted.png}
	\caption{Token-wise activation patching under multiple corruption scenarios.
	Each panel shows $\mathrm{M}|\mathrm{dPD}|$ per token across Early (L0-13), Middle (L14-23), and Late (L24-35) layer groups. The clean prompt is \texttt{A is True, B is False, ($\neg$A or $\neg$B) is}.
	Top: B changes (False$\to$True). The corrupted \texttt{True} token (position 8) exhibits high sensitivity in early/middle layers.
	Middle: A changes (True$\to$False). The corrupted \texttt{False} token shows similar early-layer dominance.
	Bottom: Both A and B change. Combined effects appear at both fact positions.
	Across all scenarios, the query token \texttt{is} shows a significant surge in late layers, confirming information convergence toward the prediction position.}
	\label{fig:infor_convergence_multi_corrupted}
\end{figure}

The results reveal several consistent patterns across all corruption scenarios:

(1) Early-layer fact encoding. In each scenario, the corrupted fact token(s) exhibit maximal patching sensitivity in early layers (L0-13). When only B changes (Top panel), position 8 dominates; when only A changes (Middle panel), position 4 dominates; when both change (Bottom panel), both positions show elevated effects. This confirms that fact tokens are primarily encoded in early layers regardless of which specific facts are manipulated.

(2) Late-layer query convergence. Across all three scenarios, the query token \texttt{is} consistently shows a pronounced surge in late layers (L24-35), with $\mathrm{M}|\mathrm{dPD}|$ values of +2.0, +0.8, and +2.5 respectively. This invariance across corruption types provides strong evidence that late-layer query convergence is a robust architectural property, not an artifact of specific input configurations.

(3) Additive effects in multi-fact corruption. The Bottom panel (both A and B corrupted) shows that patching effects at individual fact positions are approximately additive: the combined effect at positions 4 and 8 roughly equals the sum of effects observed when each is corrupted individually. This suggests that the model processes multiple facts through parallel, largely independent early-layer pathways before integrating them at later stages.

\paragraph{Extended Results.}
To verify the generalizability of the Information Transmission phenomenon, we extend our token-wise $|\mathrm{dPD}|$ analysis to \textit{Qwen3-14B} and two-hop reasoning tasks. Figure~\ref{fig:info_converge_appendix} presents the complete results.

\begin{figure*}[t]
    \centering
    \begin{subfigure}[b]{0.32\textwidth}
        \centering
        \includegraphics[width=\linewidth]{./figures/info-converge-twohop-8b.png}
        \caption{Two-hop (\textit{Qwen3-8B})}
        \label{fig:info_converge_twohop_8b}
    \end{subfigure}
    \hfill
    \begin{subfigure}[b]{0.32\textwidth}
        \centering
        \includegraphics[width=\linewidth]{./figures/info-converge-twohop-14b.png}
        \caption{Two-hop (\textit{Qwen3-14B})}
        \label{fig:info_converge_twohop_14b}
    \end{subfigure}
    \hfill
    \begin{subfigure}[b]{0.32\textwidth}
        \centering
        \includegraphics[width=\linewidth]{./figures/info-converge-onehop-14b.png}
        \caption{One-hop (\textit{Qwen3-14B})}
        \label{fig:info_converge_onehop_14b}
    \end{subfigure}

    \caption{Token-wise information convergence across model scales and reasoning depths.
    All panels show Mean $|\mathrm{dPD}|$ per token category across Early (L0-13), Middle (L14-23), and Late (L24-35) layer groups.
    (a-b) Two-hop reasoning: The convergence pattern persists with attenuated magnitudes due to increased reasoning complexity. Notably, \texttt{query\_token} remains dominant in late layers, while \texttt{expr\_last} (the segment-terminal token of the expression region) also shows elevated $|\mathrm{dPD}|$ in middle layers, confirming that segment-terminal tokens serve as information hubs.
    (c) Qwen3-14B One-hop: Consistent with the 8B model (Figure~\ref{fig:info_converge_token}), demonstrating scale-invariant information flow patterns.
    These results confirm that Information Transmission is a robust mechanism across both model scales and reasoning depths.}
    \label{fig:info_converge_appendix}
\end{figure*}


\paragraph{Cross-Model Consistency.}
Comparing Figure~\ref{fig:info_converge_token} (\textit{Qwen3-8B}) with Figure~\ref{fig:info_converge_onehop_14b} (\textit{Qwen3-14B}), both models exhibit nearly identical convergence patterns on one-hop tasks. The \texttt{facts\_value} dominance in early layers and \texttt{query\_token} surge in late layers are preserved across the 8B$\rightarrow$14B scaling, suggesting that Information Transmission is an architectural property rather than a scale-dependent phenomenon.

\paragraph{Reasoning Depth Effects.}
In two-hop reasoning (Figures~\ref{fig:info_converge_twohop_8b} and~\ref{fig:info_converge_twohop_14b}), the overall $|\mathrm{dPD}|$ magnitudes are substantially lower compared to one-hop tasks. This attenuation reflects the increased computational complexity: with additional intermediate reasoning steps, causal effects become more distributed across the network. Nevertheless, the qualitative convergence pattern, characterized by early factual encoding followed by late-layer query aggregation, remains intact. Furthermore, the \texttt{expr\_last} token (representing the final token of the expression region) exhibits moderate but consistent $|\mathrm{dPD}|$ values in middle layers across all settings, reinforcing our finding that segment-terminal tokens serve as critical information hubs throughout the network.

\subsection{Specialized Attention Heads: Robustness and Validation}
\label{appendix_heads}
To test whether Specialized Attention Heads generalize beyond the one-hop \textit{Qwen3-8B} example in the main text, we reran the same taxonomy on \textit{Qwen3-8B} and \textit{Qwen3-14B} under both \texttt{facts\_first} and \texttt{expr\_first} prompt orders, using all filtered one-hop and two-hop samples. Figure~\ref{fig:counts_14B_comparison} shows the \textit{Qwen3-14B} distributions; the corresponding \textit{Qwen3-8B} trends follow the same pattern.

The coarse organization is stable across all settings. Splitting Heads remain concentrated in earlier layers, Transmission Heads span the middle of the model, and Fact-Retrieval Heads peak in middle-to-late layers and persist deep into the network. Prompt order changes the relative abundance of roles but not the overall topology: \texttt{expr\_first} prompts produce more Splitting Heads, whereas \texttt{facts\_first} prompts produce relatively more Fact-Retrieval Heads, consistent with the changed placement of structural boundaries and fact tokens. This robustness indicates that the head taxonomy tracks the computation required by the prompt rather than artifacts of a single format.

Step 3 further provides a direct causal test. Table~\ref{tab:qwen3_head_accuracy_appendix} reports held-out clean-prompt accuracy after ablating the top-32 heads for each role, using balanced validation samples drawn from the filtered dual-correct pool. On \textit{Qwen3-8B} with \texttt{facts\_first} prompts, ablating Fact-Retrieval / Transmission / Splitting heads lowers accuracy from 96.6\% to 57.0\% / 69.2\% / 80.4\%, while the matched Random-Outside control remains at 97.2\%. The same qualitative gap appears in the other three settings: role-specific ablations reduce accuracy by 29.6 / 19.6 / 19.6 points in \textit{Qwen3-8B} with \texttt{expr\_first}, by 31.0 / 25.6 / 10.0 points in \textit{Qwen3-14B} with \texttt{facts\_first}, and by 20.4 / 21.0 / 10.8 points in \textit{Qwen3-14B} with \texttt{expr\_first}. By contrast, the Random-Outside baseline remains near the original accuracy in every setting. Tables~\ref{tab:qwen3_8b_ff_full_accuracy}--\ref{tab:qwen3_14b_ef_full_accuracy} provide the complete Qwen3 accuracy tables for all strategy types and all $k \in \{1,2,4,8,16,32,64\}$. Together, these results confirm that the identified roles are not merely visually distinctive attention patterns: they are causally important components of the mechanism uncovered in the main analysis.

\begin{table*}[t]
\centering
\small
\renewcommand{\arraystretch}{0.96}
\setlength{\tabcolsep}{5pt}
\begin{tabular}{@{} l c c c c @{}}
\toprule
\textbf{Role ($k=32$)} & \textbf{Qwen3-8B FF} & \textbf{Qwen3-8B EF} & \textbf{Qwen3-14B FF} & \textbf{Qwen3-14B EF} \\
\midrule
Base accuracy & 96.6 & 92.8 & 94.8 & 97.6 \\
Fact-Retrieval & 57.0 & 63.2 & 63.8 & 77.2 \\
Transmission & 69.2 & 73.2 & 69.2 & 76.6 \\
Splitting & 80.4 & 73.2 & 84.8 & 86.8 \\
Fact-Retrieval + Splitting & 57.6 & 64.6 & 66.2 & 75.6 \\
Fact-Retrieval + Transmission & 57.4 & 61.2 & 60.6 & 67.6 \\
Splitting + Transmission & 69.4 & 71.4 & 68.8 & 71.6 \\
Mixed-All & 62.0 & 67.4 & 61.4 & 62.0 \\
Random-Outside & 97.2 $\pm$ 0.4 & 92.5 $\pm$ 0.8 & 98.3 $\pm$ 0.2 & 98.2 $\pm$ 0.2 \\
\bottomrule
\end{tabular}
\caption{\textbf{Qwen3 accuracy table for Step 3 head validation.} Entries report post-ablation accuracy (\%) on held-out clean prompts from balanced validation samples drawn from the filtered dual-correct pool. FF = \texttt{facts\_first}; EF = \texttt{expr\_first}. Compared with the near-baseline \textit{Random-Outside} control, ablating heads selected by the discovered roles consistently produces large accuracy drops, supporting their causal relevance.}
\label{tab:qwen3_head_accuracy_appendix}
\end{table*}






\begin{figure*}[t]
    \centering
    \begin{subfigure}[b]{0.48\textwidth}
        \centering
        \includegraphics[width=\linewidth]{./figures_experiments/reports/heads_analysis_qwen3_14b_facts_first/taxonomy/head_taxonomy_line_chart.png}
        \caption{\textbf{\texttt{facts\_first}}}
        \label{fig:counts_14B_onehop}
    \end{subfigure}
    \hfill
    \begin{subfigure}[b]{0.48\textwidth}
        \centering
        \includegraphics[width=\linewidth]{./figures_experiments/reports/heads_analysis_qwen3_14b_expr_first/taxonomy/head_taxonomy_line_chart.png}
        \caption{\textbf{\texttt{expr\_first}}}
        \label{fig:counts_14B_twohop}
    \end{subfigure}
    
\caption{Layer-wise distribution of functional attention heads in \textit{Qwen3-14B}, rendered from the saved taxonomy plots in \texttt{figures\_experiments/reports/heads\_analysis\_qwen3\_14b\_facts\_first/taxonomy/head\_taxonomy\_line\_chart.png} and \texttt{figures\_experiments/reports/heads\_analysis\_qwen3\_14b\_expr\_first/taxonomy/head\_taxonomy\_line\_chart.png}. Panel (a) shows the default \texttt{facts\_first} setting, while panel (b) shows the prompt-order intervention \texttt{expr\_first}. In both cases, Splitting Heads concentrate in early layers, Transmission Heads span the middle of the network, and Fact-Retrieval Heads peak in middle-to-late layers. The main prompt-order effect is quantitative rather than qualitative: \texttt{expr\_first} increases the abundance of Splitting Heads while preserving the same coarse layer-wise topology.}
    \label{fig:counts_14B_comparison}
\end{figure*}

\subsection{Cross-Model Extensions: Llama-3.1-8B-Instruct and Mistral-7B-Instruct-v0.1}
\label{sec:appendix_llama}
We also ran the same pipeline on \textit{Llama-3.1-8B-Instruct} and \textit{Mistral-7B-Instruct-v0.1}. The committed Llama artifact snapshot covers behavioral evaluation, token analysis, MLP region ablation, attention-region ablation, joint attention+MLP region ablation, and head analysis under both \texttt{facts\_first} and \texttt{expr\_first}. The committed Mistral snapshot covers behavioral evaluation, token analysis, MLP region ablation, attention-region ablation, and head analysis under both prompt orders; the same code path also supports joint attention+MLP region ablation, but those outputs are not part of the current artifact snapshot, so we do not make numerical claims about that branch here.

\paragraph{Llama: behavioral comparison with Qwen3.}
Behaviorally, Llama underperforms the Qwen3 family on the same benchmark. Under \texttt{facts\_first}, Llama reaches $63.86\%$ clean accuracy and a $38.37\%$ dual-correct rate, compared with $86.98\%$ / $75.44\%$ for \textit{Qwen3-8B}. Under \texttt{expr\_first}, the gap remains large: Llama reaches $60.34\%$ clean accuracy and $37.02\%$ dual-correct, whereas \textit{Qwen3-8B} reaches $74.22\%$ and $53.89\%$. The hop interaction is also qualitatively different. Llama's \texttt{expr\_first} prompts are better on one-hop than \texttt{facts\_first} ($81.47\%$ vs. $69.05\%$), but substantially worse on two-hop ($41.13\%$ vs. $59.13\%$). By contrast, Qwen3 remains much more stable under the same intervention. This suggests that Llama is more brittle to prompt-order changes once multi-step composition is required.

\paragraph{Llama: token and region-level structure.}
The token-analysis branch is now complete. Figure~\ref{fig:token_facts_first_cross_model} collects the current \texttt{facts\_first} refined token-analysis outputs for all four models under a shared layout, with one subfigure per model and prompt template. In the refined \texttt{facts\_first} token analysis, the dominant Llama categories are still \texttt{query\_token} and \texttt{facts\_value}, but their total contribution is much weaker than in Qwen3-8B: $1.63$ and $0.72$ in Llama, versus $2.22$ and $2.47$ in Qwen3-8B. This suggests that the same information-transmission motif is present, but less sharply concentrated.

\begin{figure*}[t]
    \centering
    \begin{subfigure}[b]{0.48\textwidth}
        \centering
        \includegraphics[width=\linewidth]{./figures_experiments/reports/token_analysis/8b_facts_first_all_hop_refined/refined_by_stage_sum.png}
        \caption{\textit{Qwen3-8B}, \texttt{facts\_first}}
        \label{fig:token_ff_q8}
    \end{subfigure}
    \hfill
    \begin{subfigure}[b]{0.48\textwidth}
        \centering
        \includegraphics[width=\linewidth]{./figures_experiments/reports/token_analysis/14b_facts_first_all_hop_refined/refined_by_stage_sum.png}
        \caption{\textit{Qwen3-14B}, \texttt{facts\_first}}
        \label{fig:token_ff_q14}
    \end{subfigure}

    \vspace{0.8em}

    \begin{subfigure}[b]{0.48\textwidth}
        \centering
        \includegraphics[width=\linewidth]{./figures_experiments/reports_llama3/token_analysis_facts_first_refined/refined_by_stage_sum.png}
        \caption{\textit{Llama-3.1-8B-Instruct}, \texttt{facts\_first}}
        \label{fig:token_ff_llama}
    \end{subfigure}
    \hfill
    \begin{subfigure}[b]{0.48\textwidth}
        \centering
        \includegraphics[width=\linewidth]{./figures_experiments/reports_mistral/token_analysis_facts_first_refined/refined_by_stage_sum.png}
        \caption{\textit{Mistral-7B-Instruct-v0.1}, \texttt{facts\_first}}
        \label{fig:token_ff_mistral}
    \end{subfigure}
    \caption{\textbf{Cross-model refined token analysis under \texttt{facts\_first}.} Each subfigure corresponds to one model under the same prompt template, following the same grouping rule as the region-patching figures. All panels are rendered from the saved \texttt{refined\_by\_stage\_sum.png} outputs. Across all four models, \texttt{facts\_value} and \texttt{query\_token} remain the dominant categories, but their relative balance differs: Qwen3 concentrates more total mass on fact tokens, whereas Llama and Mistral show comparatively stronger query-region dominance.}
    \label{fig:token_facts_first_cross_model}
\end{figure*}

Despite weaker behavioral performance, Llama still shows a coarse staged profile. Figure~\ref{fig:llama_mlp_panel} renders the saved comparison plots from \path{figures_experiments/reports_llama3/mlp_analysis_facts_first/comparison/band_metrics_panel.png} and \path{figures_experiments/reports_llama3/mlp_analysis_expr_first/comparison/band_metrics_panel.png}. Facts and expressions remain early-to-middle dominant, while the Query Region accumulates a comparatively strong late contribution. Relative to Qwen3, the Query Region is much heavier throughout depth: for one-hop \texttt{facts\_first}, its late-band BCR is $0.22$ in Llama, versus $0.09$ in \textit{Qwen3-8B}. A plausible hypothesis is that Llama relies more heavily on answer-position scaffolding and suffix-aware integration, whereas Qwen3 exhibits a cleaner separation between semantic content regions and final answer integration.

\begin{figure*}[t]
    \centering

    \begin{subfigure}[b]{0.9\textwidth}
        \centering
        \includegraphics[width=\linewidth]{./figures_experiments/reports_llama3/mlp_analysis_facts_first/comparison/band_metrics_panel.png}
        \caption{\texttt{facts\_first}}
        \label{fig:llama_mlp_facts_first}
    \end{subfigure}

    \vspace{0.8em}

    \begin{subfigure}[b]{0.9\textwidth}
        \centering
        \includegraphics[width=\linewidth]{./figures_experiments/reports_llama3/mlp_analysis_expr_first/comparison/band_metrics_panel.png}
        \caption{\texttt{expr\_first}}
        \label{fig:llama_mlp_expr_first}
    \end{subfigure}

    \caption{\textbf{Llama-3.1-8B-Instruct: staged computation under two prompt orders.} Band-level MLP ablation summaries rendered from the completed Llama runs. Panel (a) shows the default \texttt{facts\_first} condition; panel (b) shows the prompt-order intervention \texttt{expr\_first}. The same coarse staged pattern remains visible in both cases, but compared with Qwen3 the Query Region carries more weight throughout depth, suggesting a more output-position-centered computation strategy.}
    \label{fig:llama_mlp_panel}
\end{figure*}


\begin{figure*}[t]
    \centering

    \begin{subfigure}[b]{0.48\textwidth}
        \centering
        \includegraphics[width=\linewidth]{./figures_experiments/reports_llama3/mlp_analysis_facts_first/comparison/band_metrics_panel.png}
        \caption{\texttt{facts\_first}: MLP}
        \label{fig:llama_band_ff_mlp}
    \end{subfigure}
    \hfill
    \begin{subfigure}[b]{0.48\textwidth}
        \centering
        \includegraphics[width=\linewidth]{./figures_experiments/reports_llama3/attn_analysis_facts_first/comparison/band_metrics_panel.png}
        \caption{\texttt{facts\_first}: Attention}
        \label{fig:llama_band_ff_attn}
    \end{subfigure}

    \vspace{0.8em}

    \begin{subfigure}[b]{0.48\textwidth}
        \centering
        \includegraphics[width=\linewidth]{./figures_experiments/reports_llama3/attn_mlp_analysis_facts_first/comparison/band_metrics_panel.png}
        \caption{\texttt{facts\_first}: Attention+MLP}
        \label{fig:llama_band_ff_attn_mlp}
    \end{subfigure}
    \hfill
    \begin{subfigure}[b]{0.48\textwidth}
        \centering
        \includegraphics[width=\linewidth]{./figures_experiments/reports_llama3/mlp_analysis_expr_first/comparison/band_metrics_panel.png}
        \caption{\texttt{expr\_first}: MLP}
        \label{fig:llama_band_ef_mlp}
    \end{subfigure}

    \vspace{0.8em}

    \begin{subfigure}[b]{0.48\textwidth}
        \centering
        \includegraphics[width=\linewidth]{./figures_experiments/reports_llama3/attn_analysis_expr_first/comparison/band_metrics_panel.png}
        \caption{\texttt{expr\_first}: Attention}
        \label{fig:llama_band_ef_attn}
    \end{subfigure}
    \hfill
    \begin{subfigure}[b]{0.48\textwidth}
        \centering
        \includegraphics[width=\linewidth]{./figures_experiments/reports_llama3/attn_mlp_analysis_expr_first/comparison/band_metrics_panel.png}
        \caption{\texttt{expr\_first}: Attention+MLP}
        \label{fig:llama_band_ef_attn_mlp}
    \end{subfigure}

    \caption{\textbf{Llama-3.1-8B-Instruct: band-level overviews for all three region-patching analyses.} Columns correspond to MLP, attention, and joint attention+MLP ablation; rows correspond to \texttt{facts\_first} and \texttt{expr\_first}. This collects the current Llama \texttt{band\_metrics\_panel} outputs under the same organizational rule used for Qwen3. Across branches, the Query Region remains unusually heavy relative to Qwen3, and prompt reordering continues to redistribute mass without removing the coarse staged pattern.}
    \label{fig:llama_band_overview}
\end{figure*}


\paragraph{Llama: head roles and causal validation.}
The same three head roles remain identifiable in Llama, but they are less cleanly separated than in Qwen3. Under \texttt{facts\_first}, the Llama taxonomy yields $323$ classified heads in total (Fact-Retrieval $119$, Splitting $106$, Transmission $98$); under \texttt{expr\_first}, it yields $377$ heads (Fact-Retrieval $139$, Splitting $146$, Transmission $92$). Thus, the prompt-order effect on role allocation still matches Qwen3 qualitatively: \texttt{expr\_first} again increases the weight of Splitting Heads. However, the Step 3 validation curves and accuracy tables show a different causal signature. Under \texttt{facts\_first}, the base held-out accuracy is only $60.2\%$, and all role-targeted $k=32$ ablations collapse to $55.2\%$. Under \texttt{expr\_first}, the base held-out accuracy rises to $69.2\%$, but the same saturation pattern appears, with the corresponding $k=32$ role ablations all collapsing to $58.4\%$. Moreover, the Random-Outside control improves accuracy to $72.0\%$ / $79.2\%$ for \texttt{facts\_first} / \texttt{expr\_first}, indicating that mean-ablation replacement may regularize some diffuse heads rather than simply removing irrelevant structure. This is a notable contrast with Qwen3, where Random-Outside stays near baseline and role-targeted interventions produce much cleaner separations. The complete Llama Step 3 accuracy tables are reported in Tables~\ref{tab:llama_ff_full_accuracy} and~\ref{tab:llama_ef_full_accuracy}.

\begin{figure*}[t]
    \centering
    \begin{subfigure}[b]{0.48\textwidth}
        \centering
        \includegraphics[width=\linewidth]{./figures_experiments/reports_llama3/heads_analysis_facts_first/taxonomy/head_taxonomy_line_chart.png}
        \caption{\texttt{facts\_first}: role distribution}
        \label{fig:llama_heads_facts_tax}
    \end{subfigure}
    \hfill
    \begin{subfigure}[b]{0.48\textwidth}
        \centering
        \includegraphics[width=\linewidth]{./figures_experiments/reports_llama3/heads_analysis_expr_first/taxonomy/head_taxonomy_line_chart.png}
        \caption{\texttt{expr\_first}: role distribution}
        \label{fig:llama_heads_expr_tax}
    \end{subfigure}

    \vspace{0.8em}

    \begin{subfigure}[b]{0.48\textwidth}
        \centering
        \includegraphics[width=\linewidth]{./figures_experiments/reports_llama3/heads_analysis_facts_first/validation/plots/pd_abs_ratio_curve.png}
        \caption{\texttt{facts\_first}: validation curve}
        \label{fig:llama_heads_facts_val}
    \end{subfigure}
    \hfill
    \begin{subfigure}[b]{0.48\textwidth}
        \centering
        \includegraphics[width=\linewidth]{./figures_experiments/reports_llama3/heads_analysis_expr_first/validation/plots/pd_abs_ratio_curve.png}
        \caption{\texttt{expr\_first}: validation curve}
        \label{fig:llama_heads_expr_val}
    \end{subfigure}
    \caption{\textbf{Llama-3.1-8B-Instruct: specialized head roles under two prompt orders.} Top row: layer-wise role distributions for the completed Llama head-analysis runs. The same three roles appear as in Qwen3, but their causal separation is weaker. Bottom row: Step 3 validation curves plotting $|\mathrm{dPD}|/|\mathrm{PD}_{\mathrm{clean}}|$ against the number of ablated heads. The curves show that role-targeted interventions remain harmful, but the separation from Random-Outside is much less clean than in Qwen3, consistent with a more diffuse or less robust head-level specialization.}
    \label{fig:llama_heads_panel}
\end{figure*}

\paragraph{Mistral: behavioral, token, and region-level structure.}
Mistral presents a third behavioral profile. Under \texttt{facts\_first}, it reaches $65.09\%$ clean accuracy and a $35.11\%$ dual-correct rate. Under \texttt{expr\_first}, these rise slightly to $66.73\%$ and $39.98\%$. The gain is driven by one-hop prompts ($62.88\% \rightarrow 75.69\%$), while two-hop performance drops sharply ($67.09\% \rightarrow 58.59\%$), indicating that front-loading the query helps short-range reasoning but hurts deeper composition.

The refined Mistral token analysis is also complete. As in Qwen3 and Llama, \texttt{query\_token} and \texttt{facts\_value} remain the two dominant categories, with total contributions $1.04$ and $0.75$ under \texttt{facts\_first}. Thus, the same qualitative information-transmission signature remains visible, but again with lower total mass than in Qwen3. Figure~\ref{fig:token_expr_first_cross_model} further collects the current \texttt{expr\_first} refined token-analysis outputs for all four models under the same grouped layout, making the prompt-order intervention directly comparable across model families.

\begin{figure*}[t]
    \centering
    \begin{subfigure}[b]{0.48\textwidth}
        \centering
        \includegraphics[width=\linewidth]{./figures_experiments/reports/token_analysis/8b_expr_first_all_hop_refined/refined_by_stage_sum.png}
        \caption{\textit{Qwen3-8B}, \texttt{expr\_first}}
        \label{fig:token_ef_q8}
    \end{subfigure}
    \hfill
    \begin{subfigure}[b]{0.48\textwidth}
        \centering
        \includegraphics[width=\linewidth]{./figures_experiments/reports/token_analysis/14b_expr_first_all_hop_refined/refined_by_stage_sum.png}
        \caption{\textit{Qwen3-14B}, \texttt{expr\_first}}
        \label{fig:token_ef_q14}
    \end{subfigure}

    \vspace{0.8em}

    \begin{subfigure}[b]{0.48\textwidth}
        \centering
        \includegraphics[width=\linewidth]{./figures_experiments/reports_llama3/token_analysis_expr_first_refined/refined_by_stage_sum.png}
        \caption{\textit{Llama-3.1-8B-Instruct}, \texttt{expr\_first}}
        \label{fig:token_ef_llama}
    \end{subfigure}
    \hfill
    \begin{subfigure}[b]{0.48\textwidth}
        \centering
        \includegraphics[width=\linewidth]{./figures_experiments/reports_mistral/token_analysis_expr_first_refined/refined_by_stage_sum.png}
        \caption{\textit{Mistral-7B-Instruct-v0.1}, \texttt{expr\_first}}
        \label{fig:token_ef_mistral}
    \end{subfigure}
    \caption{\textbf{Cross-model refined token analysis under \texttt{expr\_first}.} Each subfigure again corresponds to one model under one prompt template. Prompt reordering suppresses the explicit expression-token categories and shifts more visible mass toward \texttt{query\_token}, \texttt{facts\_value}, and the optional \texttt{derived\_assignment} category. The resulting cross-model comparison shows that the same coarse token-level motifs persist under \texttt{expr\_first}, but with larger differences in how strongly each model concentrates information at the query-facing portion of the prompt.}
    \label{fig:token_expr_first_cross_model}
\end{figure*}

At the region level, Mistral still shows a coarse staged profile in the MLP branch, shown in Figure~\ref{fig:mistral_mlp_panel}. The attention-region outputs indicate an even stronger late Query Region than the MLP branch: on one-hop \texttt{facts\_first}, the Query Region late-band BCR rises from $0.15$ in MLP ablation to $0.27$ in attention ablation. This suggests that Mistral may rely more heavily on late routing and integration than the main Qwen3 models.

\begin{figure*}[t]
    \centering
    \begin{subfigure}[b]{0.48\textwidth}
        \centering
        \includegraphics[width=\linewidth]{./figures_experiments/reports_mistral/mlp_analysis_facts_first/comparison/band_metrics_panel.png}
        \caption{\texttt{facts\_first}}
        \label{fig:mistral_mlp_facts_first}
    \end{subfigure}
    \hfill
    \begin{subfigure}[b]{0.48\textwidth}
        \centering
        \includegraphics[width=\linewidth]{./figures_experiments/reports_mistral/mlp_analysis_expr_first/comparison/band_metrics_panel.png}
        \caption{\texttt{expr\_first}}
        \label{fig:mistral_mlp_expr_first}
    \end{subfigure}
    \caption{\textbf{Mistral-7B-Instruct-v0.1: staged computation under two prompt orders.} Band-level MLP ablation summaries rendered from the completed Mistral runs. The same broad early-facts / middle-expression / late-query structure remains visible, but the late Query Region is heavier than in Qwen3, consistent with a stronger dependence on final integration.}
    \label{fig:mistral_mlp_panel}
\end{figure*}

\begin{figure*}[t]
      \centering
      \begin{subfigure}[b]{0.47\textwidth}
          \centering
          \includegraphics[width=\linewidth]{./figures_experiments/reports_mistral/mlp_analysis_facts_first/comparison/band_metrics_panel.png}
          \caption{\texttt{facts\_first}: MLP}
          \label{fig:mistral_band_ff_mlp}
      \end{subfigure}
      \hfill
      \begin{subfigure}[b]{0.47\textwidth}
          \centering
          \includegraphics[width=\linewidth]{./figures_experiments/reports_mistral/mlp_analysis_expr_first/comparison/band_metrics_panel.png}
          \caption{\texttt{expr\_first}: MLP}
          \label{fig:mistral_band_ef_mlp}
      \end{subfigure}

      \vspace{0.8em}

      \begin{subfigure}[b]{0.47\textwidth}
          \centering
          \includegraphics[width=\linewidth]{./figures_experiments/reports_mistral/attn_analysis_facts_first/comparison/band_metrics_panel.png}
          \caption{\texttt{facts\_first}: Attention}
          \label{fig:mistral_band_ff_attn}
      \end{subfigure}
      \hfill
      \begin{subfigure}[b]{0.47\textwidth}
          \centering
          \includegraphics[width=\linewidth]{./figures_experiments/reports_mistral/attn_analysis_expr_first/comparison/band_metrics_panel.png}
          \caption{\texttt{expr\_first}: Attention}
          \label{fig:mistral_band_ef_attn}
      \end{subfigure}

      \vspace{0.8em}

      \begin{subfigure}[b]{0.47\textwidth}
          \centering
          \includegraphics[width=\linewidth]{./figures_experiments/reports_mistral/attn_mlp_analysis_facts_first/comparison/band_metrics_panel.png}
          \caption{\texttt{facts\_first}: Attention+MLP}
          \label{fig:mistral_band_ff_attn_mlp}
      \end{subfigure}
      \hfill
      \begin{subfigure}[b]{0.47\textwidth}
          \centering
          \includegraphics[width=\linewidth]{./figures_experiments/reports_mistral/attn_mlp_analysis_expr_first/comparison/band_metrics_panel.png}
          \caption{\texttt{expr\_first}: Attention+MLP}
          \label{fig:mistral_band_ef_attn_mlp}
      \end{subfigure}

      \caption{\textbf{Mistral-7B-Instruct-v0.1: band-level overviews for all three region-patching analyses under both prompt orders.} Rows correspond to MLP, attention, and joint
  attention+MLP ablations; columns correspond to \texttt{facts\_first} and \texttt{expr\_first}. Across all six panels, Mistral preserves a coarse staged template, but remains
  substantially more query-heavy than the main Qwen3 models. Under one-hop \texttt{facts\_first}, the Query Region late-band BCR rises from $0.155$ in the MLP branch to $0.272$ in
  the attention branch, with the joint branch remaining intermediate at $0.233$; meanwhile the Facts and Expression regions remain strongly early-dominant (for example, one-hop
  early-band BCRs are $0.782/0.713/0.782$ for Facts and $0.852/0.775/0.845$ for Expression across MLP/Attention/Attention+MLP). Under \texttt{expr\_first}, prompt reordering does
  not remove the staged decomposition, but reallocates more mass toward earlier Expression processing: in one-hop, the Expression Region early-band BCR is $0.872$ / $0.822$ /
  $0.873$ for MLP / Attention / Attention+MLP, while the Query Region remains broad over the back half of the network, especially in the attention and joint branches (one-hop late-
  band BCR $0.170$ / $0.242$ / $0.242$; two-hop attention middle/late BCR $0.400/0.306$). Overall, the MLP branch is the most stage-like, the attention branch is the most query-
  weighted, and the joint branch partially restores sharper early Facts/Expression concentration without eliminating Mistral's strong query-side integration.}
      \label{fig:mistral_band_overview}
  \end{figure*}

\paragraph{Mistral: head roles and causal validation.}
Mistral also yields the same three head roles, with $314$ classified heads under \texttt{facts\_first} (Fact-Retrieval $118$, Splitting $99$, Transmission $97$) and $307$ under \texttt{expr\_first} (Fact-Retrieval $123$, Splitting $109$, Transmission $75$). The Step 3 profile is strikingly sensitive. Base held-out accuracy is $99.4\%$ / $99.8\%$ for \texttt{facts\_first} / \texttt{expr\_first}, but $k=32$ role-targeted ablations reduce this to $41.2$--$41.4\%$ and $43.2$--$44.8\%$, respectively. Unlike Qwen3, however, the Random-Outside control is not benign: it still drops accuracy to $88.4\%$ and $85.3\%$. We interpret this as evidence that Mistral's performant head subset is highly sensitive overall, but less cleanly isolated from the surrounding pool than in Qwen3. The complete Mistral Step 3 accuracy tables are reported in Tables~\ref{tab:mistral_ff_full_accuracy} and~\ref{tab:mistral_ef_full_accuracy}.

\begin{figure*}[t]
    \centering
    \begin{subfigure}[b]{0.48\textwidth}
        \centering
        \includegraphics[width=\linewidth]{./figures_experiments/reports_mistral/heads_analysis_facts_first/taxonomy/head_taxonomy_line_chart.png}
        \caption{\texttt{facts\_first}: role distribution}
        \label{fig:mistral_heads_facts_tax}
    \end{subfigure}
    \hfill
    \begin{subfigure}[b]{0.48\textwidth}
        \centering
        \includegraphics[width=\linewidth]{./figures_experiments/reports_mistral/heads_analysis_expr_first/taxonomy/head_taxonomy_line_chart.png}
        \caption{\texttt{expr\_first}: role distribution}
        \label{fig:mistral_heads_expr_tax}
    \end{subfigure}

    \vspace{0.8em}

    \begin{subfigure}[b]{0.48\textwidth}
        \centering
        \includegraphics[width=\linewidth]{./figures_experiments/reports_mistral/heads_analysis_facts_first/validation/plots/pd_abs_ratio_curve.png}
        \caption{\texttt{facts\_first}: validation curve}
        \label{fig:mistral_heads_facts_val}
    \end{subfigure}
    \hfill
    \begin{subfigure}[b]{0.48\textwidth}
        \centering
        \includegraphics[width=\linewidth]{./figures_experiments/reports_mistral/heads_analysis_expr_first/validation/plots/pd_abs_ratio_curve.png}
        \caption{\texttt{expr\_first}: validation curve}
        \label{fig:mistral_heads_expr_val}
    \end{subfigure}
    \caption{\textbf{Mistral-7B-Instruct-v0.1: specialized head roles under two prompt orders.} Top row: layer-wise role distributions for the completed Mistral head-analysis runs. Bottom row: Step 3 validation curves plotting $|\mathrm{dPD}|/|\mathrm{PD}_{\mathrm{clean}}|$ against the number of ablated heads. The three roles remain recoverable, but both role-targeted and random-outside ablations are more damaging than in Qwen3, indicating a highly sensitive yet less cleanly separable head-level substrate.}
    \label{fig:mistral_heads_panel}
\end{figure*}

\paragraph{Current cross-model hypothesis.}
Taken together, the Llama and Mistral extensions strengthen the claim that the high-level motifs we identify in Qwen3 --- staged region sensitivity, prompt-order-dependent reallocation, and reusable head roles --- are not unique to a single model family. At the same time, both extensions look less cleanly factorized than Qwen3. Llama shows weaker token-level concentration and strong improvement under random-outside head ablation; Mistral shows extremely strong head sensitivity overall, including for the random control. One possible interpretation is that Qwen3 implements a sharper reasoning scaffold, whereas Llama and Mistral rely on noisier mixtures of semantic computation and output-position heuristics.

\clearpage
\onecolumn
\subsection{Additional Band-Composition and Region Galleries}
\label{sec:appendix_extra_region_gallery}

\captionsetup[figure]{font=small}
\captionsetup[subfigure]{font=small}

The figures below extend the staged-computation analysis from the main text through full region-wise curves rather than through band summaries. The retained \texttt{expr\_first} galleries show the entire one-hop/two-hop profile for each region and branch, which makes two points visible at layer resolution: prompt reordering moves Expression earlier, and Query remains the main integration sink even after that reordering.

\paragraph{How to read these region figures.}
Each panel should be read as a layer-wise sensitivity profile for one semantic region. A sharp early peak indicates a concentrated front-of-network computation stage; a broader middle or late plateau indicates that the same region remains active deeper into the network. The most informative comparison is within a branch: Facts versus Expression versus Query under the same model and prompt order.

%\begin{figure*}[p]
%    \centering
%    \begin{subfigure}[b]{0.29\textwidth}
%        \centering
%        \includegraphics[width=\linewidth]{./figures_experiments/reports/mlp_analysis/comparison/bcr_stacked.png}
%        \caption{\textit{Qwen3-8B}: MLP}
%    \end{subfigure}
%    \hfill
%    \begin{subfigure}[b]{0.29\textwidth}
%        \centering
%        \includegraphics[width=\linewidth]{./figures_experiments/reports/attn_analysis/comparison/bcr_stacked.png}
%        \caption{\textit{Qwen3-8B}: Attention}
%    \end{subfigure}
%    \hfill
%    \begin{subfigure}[b]{0.29\textwidth}
%        \centering
%        \includegraphics[width=\linewidth]{./figures_experiments/reports/attn_mlp_analysis/comparison/bcr_stacked.png}
%        \caption{\textit{Qwen3-8B}: Attention+MLP}
%    \end{subfigure}
%
%    \vspace{0.2em}
%
%    \begin{subfigure}[b]{0.29\textwidth}
%        \centering
%        \includegraphics[width=\linewidth]{./figures_experiments/reports/mlp_analysis/comparison_14B/bcr_stacked.png}
%        \caption{\textit{Qwen3-14B}: MLP}
%    \end{subfigure}
%    \hfill
%    \begin{subfigure}[b]{0.29\textwidth}
%        \centering
%        \includegraphics[width=\linewidth]{./figures_experiments/reports/attn_analysis/comparison_14B/bcr_stacked.png}
%        \caption{\textit{Qwen3-14B}: Attention}
%    \end{subfigure}
%    \hfill
%    \begin{subfigure}[b]{0.29\textwidth}
%        \centering
%        \includegraphics[width=\linewidth]{./figures_experiments/reports/attn_mlp_analysis/comparison_14B/bcr_stacked.png}
%        \caption{\textit{Qwen3-14B}: Attention+MLP}
%    \end{subfigure}
%    \caption{\textbf{Qwen3 \texttt{facts\_first}: stacked BCR summaries across scales and branches.} Each panel shows the early/middle/late contribution decomposition within the Facts, Expression, and Query regions. The MLP and joint branches maintain the clearest early-Facts profile, while the attention branch spreads more mass into middle-layer Expression and Query routing. The 14B results preserve the same ordering while reducing some of the 8B late-query mass.}
%    \label{fig:qwen3_bcr_stacked_facts}
%\end{figure*}

%\begin{figure*}[p]
%    \centering
%    \begin{subfigure}[b]{0.29\textwidth}
%        \centering
%        \includegraphics[width=\linewidth]{./figures_experiments/reports/mlp_analysis_expr_first/comparison_expr_first/bcr_stacked.png}
%        \caption{\textit{Qwen3-8B}: MLP}
%    \end{subfigure}
%    \hfill
%    \begin{subfigure}[b]{0.29\textwidth}
%        \centering
%        \includegraphics[width=\linewidth]{./figures_experiments/reports/attn_analysis_expr_first/comparison_expr_first/bcr_stacked.png}
%        \caption{\textit{Qwen3-8B}: Attention}
%    \end{subfigure}
%    \hfill
%    \begin{subfigure}[b]{0.29\textwidth}
%        \centering
%        \includegraphics[width=\linewidth]{./figures_experiments/reports/attn_mlp_analysis_expr_first/comparison_expr_first/bcr_stacked.png}
%        \caption{\textit{Qwen3-8B}: Attention+MLP}
%    \end{subfigure}
%
%    \vspace{0.2em}
%
%    \begin{subfigure}[b]{0.29\textwidth}
%        \centering
%        \includegraphics[width=\linewidth]{./figures_experiments/reports/mlp_analysis_expr_first/comparison_expr_first_14B/bcr_stacked.png}
%        \caption{\textit{Qwen3-14B}: MLP}
%    \end{subfigure}
%    \hfill
%    \begin{subfigure}[b]{0.29\textwidth}
%        \centering
%        \includegraphics[width=\linewidth]{./figures_experiments/reports/attn_analysis_expr_first/comparison_expr_first_14B/bcr_stacked.png}
%        \caption{\textit{Qwen3-14B}: Attention}
%    \end{subfigure}
%    \hfill
%    \begin{subfigure}[b]{0.29\textwidth}
%        \centering
%        \includegraphics[width=\linewidth]{./figures_experiments/reports/attn_mlp_analysis_expr_first/comparison_expr_first_14B/bcr_stacked.png}
%        \caption{\textit{Qwen3-14B}: Attention+MLP}
%    \end{subfigure}
%    \caption{\textbf{Qwen3 \texttt{expr\_first}: stacked BCR summaries across scales and branches.} Compared with Figure~\ref{fig:qwen3_bcr_stacked_facts}, the Expression Region acquires much more early-band mass, while the Facts Region becomes less exclusively early-layer concentrated. The effect is strongest in the attention and joint branches, but every branch still preserves a distinct query-facing aggregation component rather than collapsing into a single early computation stage.}
%    \label{fig:qwen3_bcr_stacked_expr}
%\end{figure*}

\begin{figure*}[t]
    \centering
    \begin{subfigure}[b]{0.285\textwidth}
        \centering
        \includegraphics[width=\linewidth]{./figures_experiments/reports/mlp_analysis_expr_first/mlp_8b_one_hop_expr_first/facts_region/mlp_facts_region_scores.png}
        \caption{Facts (One-hop)}
    \end{subfigure}
    \hfill
    \begin{subfigure}[b]{0.285\textwidth}
        \centering
        \includegraphics[width=\linewidth]{./figures_experiments/reports/mlp_analysis_expr_first/mlp_8b_one_hop_expr_first/expression_region/mlp_expression_region_scores.png}
        \caption{Expression (One-hop)}
    \end{subfigure}
    \hfill
    \begin{subfigure}[b]{0.285\textwidth}
        \centering
        \includegraphics[width=\linewidth]{./figures_experiments/reports/mlp_analysis_expr_first/mlp_8b_one_hop_expr_first/query_region/mlp_query_region_scores.png}
        \caption{Query (One-hop)}
    \end{subfigure}

    \vspace{0.2em}

    \begin{subfigure}[b]{0.285\textwidth}
        \centering
        \includegraphics[width=\linewidth]{./figures_experiments/reports/mlp_analysis_expr_first/mlp_8b_two_hop_expr_first/facts_region/mlp_facts_region_scores.png}
        \caption{Facts (Two-hop)}
    \end{subfigure}
    \hfill
    \begin{subfigure}[b]{0.285\textwidth}
        \centering
        \includegraphics[width=\linewidth]{./figures_experiments/reports/mlp_analysis_expr_first/mlp_8b_two_hop_expr_first/expression_region/mlp_expression_region_scores.png}
        \caption{Expression (Two-hop)}
    \end{subfigure}
    \hfill
    \begin{subfigure}[b]{0.285\textwidth}
        \centering
        \includegraphics[width=\linewidth]{./figures_experiments/reports/mlp_analysis_expr_first/mlp_8b_two_hop_expr_first/query_region/mlp_query_region_scores.png}
        \caption{Query (Two-hop)}
    \end{subfigure}
    \caption{\textbf{Complete \textit{Qwen3-8B} MLP region gallery under \texttt{expr\_first}.} The full one-hop/two-hop curves show that reordering the prompt moves expression processing earlier without removing the later query-facing integration stage. In one-hop, the Facts Region no longer dominates the earliest layers as strongly as under \texttt{facts\_first}, while the Expression Region becomes the sharpest early-layer signal. In two-hop, the same ordering survives but the Query Region broadens into a stronger middle-to-late bridge.}
    \label{fig:qwen3_8b_mlp_expr_gallery}
\end{figure*}

\begin{figure*}[t]
    \centering
    \begin{subfigure}[b]{0.285\textwidth}
        \centering
        \includegraphics[width=\linewidth]{./figures_experiments/reports/mlp_analysis_expr_first/mlp_14b_one_hop_expr_first/facts_region/mlp_facts_region_scores.png}
        \caption{Facts (One-hop)}
    \end{subfigure}
    \hfill
    \begin{subfigure}[b]{0.285\textwidth}
        \centering
        \includegraphics[width=\linewidth]{./figures_experiments/reports/mlp_analysis_expr_first/mlp_14b_one_hop_expr_first/expression_region/mlp_expression_region_scores.png}
        \caption{Expression (One-hop)}
    \end{subfigure}
    \hfill
    \begin{subfigure}[b]{0.285\textwidth}
        \centering
        \includegraphics[width=\linewidth]{./figures_experiments/reports/mlp_analysis_expr_first/mlp_14b_one_hop_expr_first/query_region/mlp_query_region_scores.png}
        \caption{Query (One-hop)}
    \end{subfigure}

    \vspace{0.2em}

    \begin{subfigure}[b]{0.285\textwidth}
        \centering
        \includegraphics[width=\linewidth]{./figures_experiments/reports/mlp_analysis_expr_first/mlp_14b_two_hop_expr_first/facts_region/mlp_facts_region_scores.png}
        \caption{Facts (Two-hop)}
    \end{subfigure}
    \hfill
    \begin{subfigure}[b]{0.285\textwidth}
        \centering
        \includegraphics[width=\linewidth]{./figures_experiments/reports/mlp_analysis_expr_first/mlp_14b_two_hop_expr_first/expression_region/mlp_expression_region_scores.png}
        \caption{Expression (Two-hop)}
    \end{subfigure}
    \hfill
    \begin{subfigure}[b]{0.285\textwidth}
        \centering
        \includegraphics[width=\linewidth]{./figures_experiments/reports/mlp_analysis_expr_first/mlp_14b_two_hop_expr_first/query_region/mlp_query_region_scores.png}
        \caption{Query (Two-hop)}
    \end{subfigure}
    \caption{\textbf{Complete \textit{Qwen3-14B} MLP region gallery under \texttt{expr\_first}.} The larger model exhibits the same prompt-order effect as 8B, but more cleanly: one-hop Expression is the most concentrated early-layer signal, Facts become broader and less purely early-band dominated, and Query remains the terminal aggregation region. Two-hop again lowers the absolute sharpness of every region while preserving the same coarse temporal ordering.}
    \label{fig:qwen3_14b_mlp_expr_gallery}
\end{figure*}

\begin{figure*}[t]
    \centering
    \begin{subfigure}[b]{0.285\textwidth}
        \centering
        \includegraphics[width=\linewidth]{./figures_experiments/reports/attn_analysis_expr_first/attn_8b_one_hop_expr_first/facts_region/attn_facts_region_scores.png}
        \caption{Facts (One-hop)}
    \end{subfigure}
    \hfill
    \begin{subfigure}[b]{0.285\textwidth}
        \centering
        \includegraphics[width=\linewidth]{./figures_experiments/reports/attn_analysis_expr_first/attn_8b_one_hop_expr_first/expression_region/attn_expression_region_scores.png}
        \caption{Expression (One-hop)}
    \end{subfigure}
    \hfill
    \begin{subfigure}[b]{0.285\textwidth}
        \centering
        \includegraphics[width=\linewidth]{./figures_experiments/reports/attn_analysis_expr_first/attn_8b_one_hop_expr_first/query_region/attn_query_region_scores.png}
        \caption{Query (One-hop)}
    \end{subfigure}

    \vspace{0.2em}

    \begin{subfigure}[b]{0.285\textwidth}
        \centering
        \includegraphics[width=\linewidth]{./figures_experiments/reports/attn_analysis_expr_first/attn_8b_two_hop_expr_first/facts_region/attn_facts_region_scores.png}
        \caption{Facts (Two-hop)}
    \end{subfigure}
    \hfill
    \begin{subfigure}[b]{0.285\textwidth}
        \centering
        \includegraphics[width=\linewidth]{./figures_experiments/reports/attn_analysis_expr_first/attn_8b_two_hop_expr_first/expression_region/attn_expression_region_scores.png}
        \caption{Expression (Two-hop)}
    \end{subfigure}
    \hfill
    \begin{subfigure}[b]{0.285\textwidth}
        \centering
        \includegraphics[width=\linewidth]{./figures_experiments/reports/attn_analysis_expr_first/attn_8b_two_hop_expr_first/query_region/attn_query_region_scores.png}
        \caption{Query (Two-hop)}
    \end{subfigure}
    \caption{\textbf{Complete \textit{Qwen3-8B} attention-region gallery under \texttt{expr\_first}.} Relative to the MLP branch, attention patching shifts more mass toward middle-layer routing. The Facts Region is no longer purely early-layer dominated in one-hop, while the Query Region broadens into a larger middle-to-late control surface. This supports the interpretation that attention primarily redistributes and routes information once prompt order front-loads the expression.}
    \label{fig:qwen3_8b_attn_expr_gallery}
\end{figure*}

\begin{figure*}[t]
    \centering
    \begin{subfigure}[b]{0.285\textwidth}
        \centering
        \includegraphics[width=\linewidth]{./figures_experiments/reports/attn_analysis_expr_first/attn_14b_one_hop_expr_first/facts_region/attn_facts_region_scores.png}
        \caption{Facts (One-hop)}
    \end{subfigure}
    \hfill
    \begin{subfigure}[b]{0.285\textwidth}
        \centering
        \includegraphics[width=\linewidth]{./figures_experiments/reports/attn_analysis_expr_first/attn_14b_one_hop_expr_first/expression_region/attn_expression_region_scores.png}
        \caption{Expression (One-hop)}
    \end{subfigure}
    \hfill
    \begin{subfigure}[b]{0.285\textwidth}
        \centering
        \includegraphics[width=\linewidth]{./figures_experiments/reports/attn_analysis_expr_first/attn_14b_one_hop_expr_first/query_region/attn_query_region_scores.png}
        \caption{Query (One-hop)}
    \end{subfigure}

    \vspace{0.2em}

    \begin{subfigure}[b]{0.285\textwidth}
        \centering
        \includegraphics[width=\linewidth]{./figures_experiments/reports/attn_analysis_expr_first/attn_14b_two_hop_expr_first/facts_region/attn_facts_region_scores.png}
        \caption{Facts (Two-hop)}
    \end{subfigure}
    \hfill
    \begin{subfigure}[b]{0.285\textwidth}
        \centering
        \includegraphics[width=\linewidth]{./figures_experiments/reports/attn_analysis_expr_first/attn_14b_two_hop_expr_first/expression_region/attn_expression_region_scores.png}
        \caption{Expression (Two-hop)}
    \end{subfigure}
    \hfill
    \begin{subfigure}[b]{0.285\textwidth}
        \centering
        \includegraphics[width=\linewidth]{./figures_experiments/reports/attn_analysis_expr_first/attn_14b_two_hop_expr_first/query_region/attn_query_region_scores.png}
        \caption{Query (Two-hop)}
    \end{subfigure}
    \caption{\textbf{Complete \textit{Qwen3-14B} attention-region gallery under \texttt{expr\_first}.} The larger attention branch shows the same structure as 8B, but with an even clearer mid-layer Facts bulge in one-hop and a stronger middle-band Query component in two-hop. These curves make explicit that attention is the branch most sensitive to prompt-layout changes, especially in how it redistributes factual information away from the earliest layers.}
    \label{fig:qwen3_14b_attn_expr_gallery}
\end{figure*}

\begin{figure*}[t]
    \centering
    \begin{subfigure}[b]{0.285\textwidth}
        \centering
        \includegraphics[width=\linewidth]{./figures_experiments/reports/attn_mlp_analysis_expr_first/attn_mlp_8b_one_hop_expr_first/facts_region/attn_mlp_facts_region_scores.png}
        \caption{Facts (One-hop)}
    \end{subfigure}
    \hfill
    \begin{subfigure}[b]{0.285\textwidth}
        \centering
        \includegraphics[width=\linewidth]{./figures_experiments/reports/attn_mlp_analysis_expr_first/attn_mlp_8b_one_hop_expr_first/expression_region/attn_mlp_expression_region_scores.png}
        \caption{Expression (One-hop)}
    \end{subfigure}
    \hfill
    \begin{subfigure}[b]{0.285\textwidth}
        \centering
        \includegraphics[width=\linewidth]{./figures_experiments/reports/attn_mlp_analysis_expr_first/attn_mlp_8b_one_hop_expr_first/query_region/attn_mlp_query_region_scores.png}
        \caption{Query (One-hop)}
    \end{subfigure}

    \vspace{0.2em}

    \begin{subfigure}[b]{0.285\textwidth}
        \centering
        \includegraphics[width=\linewidth]{./figures_experiments/reports/attn_mlp_analysis_expr_first/attn_mlp_8b_two_hop_expr_first/facts_region/attn_mlp_facts_region_scores.png}
        \caption{Facts (Two-hop)}
    \end{subfigure}
    \hfill
    \begin{subfigure}[b]{0.285\textwidth}
        \centering
        \includegraphics[width=\linewidth]{./figures_experiments/reports/attn_mlp_analysis_expr_first/attn_mlp_8b_two_hop_expr_first/expression_region/attn_mlp_expression_region_scores.png}
        \caption{Expression (Two-hop)}
    \end{subfigure}
    \hfill
    \begin{subfigure}[b]{0.285\textwidth}
        \centering
        \includegraphics[width=\linewidth]{./figures_experiments/reports/attn_mlp_analysis_expr_first/attn_mlp_8b_two_hop_expr_first/query_region/attn_mlp_query_region_scores.png}
        \caption{Query (Two-hop)}
    \end{subfigure}
    \caption{\textbf{Complete \textit{Qwen3-8B} joint attention+MLP gallery under \texttt{expr\_first}.} The joint intervention remains more stage-like than the pure attention branch. Facts and expressions both retain strong early peaks, while Query continues to broaden over the middle and late bands. The figure therefore shows that coupling routing with nonlinear transformation partially stabilizes the staged decomposition under prompt reordering.}
    \label{fig:qwen3_8b_attn_mlp_expr_gallery}
\end{figure*}

\begin{figure*}[t]
    \centering
    \begin{subfigure}[b]{0.285\textwidth}
        \centering
        \includegraphics[width=\linewidth]{./figures_experiments/reports/attn_mlp_analysis_expr_first/attn_mlp_14b_one_hop_expr_first/facts_region/attn_mlp_facts_region_scores.png}
        \caption{Facts (One-hop)}
    \end{subfigure}
    \hfill
    \begin{subfigure}[b]{0.285\textwidth}
        \centering
        \includegraphics[width=\linewidth]{./figures_experiments/reports/attn_mlp_analysis_expr_first/attn_mlp_14b_one_hop_expr_first/expression_region/attn_mlp_expression_region_scores.png}
        \caption{Expression (One-hop)}
    \end{subfigure}
    \hfill
    \begin{subfigure}[b]{0.285\textwidth}
        \centering
        \includegraphics[width=\linewidth]{./figures_experiments/reports/attn_mlp_analysis_expr_first/attn_mlp_14b_one_hop_expr_first/query_region/attn_mlp_query_region_scores.png}
        \caption{Query (One-hop)}
    \end{subfigure}

    \vspace{0.2em}

    \begin{subfigure}[b]{0.285\textwidth}
        \centering
        \includegraphics[width=\linewidth]{./figures_experiments/reports/attn_mlp_analysis_expr_first/attn_mlp_14b_two_hop_expr_first/facts_region/attn_mlp_facts_region_scores.png}
        \caption{Facts (Two-hop)}
    \end{subfigure}
    \hfill
    \begin{subfigure}[b]{0.285\textwidth}
        \centering
        \includegraphics[width=\linewidth]{./figures_experiments/reports/attn_mlp_analysis_expr_first/attn_mlp_14b_two_hop_expr_first/expression_region/attn_mlp_expression_region_scores.png}
        \caption{Expression (Two-hop)}
    \end{subfigure}
    \hfill
    \begin{subfigure}[b]{0.285\textwidth}
        \centering
        \includegraphics[width=\linewidth]{./figures_experiments/reports/attn_mlp_analysis_expr_first/attn_mlp_14b_two_hop_expr_first/query_region/attn_mlp_query_region_scores.png}
        \caption{Query (Two-hop)}
    \end{subfigure}
    \caption{\textbf{Complete \textit{Qwen3-14B} joint attention+MLP gallery under \texttt{expr\_first}.} The 14B joint branch amplifies the one-hop Expression peak while keeping Facts strongly early-layer concentrated. Two-hop curves remain broader, but still preserve the same order: Facts and Expression dominate the front half of the network, and Query remains the principal sink for later integration.}
    \label{fig:qwen3_14b_attn_mlp_expr_gallery}
\end{figure*}

%\begin{figure*}[p]
%    \centering
%    \begin{subfigure}[b]{0.29\textwidth}
%        \centering
%        \includegraphics[width=\linewidth]{./figures_experiments/reports_llama3/mlp_analysis_facts_first/comparison/bcr_stacked.png}
%        \caption{Llama: MLP}
%    \end{subfigure}
%    \hfill
%    \begin{subfigure}[b]{0.29\textwidth}
%        \centering
%        \includegraphics[width=\linewidth]{./figures_experiments/reports_llama3/attn_analysis_facts_first/comparison/bcr_stacked.png}
%        \caption{Llama: Attention}
%    \end{subfigure}
%    \hfill
%    \begin{subfigure}[b]{0.29\textwidth}
%        \centering
%        \includegraphics[width=\linewidth]{./figures_experiments/reports_llama3/attn_mlp_analysis_facts_first/comparison/bcr_stacked.png}
%        \caption{Llama: Attention+MLP}
%    \end{subfigure}
%
%    \vspace{0.2em}
%
%    \begin{subfigure}[b]{0.29\textwidth}
%        \centering
%        \includegraphics[width=\linewidth]{./figures_experiments/reports_mistral/mlp_analysis_facts_first/comparison/bcr_stacked.png}
%        \caption{Mistral: MLP}
%    \end{subfigure}
%    \hfill
%    \begin{subfigure}[b]{0.29\textwidth}
%        \centering
%        \includegraphics[width=\linewidth]{./figures_experiments/reports_mistral/attn_analysis_facts_first/comparison/bcr_stacked.png}
%        \caption{Mistral: Attention}
%    \end{subfigure}
%    \hfill
%    \begin{subfigure}[b]{0.29\textwidth}
%        \centering
%        \includegraphics[width=\linewidth]{./figures_experiments/reports_mistral/attn_mlp_analysis_facts_first/comparison/bcr_stacked.png}
%        \caption{Mistral: Attention+MLP}
%    \end{subfigure}
%    \caption{\textbf{Cross-model stacked BCR summaries under \texttt{facts\_first}.} Both Llama and Mistral preserve the same coarse facts-to-expression-to-query ordering, but the stacked decompositions reveal much heavier query-facing mass than in Qwen3. Llama spreads more mass into the Query Region across all three branches, while Mistral retains a comparatively strong middle-band Query component even in the MLP branch.}
%    \label{fig:cross_model_bcr_stacked_facts}
%\end{figure*}

%\begin{figure*}[p]
%    \centering
%    \begin{subfigure}[b]{0.29\textwidth}
%        \centering
%        \includegraphics[width=\linewidth]{./figures_experiments/reports_llama3/mlp_analysis_expr_first/comparison/bcr_stacked.png}
%        \caption{Llama: MLP}
%    \end{subfigure}
%    \hfill
%    \begin{subfigure}[b]{0.29\textwidth}
%        \centering
%        \includegraphics[width=\linewidth]{./figures_experiments/reports_llama3/attn_analysis_expr_first/comparison/bcr_stacked.png}
%        \caption{Llama: Attention}
%    \end{subfigure}
%    \hfill
%    \begin{subfigure}[b]{0.29\textwidth}
%        \centering
%        \includegraphics[width=\linewidth]{./figures_experiments/reports_llama3/attn_mlp_analysis_expr_first/comparison/bcr_stacked.png}
%        \caption{Llama: Attention+MLP}
%    \end{subfigure}
%
%    \vspace{0.2em}
%
%    \begin{subfigure}[b]{0.29\textwidth}
%        \centering
%        \includegraphics[width=\linewidth]{./figures_experiments/reports_mistral/mlp_analysis_expr_first/comparison/bcr_stacked.png}
%        \caption{Mistral: MLP}
%    \end{subfigure}
%    \hfill
%    \begin{subfigure}[b]{0.29\textwidth}
%        \centering
%        \includegraphics[width=\linewidth]{./figures_experiments/reports_mistral/attn_analysis_expr_first/comparison/bcr_stacked.png}
%        \caption{Mistral: Attention}
%    \end{subfigure}
%    \hfill
%    \begin{subfigure}[b]{0.29\textwidth}
%        \centering
%        \includegraphics[width=\linewidth]{./figures_experiments/reports_mistral/attn_mlp_analysis_expr_first/comparison/bcr_stacked.png}
%        \caption{Mistral: Attention+MLP}
%    \end{subfigure}
%    \caption{\textbf{Cross-model stacked BCR summaries under \texttt{expr\_first}.} Prompt reordering again increases early Expression mass, but the cross-model gap persists. Llama keeps a notably large early-plus-late Query component, while Mistral preserves a broad Query distribution across the back half of the network. These summaries reinforce the interpretation that both non-Qwen families rely more heavily on query-facing scaffolding during final integration.}
%    \label{fig:cross_model_bcr_stacked_expr}
%\end{figure*}

\paragraph{How to compare the cross-model region figures.}
In the cross-model panels below, the safest comparison is to hold the branch fixed and then compare families: first compare Qwen3/Llama/Mistral within MLP, then within attention, then within the joint branch. A family with more late Query mass is relying more heavily on query-facing scaffolding during final integration, whereas a family with cleaner early Facts and middle Expression concentration is implementing a sharper staged decomposition.

The full cross-model MLP galleries sharpen the same qualitative point. In Llama, even the core MLP branch allocates unusually large mass to the Query Region: under one-hop \texttt{facts\_first}, the query profile is split across early ($0.51$), middle ($0.27$), and late ($0.22$) bands rather than being primarily a late aggregator. Mistral looks somewhat cleaner than Llama, but still remains more query-heavy than Qwen3, with one-hop \texttt{facts\_first} query shares of $0.48$ / $0.37$ / $0.15$ (early / middle / late). The full curves below make these differences visible at layer resolution rather than only through compact band summaries.

\begin{figure*}[t]
    \centering
    \begin{subfigure}[b]{0.285\textwidth}
        \centering
        \includegraphics[width=\linewidth]{./figures_experiments/reports_llama3/mlp_analysis_facts_first/mlp_one_hop_facts_first/facts_region/mlp_facts_region_scores.png}
        \caption{Facts (One-hop)}
    \end{subfigure}
    \hfill
    \begin{subfigure}[b]{0.285\textwidth}
        \centering
        \includegraphics[width=\linewidth]{./figures_experiments/reports_llama3/mlp_analysis_facts_first/mlp_one_hop_facts_first/expression_region/mlp_expression_region_scores.png}
        \caption{Expression (One-hop)}
    \end{subfigure}
    \hfill
    \begin{subfigure}[b]{0.285\textwidth}
        \centering
        \includegraphics[width=\linewidth]{./figures_experiments/reports_llama3/mlp_analysis_facts_first/mlp_one_hop_facts_first/query_region/mlp_query_region_scores.png}
        \caption{Query (One-hop)}
    \end{subfigure}

    \vspace{0.2em}

    \begin{subfigure}[b]{0.285\textwidth}
        \centering
        \includegraphics[width=\linewidth]{./figures_experiments/reports_llama3/mlp_analysis_facts_first/mlp_two_hop_facts_first/facts_region/mlp_facts_region_scores.png}
        \caption{Facts (Two-hop)}
    \end{subfigure}
    \hfill
    \begin{subfigure}[b]{0.285\textwidth}
        \centering
        \includegraphics[width=\linewidth]{./figures_experiments/reports_llama3/mlp_analysis_facts_first/mlp_two_hop_facts_first/expression_region/mlp_expression_region_scores.png}
        \caption{Expression (Two-hop)}
    \end{subfigure}
    \hfill
    \begin{subfigure}[b]{0.285\textwidth}
        \centering
        \includegraphics[width=\linewidth]{./figures_experiments/reports_llama3/mlp_analysis_facts_first/mlp_two_hop_facts_first/query_region/mlp_query_region_scores.png}
        \caption{Query (Two-hop)}
    \end{subfigure}
    \caption{\textbf{Complete Llama MLP region gallery under \texttt{facts\_first}.} Llama preserves the broad staged template but with a much heavier Query Region than Qwen3. The query curve is already substantial in one-hop middle layers and remains comparably strong in late layers, which is consistent with a more output-position-centered computation strategy.}
    \label{fig:llama_mlp_facts_gallery}
\end{figure*}

\begin{figure*}[t]
    \centering
    \begin{subfigure}[b]{0.285\textwidth}
        \centering
        \includegraphics[width=\linewidth]{./figures_experiments/reports_llama3/mlp_analysis_expr_first/mlp_one_hop_expr_first/facts_region/mlp_facts_region_scores.png}
        \caption{Facts (One-hop)}
    \end{subfigure}
    \hfill
    \begin{subfigure}[b]{0.285\textwidth}
        \centering
        \includegraphics[width=\linewidth]{./figures_experiments/reports_llama3/mlp_analysis_expr_first/mlp_one_hop_expr_first/expression_region/mlp_expression_region_scores.png}
        \caption{Expression (One-hop)}
    \end{subfigure}
    \hfill
    \begin{subfigure}[b]{0.285\textwidth}
        \centering
        \includegraphics[width=\linewidth]{./figures_experiments/reports_llama3/mlp_analysis_expr_first/mlp_one_hop_expr_first/query_region/mlp_query_region_scores.png}
        \caption{Query (One-hop)}
    \end{subfigure}

    \vspace{0.2em}

    \begin{subfigure}[b]{0.285\textwidth}
        \centering
        \includegraphics[width=\linewidth]{./figures_experiments/reports_llama3/mlp_analysis_expr_first/mlp_two_hop_expr_first/facts_region/mlp_facts_region_scores.png}
        \caption{Facts (Two-hop)}
    \end{subfigure}
    \hfill
    \begin{subfigure}[b]{0.285\textwidth}
        \centering
        \includegraphics[width=\linewidth]{./figures_experiments/reports_llama3/mlp_analysis_expr_first/mlp_two_hop_expr_first/expression_region/mlp_expression_region_scores.png}
        \caption{Expression (Two-hop)}
    \end{subfigure}
    \hfill
    \begin{subfigure}[b]{0.285\textwidth}
        \centering
        \includegraphics[width=\linewidth]{./figures_experiments/reports_llama3/mlp_analysis_expr_first/mlp_two_hop_expr_first/query_region/mlp_query_region_scores.png}
        \caption{Query (Two-hop)}
    \end{subfigure}
    \caption{\textbf{Complete Llama MLP region gallery under \texttt{expr\_first}.} Prompt reordering shifts expression processing earlier as expected, but the Query Region remains unusually broad and strong. Even after moving the queried expression to the front, Llama still allocates a large amount of MLP sensitivity to the query-facing suffix and answer anchor.}
    \label{fig:llama_mlp_expr_gallery}
\end{figure*}

\begin{figure*}[t]
    \centering
    \begin{subfigure}[b]{0.285\textwidth}
        \centering
        \includegraphics[width=\linewidth]{./figures_experiments/reports_mistral/mlp_analysis_facts_first/mlp_one_hop_facts_first/facts_region/mlp_facts_region_scores.png}
        \caption{Facts (One-hop)}
    \end{subfigure}
    \hfill
    \begin{subfigure}[b]{0.285\textwidth}
        \centering
        \includegraphics[width=\linewidth]{./figures_experiments/reports_mistral/mlp_analysis_facts_first/mlp_one_hop_facts_first/expression_region/mlp_expression_region_scores.png}
        \caption{Expression (One-hop)}
    \end{subfigure}
    \hfill
    \begin{subfigure}[b]{0.285\textwidth}
        \centering
        \includegraphics[width=\linewidth]{./figures_experiments/reports_mistral/mlp_analysis_facts_first/mlp_one_hop_facts_first/query_region/mlp_query_region_scores.png}
        \caption{Query (One-hop)}
    \end{subfigure}

    \vspace{0.2em}

    \begin{subfigure}[b]{0.285\textwidth}
        \centering
        \includegraphics[width=\linewidth]{./figures_experiments/reports_mistral/mlp_analysis_facts_first/mlp_two_hop_facts_first/facts_region/mlp_facts_region_scores.png}
        \caption{Facts (Two-hop)}
    \end{subfigure}
    \hfill
    \begin{subfigure}[b]{0.285\textwidth}
        \centering
        \includegraphics[width=\linewidth]{./figures_experiments/reports_mistral/mlp_analysis_facts_first/mlp_two_hop_facts_first/expression_region/mlp_expression_region_scores.png}
        \caption{Expression (Two-hop)}
    \end{subfigure}
    \hfill
    \begin{subfigure}[b]{0.285\textwidth}
        \centering
        \includegraphics[width=\linewidth]{./figures_experiments/reports_mistral/mlp_analysis_facts_first/mlp_two_hop_facts_first/query_region/mlp_query_region_scores.png}
        \caption{Query (Two-hop)}
    \end{subfigure}
    \caption{\textbf{Complete Mistral MLP region gallery under \texttt{facts\_first}.} Mistral retains a clearer early-Facts and early-Expression ordering than Llama, but again assigns more mass to the Query Region than Qwen3. The one-hop Query curve remains substantial throughout the back half of the network, and the two-hop curves further broaden this integration profile.}
    \label{fig:mistral_mlp_facts_gallery}
\end{figure*}

\begin{figure*}[t]
    \centering
    \begin{subfigure}[b]{0.285\textwidth}
        \centering
        \includegraphics[width=\linewidth]{./figures_experiments/reports_mistral/mlp_analysis_expr_first/mlp_one_hop_expr_first/facts_region/mlp_facts_region_scores.png}
        \caption{Facts (One-hop)}
    \end{subfigure}
    \hfill
    \begin{subfigure}[b]{0.285\textwidth}
        \centering
        \includegraphics[width=\linewidth]{./figures_experiments/reports_mistral/mlp_analysis_expr_first/mlp_one_hop_expr_first/expression_region/mlp_expression_region_scores.png}
        \caption{Expression (One-hop)}
    \end{subfigure}
    \hfill
    \begin{subfigure}[b]{0.285\textwidth}
        \centering
        \includegraphics[width=\linewidth]{./figures_experiments/reports_mistral/mlp_analysis_expr_first/mlp_one_hop_expr_first/query_region/mlp_query_region_scores.png}
        \caption{Query (One-hop)}
    \end{subfigure}

    \vspace{0.2em}

    \begin{subfigure}[b]{0.285\textwidth}
        \centering
        \includegraphics[width=\linewidth]{./figures_experiments/reports_mistral/mlp_analysis_expr_first/mlp_two_hop_expr_first/facts_region/mlp_facts_region_scores.png}
        \caption{Facts (Two-hop)}
    \end{subfigure}
    \hfill
    \begin{subfigure}[b]{0.285\textwidth}
        \centering
        \includegraphics[width=\linewidth]{./figures_experiments/reports_mistral/mlp_analysis_expr_first/mlp_two_hop_expr_first/expression_region/mlp_expression_region_scores.png}
        \caption{Expression (Two-hop)}
    \end{subfigure}
    \hfill
    \begin{subfigure}[b]{0.285\textwidth}
        \centering
        \includegraphics[width=\linewidth]{./figures_experiments/reports_mistral/mlp_analysis_expr_first/mlp_two_hop_expr_first/query_region/mlp_query_region_scores.png}
        \caption{Query (Two-hop)}
    \end{subfigure}
    \caption{\textbf{Complete Mistral MLP region gallery under \texttt{expr\_first}.} Mistral again shifts expression processing forward under prompt reordering, but the Query Region remains broad and relatively heavy. Compared with Qwen3, the figure suggests a weaker separation between semantic content processing and final answer integration.}
    \label{fig:mistral_mlp_expr_gallery}
\end{figure*}

\subsection{Additional Token Diagnostics}
\label{sec:appendix_extra_token_gallery}

The token diagnostics below sharpen two points left implicit in the main appendix. First, the refined mean and signed plots separate large-magnitude categories from categories that consistently move the model toward the correct answer. Second, the cross-model contrast remains stable across different normalizations. In unsigned totals, Qwen3 allocates substantially more mass to fact tokens than either non-Qwen family, whereas Llama and Mistral put relatively more mass on \texttt{query\_token} (under \texttt{facts\_first}: \textit{Qwen3-8B} \texttt{facts\_value} $2.47$ vs. \texttt{query\_token} $2.22$; Llama $0.72$ vs. $1.63$; Mistral $0.75$ vs. $1.04$). In signed totals, Qwen3 remains fact-dominant, Llama remains query-dominant, and Mistral falls between those extremes: under signed \texttt{facts\_first}, its \texttt{query\_token} and \texttt{facts\_value} totals are nearly identical ($0.731$ vs. $0.729$).

\paragraph{How to read these token figures.}
The refined \textbf{sum} plots answer a total-contribution question and therefore reward categories that accumulate substantial mass over multiple tokens, whereas the refined \textbf{mean} plots normalize by category size and answer a per-token-importance question. Signed variants preserve whether the intervention shifts the model in the correct direction, so a category can be large in unsigned magnitude but less dominant once directional consistency is required.

For Qwen3, the retained refined views are given in Figures~\ref{fig:qwen_token_refined_mean}, \ref{fig:qwen_token_refined_sum_signed}, and \ref{fig:qwen_token_refined_mean_signed}. The corresponding cross-model comparisons are collected in Figures~\ref{fig:cross_model_token_refined_mean}, \ref{fig:cross_model_token_refined_sum_signed}, and \ref{fig:cross_model_token_refined_mean_signed}.

%\begin{figure*}[p]
%    \centering
%    \begin{subfigure}[b]{0.44\textwidth}
%        \centering
%        \includegraphics[width=\linewidth]{./figures_experiments/reports/token_analysis/8b_facts_first_all_hop_raw/heatmap_simple.png}
%        \caption{\textit{Qwen3-8B}, \texttt{facts\_first}}
%    \end{subfigure}
%    \hfill
%    \begin{subfigure}[b]{0.44\textwidth}
%        \centering
%        \includegraphics[width=\linewidth]{./figures_experiments/reports/token_analysis/14b_facts_first_all_hop_raw/heatmap_simple.png}
%        \caption{\textit{Qwen3-14B}, \texttt{facts\_first}}
%    \end{subfigure}
%
%    \vspace{0.2em}
%
%    \begin{subfigure}[b]{0.44\textwidth}
%        \centering
%        \includegraphics[width=\linewidth]{./figures_experiments/reports/token_analysis/8b_expr_first_all_hop_raw/heatmap_simple.png}
%        \caption{\textit{Qwen3-8B}, \texttt{expr\_first}}
%    \end{subfigure}
%    \hfill
%    \begin{subfigure}[b]{0.44\textwidth}
%        \centering
%        \includegraphics[width=\linewidth]{./figures_experiments/reports/token_analysis/14b_expr_first_all_hop_raw/heatmap_simple.png}
%        \caption{\textit{Qwen3-14B}, \texttt{expr\_first}}
%    \end{subfigure}
%    \caption{\textbf{Qwen3 raw token heatmaps across scales and prompt orders.} The raw category-by-layer heatmaps make the same stage structure visible before refinement. \texttt{facts\_first} concentrates large early-layer effects on fact tokens, while \texttt{expr\_first} shifts more structure toward query-facing and derived-assignment categories. The 14B model keeps the same topology while producing a slightly cleaner early-to-late separation.}
%    \label{fig:qwen_token_raw_heatmaps}
%\end{figure*}

\begin{figure*}[t]
    \centering
    \begin{subfigure}[b]{0.44\textwidth}
        \centering
        \includegraphics[width=\linewidth]{./figures_experiments/reports/token_analysis/8b_facts_first_all_hop_refined/refined_by_stage_mean.png}
        \caption{\textit{Qwen3-8B}, \texttt{facts\_first}}
    \end{subfigure}
    \hfill
    \begin{subfigure}[b]{0.44\textwidth}
        \centering
        \includegraphics[width=\linewidth]{./figures_experiments/reports/token_analysis/14b_facts_first_all_hop_refined/refined_by_stage_mean.png}
        \caption{\textit{Qwen3-14B}, \texttt{facts\_first}}
    \end{subfigure}

    \vspace{0.2em}

    \begin{subfigure}[b]{0.44\textwidth}
        \centering
        \includegraphics[width=\linewidth]{./figures_experiments/reports/token_analysis/8b_expr_first_all_hop_refined/refined_by_stage_mean.png}
        \caption{\textit{Qwen3-8B}, \texttt{expr\_first}}
    \end{subfigure}
    \hfill
    \begin{subfigure}[b]{0.44\textwidth}
        \centering
        \includegraphics[width=\linewidth]{./figures_experiments/reports/token_analysis/14b_expr_first_all_hop_refined/refined_by_stage_mean.png}
        \caption{\textit{Qwen3-14B}, \texttt{expr\_first}}
    \end{subfigure}
    \caption{\textbf{Qwen3 refined mean token diagnostics.} Averaging by token rather than summing by category preserves the same hierarchy as the sum-based plots while shrinking the advantage of categories that occupy more positions. Facts and query remain the dominant token types in all four settings, which shows that the main result is not an artifact of category cardinality.}
    \label{fig:qwen_token_refined_mean}
\end{figure*}

\begin{figure*}[t]
    \centering
    \begin{subfigure}[b]{0.44\textwidth}
        \centering
        \includegraphics[width=\linewidth]{./figures_experiments/reports/token_analysis/8b_facts_first_all_hop_refined/refined_by_stage_sum_signed.png}
        \caption{\textit{Qwen3-8B}, \texttt{facts\_first}}
    \end{subfigure}
    \hfill
    \begin{subfigure}[b]{0.44\textwidth}
        \centering
        \includegraphics[width=\linewidth]{./figures_experiments/reports/token_analysis/14b_facts_first_all_hop_refined/refined_by_stage_sum_signed.png}
        \caption{\textit{Qwen3-14B}, \texttt{facts\_first}}
    \end{subfigure}

    \vspace{0.2em}

    \begin{subfigure}[b]{0.44\textwidth}
        \centering
        \includegraphics[width=\linewidth]{./figures_experiments/reports/token_analysis/8b_expr_first_all_hop_refined/refined_by_stage_sum_signed.png}
        \caption{\textit{Qwen3-8B}, \texttt{expr\_first}}
    \end{subfigure}
    \hfill
    \begin{subfigure}[b]{0.44\textwidth}
        \centering
        \includegraphics[width=\linewidth]{./figures_experiments/reports/token_analysis/14b_expr_first_all_hop_refined/refined_by_stage_sum_signed.png}
        \caption{\textit{Qwen3-14B}, \texttt{expr\_first}}
    \end{subfigure}
    \caption{\textbf{Qwen3 refined signed-sum token diagnostics.} These plots retain the sign of the token-level shift rather than using absolute magnitude. The dominant categories remain the same as in the unsigned plots, which shows that fact and query tokens are not merely high-variance sites: they push the model in a directionally meaningful way for the final decision.}
    \label{fig:qwen_token_refined_sum_signed}
\end{figure*}

\begin{figure*}[t]
    \centering
    \begin{subfigure}[b]{0.44\textwidth}
        \centering
        \includegraphics[width=\linewidth]{./figures_experiments/reports/token_analysis/8b_facts_first_all_hop_refined/refined_by_stage_mean_signed.png}
        \caption{\textit{Qwen3-8B}, \texttt{facts\_first}}
    \end{subfigure}
    \hfill
    \begin{subfigure}[b]{0.44\textwidth}
        \centering
        \includegraphics[width=\linewidth]{./figures_experiments/reports/token_analysis/14b_facts_first_all_hop_refined/refined_by_stage_mean_signed.png}
        \caption{\textit{Qwen3-14B}, \texttt{facts\_first}}
    \end{subfigure}

    \vspace{0.2em}

    \begin{subfigure}[b]{0.44\textwidth}
        \centering
        \includegraphics[width=\linewidth]{./figures_experiments/reports/token_analysis/8b_expr_first_all_hop_refined/refined_by_stage_mean_signed.png}
        \caption{\textit{Qwen3-8B}, \texttt{expr\_first}}
    \end{subfigure}
    \hfill
    \begin{subfigure}[b]{0.44\textwidth}
        \centering
        \includegraphics[width=\linewidth]{./figures_experiments/reports/token_analysis/14b_expr_first_all_hop_refined/refined_by_stage_mean_signed.png}
        \caption{\textit{Qwen3-14B}, \texttt{expr\_first}}
    \end{subfigure}
    \caption{\textbf{Qwen3 refined signed-mean token diagnostics.} The signed mean-per-token view confirms that the facts-versus-query balance is robust to both normalization choices. Under \texttt{expr\_first}, the facts category remains directionally dominant in 14B even though query-facing mass is more visible in the unsigned plots, indicating that prompt reordering broadens computation without eliminating fact access.}
    \label{fig:qwen_token_refined_mean_signed}
\end{figure*}

%\begin{figure*}[p]
%    \centering
%    \begin{subfigure}[b]{0.44\textwidth}
%        \centering
%        \includegraphics[width=\linewidth]{./figures_experiments/reports_llama3/token_analysis_facts_first_raw/heatmap_simple.png}
%        \caption{\textit{Llama-3.1-8B-Instruct}, \texttt{facts\_first}}
%    \end{subfigure}
%    \hfill
%    \begin{subfigure}[b]{0.44\textwidth}
%        \centering
%        \includegraphics[width=\linewidth]{./figures_experiments/reports_llama3/token_analysis_expr_first_raw/heatmap_simple.png}
%        \caption{\textit{Llama-3.1-8B-Instruct}, \texttt{expr\_first}}
%    \end{subfigure}
%
%    \vspace{0.2em}
%
%    \begin{subfigure}[b]{0.44\textwidth}
%        \centering
%        \includegraphics[width=\linewidth]{./figures_experiments/reports_mistral/token_analysis_facts_first_raw/heatmap_simple.png}
%        \caption{\textit{Mistral-7B-Instruct-v0.1}, \texttt{facts\_first}}
%    \end{subfigure}
%    \hfill
%    \begin{subfigure}[b]{0.44\textwidth}
%        \centering
%        \includegraphics[width=\linewidth]{./figures_experiments/reports_mistral/token_analysis_expr_first_raw/heatmap_simple.png}
%        \caption{\textit{Mistral-7B-Instruct-v0.1}, \texttt{expr\_first}}
%    \end{subfigure}
%    \caption{\textbf{Cross-model raw token heatmaps.} Before any category refinement, both non-Qwen families already show heavier query-facing structure than Qwen3. Llama is especially diffuse, with broad late-layer activity over the query-facing portion of the prompt, while Mistral keeps a somewhat cleaner early-to-late progression but still concentrates less sharply on fact tokens than Qwen3.}
%    \label{fig:cross_model_token_raw_heatmaps}
%\end{figure*}

\begin{figure*}[t]
    \centering
    \begin{subfigure}[b]{0.44\textwidth}
        \centering
        \includegraphics[width=\linewidth]{./figures_experiments/reports_llama3/token_analysis_facts_first_refined/refined_by_stage_mean.png}
        \caption{\textit{Llama-3.1-8B-Instruct}, \texttt{facts\_first}}
    \end{subfigure}
    \hfill
    \begin{subfigure}[b]{0.44\textwidth}
        \centering
        \includegraphics[width=\linewidth]{./figures_experiments/reports_llama3/token_analysis_expr_first_refined/refined_by_stage_mean.png}
        \caption{\textit{Llama-3.1-8B-Instruct}, \texttt{expr\_first}}
    \end{subfigure}

    \vspace{0.2em}

    \begin{subfigure}[b]{0.44\textwidth}
        \centering
        \includegraphics[width=\linewidth]{./figures_experiments/reports_mistral/token_analysis_facts_first_refined/refined_by_stage_mean.png}
        \caption{\textit{Mistral-7B-Instruct-v0.1}, \texttt{facts\_first}}
    \end{subfigure}
    \hfill
    \begin{subfigure}[b]{0.44\textwidth}
        \centering
        \includegraphics[width=\linewidth]{./figures_experiments/reports_mistral/token_analysis_expr_first_refined/refined_by_stage_mean.png}
        \caption{\textit{Mistral-7B-Instruct-v0.1}, \texttt{expr\_first}}
    \end{subfigure}
    \caption{\textbf{Cross-model refined mean token diagnostics.} In both non-Qwen families, \texttt{query\_token} remains the largest mean contribution under both prompt orders. This is a compact quantitative diagnostic of the more output-position-centered strategy already visible in the band-composition summaries.}
    \label{fig:cross_model_token_refined_mean}
\end{figure*}

\begin{figure*}[t]
    \centering
    \begin{subfigure}[b]{0.44\textwidth}
        \centering
        \includegraphics[width=\linewidth]{./figures_experiments/reports_llama3/token_analysis_facts_first_refined/refined_by_stage_sum_signed.png}
        \caption{\textit{Llama-3.1-8B-Instruct}, \texttt{facts\_first}}
    \end{subfigure}
    \hfill
    \begin{subfigure}[b]{0.44\textwidth}
        \centering
        \includegraphics[width=\linewidth]{./figures_experiments/reports_llama3/token_analysis_expr_first_refined/refined_by_stage_sum_signed.png}
        \caption{\textit{Llama-3.1-8B-Instruct}, \texttt{expr\_first}}
    \end{subfigure}

    \vspace{0.2em}

    \begin{subfigure}[b]{0.44\textwidth}
        \centering
        \includegraphics[width=\linewidth]{./figures_experiments/reports_mistral/token_analysis_facts_first_refined/refined_by_stage_sum_signed.png}
        \caption{\textit{Mistral-7B-Instruct-v0.1}, \texttt{facts\_first}}
    \end{subfigure}
    \hfill
    \begin{subfigure}[b]{0.44\textwidth}
        \centering
        \includegraphics[width=\linewidth]{./figures_experiments/reports_mistral/token_analysis_expr_first_refined/refined_by_stage_sum_signed.png}
        \caption{\textit{Mistral-7B-Instruct-v0.1}, \texttt{expr\_first}}
    \end{subfigure}
    \caption{\textbf{Cross-model refined signed-sum token diagnostics.} Retaining the sign sharpens the inter-family contrast, but not uniformly. Llama remains clearly query-dominant, Qwen3 remains clearly fact-dominant, and Mistral falls between the two: it is nearly balanced under \texttt{facts\_first} and becomes more query-heavy under \texttt{expr\_first}. This is why Mistral looks more query-facing than Qwen3 without being as one-sided as Llama in every signed view.}
    \label{fig:cross_model_token_refined_sum_signed}
\end{figure*}

\begin{figure*}[t]
    \centering
    \begin{subfigure}[b]{0.44\textwidth}
        \centering
        \includegraphics[width=\linewidth]{./figures_experiments/reports_llama3/token_analysis_facts_first_refined/refined_by_stage_mean_signed.png}
        \caption{\textit{Llama-3.1-8B-Instruct}, \texttt{facts\_first}}
    \end{subfigure}
    \hfill
    \begin{subfigure}[b]{0.44\textwidth}
        \centering
        \includegraphics[width=\linewidth]{./figures_experiments/reports_llama3/token_analysis_expr_first_refined/refined_by_stage_mean_signed.png}
        \caption{\textit{Llama-3.1-8B-Instruct}, \texttt{expr\_first}}
    \end{subfigure}

    \vspace{0.2em}

    \begin{subfigure}[b]{0.44\textwidth}
        \centering
        \includegraphics[width=\linewidth]{./figures_experiments/reports_mistral/token_analysis_facts_first_refined/refined_by_stage_mean_signed.png}
        \caption{\textit{Mistral-7B-Instruct-v0.1}, \texttt{facts\_first}}
    \end{subfigure}
    \hfill
    \begin{subfigure}[b]{0.44\textwidth}
        \centering
        \includegraphics[width=\linewidth]{./figures_experiments/reports_mistral/token_analysis_expr_first_refined/refined_by_stage_mean_signed.png}
        \caption{\textit{Mistral-7B-Instruct-v0.1}, \texttt{expr\_first}}
    \end{subfigure}
    \caption{\textbf{Cross-model refined signed-mean token diagnostics.} The per-token normalization leaves the same ordering intact: Llama and Mistral remain query-heavy even after accounting for category size. This confirms that the cross-model difference is not driven by having more query-region tokens in the prompt template, but by where the models actually place their causal sensitivity.}
    \label{fig:cross_model_token_refined_mean_signed}
\end{figure*}

\subsection{Additional Head Diagnostics}
\label{sec:appendix_extra_heads_gallery}

The omitted head-analysis figures extend the taxonomy line charts already shown in the appendix. The Step~1 discovery maps show whether candidate heads are spatially clustered before any role label is assigned. The Step~3 plots then expose both the magnitude and direction of the validation effect through absolute-ratio and raw dPD-shift curves. For Qwen3, these diagnostics align closely: candidate heads cluster in the same early/middle/late bands later recovered by the taxonomy, and role-targeted curves remain much more damaging than \textit{Random-Outside}. For Llama and Mistral, the same views show weaker separation and larger side effects from the random control.

\paragraph{How to read these head figures.}
The discovery maps should be read as \emph{candidate density over layers and heads}: clustering means that promising heads are already spatially organized before the taxonomy rules are applied. The absolute-ratio validation curves answer ``how damaging is this ablation relative to the model's original margin?'', so higher curves are worse. The raw dPD-shift curves add direction: more negative values mean the ablation pushes the model farther away from the correct answer, while values near zero indicate weak causal effect.

The Qwen3-specific head extensions appear in Figures~\ref{fig:qwen_heads_step1_gallery} and~\ref{fig:qwen_heads_dpd_shift_gallery}, while the corresponding cross-model diagnostics appear in Figures~\ref{fig:cross_model_heads_step1_gallery} and~\ref{fig:cross_model_heads_dpd_shift_gallery}.

\begin{figure*}[t]
    \centering
    \begin{subfigure}[b]{0.44\textwidth}
        \centering
        \includegraphics[width=\linewidth]{./figures_experiments/reports/heads_analysis_qwen3_8b_facts_first/classify/layer_head_role_distribution.png}
        \caption{\textit{Qwen3-8B}, \texttt{facts\_first}}
    \end{subfigure}
    \hfill
    \begin{subfigure}[b]{0.44\textwidth}
        \centering
        \includegraphics[width=\linewidth]{./figures_experiments/reports/heads_analysis_qwen3_8b_expr_first/classify/layer_head_role_distribution.png}
        \caption{\textit{Qwen3-8B}, \texttt{expr\_first}}
    \end{subfigure}

    \vspace{0.2em}

    \begin{subfigure}[b]{0.44\textwidth}
        \centering
        \includegraphics[width=\linewidth]{./figures_experiments/reports/heads_analysis_qwen3_14b_facts_first/classify/layer_head_role_distribution.png}
        \caption{\textit{Qwen3-14B}, \texttt{facts\_first}}
    \end{subfigure}
    \hfill
    \begin{subfigure}[b]{0.44\textwidth}
        \centering
        \includegraphics[width=\linewidth]{./figures_experiments/reports/heads_analysis_qwen3_14b_expr_first/classify/layer_head_role_distribution.png}
        \caption{\textit{Qwen3-14B}, \texttt{expr\_first}}
    \end{subfigure}
    \caption{\textbf{Qwen3 Step~1 discovery maps across scales and prompt orders.} Even before taxonomy assignment, the candidate heads are not uniformly spread across the network. Early-layer candidates cluster near boundary-sensitive regions, middle-layer candidates align with transmission bands, and later candidates shift toward fact-retrieval behavior. The \texttt{expr\_first} intervention increases the density of early-layer candidates, matching the higher abundance of Splitting Heads reported in Step~2.}
    \label{fig:qwen_heads_step1_gallery}
\end{figure*}

\begin{figure*}[t]
    \centering
    \begin{subfigure}[b]{0.44\textwidth}
        \centering
        \includegraphics[width=\linewidth]{./figures_experiments/reports/heads_analysis_qwen3_8b_facts_first/validation/plots/dpd_shift_curve.png}
        \caption{\textit{Qwen3-8B}, \texttt{facts\_first}}
    \end{subfigure}
    \hfill
    \begin{subfigure}[b]{0.44\textwidth}
        \centering
        \includegraphics[width=\linewidth]{./figures_experiments/reports/heads_analysis_qwen3_8b_expr_first/validation/plots/dpd_shift_curve.png}
        \caption{\textit{Qwen3-8B}, \texttt{expr\_first}}
    \end{subfigure}

    \vspace{0.2em}

    \begin{subfigure}[b]{0.44\textwidth}
        \centering
        \includegraphics[width=\linewidth]{./figures_experiments/reports/heads_analysis_qwen3_14b_facts_first/validation/plots/dpd_shift_curve.png}
        \caption{\textit{Qwen3-14B}, \texttt{facts\_first}}
    \end{subfigure}
    \hfill
    \begin{subfigure}[b]{0.44\textwidth}
        \centering
        \includegraphics[width=\linewidth]{./figures_experiments/reports/heads_analysis_qwen3_14b_expr_first/validation/plots/dpd_shift_curve.png}
        \caption{\textit{Qwen3-14B}, \texttt{expr\_first}}
    \end{subfigure}
    \caption{\textbf{Qwen3 Step~3 dPD-shift diagnostics.} The raw probability-difference shifts are again strongly negative for the discovered roles and remain close to zero for the random baseline. The effect is strongest in \textit{Qwen3-8B} \texttt{facts\_first}, where Fact-Retrieval heads cause the largest downward shift, but the same ordering remains visible in the other three settings.}
    \label{fig:qwen_heads_dpd_shift_gallery}
\end{figure*}

\begin{figure*}[t]
    \centering
    \begin{subfigure}[b]{0.44\textwidth}
        \centering
        \includegraphics[width=\linewidth]{./figures_experiments/reports_llama3/heads_analysis_facts_first/classify/layer_head_role_distribution.png}
        \caption{Llama, \texttt{facts\_first}}
    \end{subfigure}
    \hfill
    \begin{subfigure}[b]{0.44\textwidth}
        \centering
        \includegraphics[width=\linewidth]{./figures_experiments/reports_llama3/heads_analysis_expr_first/classify/layer_head_role_distribution.png}
        \caption{Llama, \texttt{expr\_first}}
    \end{subfigure}

    \vspace{0.2em}

    \begin{subfigure}[b]{0.44\textwidth}
        \centering
        \includegraphics[width=\linewidth]{./figures_experiments/reports_mistral/heads_analysis_facts_first/classify/layer_head_role_distribution.png}
        \caption{Mistral, \texttt{facts\_first}}
    \end{subfigure}
    \hfill
    \begin{subfigure}[b]{0.44\textwidth}
        \centering
        \includegraphics[width=\linewidth]{./figures_experiments/reports_mistral/heads_analysis_expr_first/classify/layer_head_role_distribution.png}
        \caption{Mistral, \texttt{expr\_first}}
    \end{subfigure}
    \caption{\textbf{Cross-model Step~1 discovery maps.} The same coarse early/middle/late clustering is still visible in Llama and Mistral, but the candidate sets are visibly more diffuse than in Qwen3. This foreshadows the weaker role separation later visible in the absolute-ratio and dPD-shift validation curves.}
    \label{fig:cross_model_heads_step1_gallery}
\end{figure*}

\begin{figure*}[t]
    \centering
    \begin{subfigure}[b]{0.44\textwidth}
        \centering
        \includegraphics[width=\linewidth]{./figures_experiments/reports_llama3/heads_analysis_facts_first/validation/plots/dpd_shift_curve.png}
        \caption{Llama, \texttt{facts\_first}}
    \end{subfigure}
    \hfill
    \begin{subfigure}[b]{0.44\textwidth}
        \centering
        \includegraphics[width=\linewidth]{./figures_experiments/reports_llama3/heads_analysis_expr_first/validation/plots/dpd_shift_curve.png}
        \caption{Llama, \texttt{expr\_first}}
    \end{subfigure}

    \vspace{0.2em}

    \begin{subfigure}[b]{0.44\textwidth}
        \centering
        \includegraphics[width=\linewidth]{./figures_experiments/reports_mistral/heads_analysis_facts_first/validation/plots/dpd_shift_curve.png}
        \caption{Mistral, \texttt{facts\_first}}
    \end{subfigure}
    \hfill
    \begin{subfigure}[b]{0.44\textwidth}
        \centering
        \includegraphics[width=\linewidth]{./figures_experiments/reports_mistral/heads_analysis_expr_first/validation/plots/dpd_shift_curve.png}
        \caption{Mistral, \texttt{expr\_first}}
    \end{subfigure}
    \caption{\textbf{Cross-model Step~3 dPD-shift diagnostics.} The raw dPD shifts confirm the same story. Llama's targeted interventions are harmful but shallow, and the random baseline can even move in the helpful direction. Mistral's targeted interventions remain sharply harmful, yet the random baseline also produces a noticeable positive displacement, consistent with a highly sensitive but less cleanly isolated head substrate.}
    \label{fig:cross_model_heads_dpd_shift_gallery}
\end{figure*}

\subsection{Further Region and Magnitude Atlases}
\label{sec:appendix_further_atlas}

This subsection collects the remaining retained region galleries. They are included to complete the omitted hop- and branch-specific views, especially where the earlier appendix sections only gave a partial picture.

\paragraph{Qwen3 facts-first supplements.}
The remaining Qwen3 figures fill three specific gaps: the omitted two-hop attention gallery, the omitted two-hop joint attention+MLP gallery, and the retained facts-first two-hop MLP gallery. Together they show that increased reasoning depth broadens the Query Region without changing the overall Facts $\rightarrow$ Expression $\rightarrow$ Query order.

\begin{figure*}[t]
    \centering
    \begin{subfigure}[b]{0.285\textwidth}
        \centering
        \includegraphics[width=\linewidth]{./figures_experiments/reports/attn_analysis/attn_8b_two_hop/facts_region/attn_facts_region_scores.png}
        \caption{\textit{Qwen3-8B}: Facts}
    \end{subfigure}
    \hfill
    \begin{subfigure}[b]{0.285\textwidth}
        \centering
        \includegraphics[width=\linewidth]{./figures_experiments/reports/attn_analysis/attn_8b_two_hop/expression_region/attn_expression_region_scores.png}
        \caption{\textit{Qwen3-8B}: Expression}
    \end{subfigure}
    \hfill
    \begin{subfigure}[b]{0.285\textwidth}
        \centering
        \includegraphics[width=\linewidth]{./figures_experiments/reports/attn_analysis/attn_8b_two_hop/query_region/attn_query_region_scores.png}
        \caption{\textit{Qwen3-8B}: Query}
    \end{subfigure}

    \vspace{0.2em}

    \begin{subfigure}[b]{0.285\textwidth}
        \centering
        \includegraphics[width=\linewidth]{./figures_experiments/reports/attn_analysis/attn_14b_two_hop/facts_region/attn_facts_region_scores.png}
        \caption{\textit{Qwen3-14B}: Facts}
    \end{subfigure}
    \hfill
    \begin{subfigure}[b]{0.285\textwidth}
        \centering
        \includegraphics[width=\linewidth]{./figures_experiments/reports/attn_analysis/attn_14b_two_hop/expression_region/attn_expression_region_scores.png}
        \caption{\textit{Qwen3-14B}: Expression}
    \end{subfigure}
    \hfill
    \begin{subfigure}[b]{0.285\textwidth}
        \centering
        \includegraphics[width=\linewidth]{./figures_experiments/reports/attn_analysis/attn_14b_two_hop/query_region/attn_query_region_scores.png}
        \caption{\textit{Qwen3-14B}: Query}
    \end{subfigure}
    \caption{\textbf{Qwen3 facts-first attention-region galleries for two-hop reasoning.} Two-hop attention patching preserves the same semantic order as one-hop, but widens the query-facing integration window. In both model sizes, the Facts Region remains early-dominant, the Expression Region splits between early and middle layers, and the Query Region acquires its largest mass in the middle-to-late bands.}
    \label{fig:qwen_attn_twohop_facts_gallery}
\end{figure*}

\begin{figure*}[t]
    \centering
    \begin{subfigure}[b]{0.285\textwidth}
        \centering
        \includegraphics[width=\linewidth]{./figures_experiments/reports/attn_mlp_analysis/attn_mlp_8b_two_hop/facts_region/attn_mlp_facts_region_scores.png}
        \caption{\textit{Qwen3-8B}: Facts}
    \end{subfigure}
    \hfill
    \begin{subfigure}[b]{0.285\textwidth}
        \centering
        \includegraphics[width=\linewidth]{./figures_experiments/reports/attn_mlp_analysis/attn_mlp_8b_two_hop/expression_region/attn_mlp_expression_region_scores.png}
        \caption{\textit{Qwen3-8B}: Expression}
    \end{subfigure}
    \hfill
    \begin{subfigure}[b]{0.285\textwidth}
        \centering
        \includegraphics[width=\linewidth]{./figures_experiments/reports/attn_mlp_analysis/attn_mlp_8b_two_hop/query_region/attn_mlp_query_region_scores.png}
        \caption{\textit{Qwen3-8B}: Query}
    \end{subfigure}

    \vspace{0.2em}

    \begin{subfigure}[b]{0.285\textwidth}
        \centering
        \includegraphics[width=\linewidth]{./figures_experiments/reports/attn_mlp_analysis/attn_mlp_14b_two_hop/facts_region/attn_mlp_facts_region_scores.png}
        \caption{\textit{Qwen3-14B}: Facts}
    \end{subfigure}
    \hfill
    \begin{subfigure}[b]{0.285\textwidth}
        \centering
        \includegraphics[width=\linewidth]{./figures_experiments/reports/attn_mlp_analysis/attn_mlp_14b_two_hop/expression_region/attn_mlp_expression_region_scores.png}
        \caption{\textit{Qwen3-14B}: Expression}
    \end{subfigure}
    \hfill
    \begin{subfigure}[b]{0.285\textwidth}
        \centering
        \includegraphics[width=\linewidth]{./figures_experiments/reports/attn_mlp_analysis/attn_mlp_14b_two_hop/query_region/attn_mlp_query_region_scores.png}
        \caption{\textit{Qwen3-14B}: Query}
    \end{subfigure}
    \caption{\textbf{Qwen3 facts-first joint attention+MLP galleries for two-hop reasoning.} The joint branch remains closer to the original staged decomposition than pure attention. Facts keep the strongest early-layer dominance, Expression spans the early-to-middle transition, and Query widens mainly in middle and late layers. The same ordering appears in both scales, which supports the view that routing and nonlinear transformation cooperate rather than compete in the joint intervention.}
    \label{fig:qwen_attn_mlp_twohop_facts_gallery}
\end{figure*}

%\begin{figure*}[p]
%    \centering
%    \begin{subfigure}[b]{0.29\textwidth}
%        \centering
%        \includegraphics[width=\linewidth]{./figures_experiments/reports/mlp_analysis/comparison_14B/band_metrics_panel.png}
%        \caption{MLP, 14B \texttt{facts\_first}}
%    \end{subfigure}
%    \hfill
%    \begin{subfigure}[b]{0.29\textwidth}
%        \centering
%        \includegraphics[width=\linewidth]{./figures_experiments/reports/attn_analysis/comparison_14B/band_metrics_panel.png}
%        \caption{Attention, 14B \texttt{facts\_first}}
%    \end{subfigure}
%    \hfill
%    \begin{subfigure}[b]{0.29\textwidth}
%        \centering
%        \includegraphics[width=\linewidth]{./figures_experiments/reports/attn_mlp_analysis/comparison_14B/band_metrics_panel.png}
%        \caption{Attention+MLP, 14B \texttt{facts\_first}}
%    \end{subfigure}
%
%    \vspace{0.2em}
%
%    \begin{subfigure}[b]{0.44\textwidth}
%        \centering
%        \includegraphics[width=\linewidth]{./figures_experiments/reports/attn_analysis_expr_first/comparison_expr_first_14B/band_metrics_panel.png}
%        \caption{Attention, 14B \texttt{expr\_first}}
%    \end{subfigure}
%    \hfill
%    \begin{subfigure}[b]{0.44\textwidth}
%        \centering
%        \includegraphics[width=\linewidth]{./figures_experiments/reports/attn_mlp_analysis_expr_first/comparison_expr_first_14B/band_metrics_panel.png}
%        \caption{Attention+MLP, 14B \texttt{expr\_first}}
%    \end{subfigure}
%    \caption{\textbf{Remaining Qwen3-14B branch-level band summaries.} These panels complete the set of 14B branch summaries across prompt orders. They show the same main tendency already established elsewhere: \texttt{expr\_first} moves more mass into early Expression processing, while the attention branch remains the most middle-heavy and the joint branch the most stage-like.}
%    \label{fig:qwen14_remaining_band_panels}
%\end{figure*}

%\begin{figure*}[p]
%    \centering
%    \begin{subfigure}[b]{0.44\textwidth}
%        \centering
%        \includegraphics[width=\linewidth]{./figures_experiments/reports/attn_analysis_expr_first/comparison_expr_first/band_metrics_panel.png}
%        \caption{Attention, 8B \texttt{expr\_first}}
%    \end{subfigure}
%    \hfill
%    \begin{subfigure}[b]{0.44\textwidth}
%        \centering
%        \includegraphics[width=\linewidth]{./figures_experiments/reports/attn_mlp_analysis_expr_first/comparison_expr_first/band_metrics_panel.png}
%        \caption{Attention+MLP, 8B \texttt{expr\_first}}
%    \end{subfigure}
%    \caption{\textbf{Remaining Qwen3-8B \texttt{expr\_first} branch summaries.} These two omitted panels complete the 8B prompt-order intervention picture for the non-MLP branches. They again show that attention makes the prompt-order shift more structural, while the joint branch retains a cleaner staged decomposition.}
%    \label{fig:qwen8_expr_remaining_band_panels}
%\end{figure*}

%\begin{figure*}[t]
%    \centering
%    \begin{subfigure}[b]{0.44\textwidth}
%        \centering
%        \includegraphics[width=\linewidth]{./figures_experiments/reports/heads_analysis_qwen3_8b_facts_first/validation/plots/pd_abs_ratio_curve.png}
%        \caption{\textit{Qwen3-8B}, \texttt{facts\_first}}
%    \end{subfigure}
%    \hfill
%    \begin{subfigure}[b]{0.44\textwidth}
%        \centering
%        \includegraphics[width=\linewidth]{./figures_experiments/reports/heads_analysis_qwen3_14b_facts_first/validation/plots/pd_abs_ratio_curve.png}
%        \caption{\textit{Qwen3-14B}, \texttt{facts\_first}}
%    \end{subfigure}
%    \caption{\textbf{Qwen3 facts-first absolute validation curves.} These are the two remaining absolute-ratio head-validation plots not previously shown. They complete the Qwen3 Step~3 atlas by pairing the conventional absolute-ratio view with the dPD-shift diagnostics under \texttt{facts\_first}. Both scales preserve the same ordering: Fact-Retrieval is the most damaging single-role ablation, Transmission is second, Splitting is weakest, and \textit{Random-Outside} remains far less harmful.}
%    \label{fig:qwen_heads_abs_ratio_facts}
%\end{figure*}

\begin{figure*}[t]
    \centering
    \begin{subfigure}[b]{0.285\textwidth}
        \centering
        \includegraphics[width=\linewidth]{./figures_experiments/reports/mlp_analysis/mlp_8b_two_hop/facts_region/mlp_facts_region_scores.png}
        \caption{Facts}
    \end{subfigure}
    \hfill
    \begin{subfigure}[b]{0.285\textwidth}
        \centering
        \includegraphics[width=\linewidth]{./figures_experiments/reports/mlp_analysis/mlp_8b_two_hop/expression_region/mlp_expression_region_scores.png}
        \caption{Expression}
    \end{subfigure}
    \hfill
    \begin{subfigure}[b]{0.285\textwidth}
        \centering
        \includegraphics[width=\linewidth]{./figures_experiments/reports/mlp_analysis/mlp_8b_two_hop/query_region/mlp_query_region_scores.png}
        \caption{Query}
    \end{subfigure}
    \caption{\textbf{Qwen3-8B facts-first MLP gallery for two-hop reasoning.} These three curves complete the 8B facts-first MLP atlas. Compared with one-hop, two-hop reasoning weakens the sharpness of every region but preserves the same order: facts dominate earliest, expression remains concentrated in the middle of the network, and query stays strongest in the middle-to-late transition.}
    \label{fig:qwen8_mlp_twohop_facts_gallery}
\end{figure*}

%\begin{figure*}[p]
%    \centering
%    \begin{subfigure}[b]{0.285\textwidth}
%        \centering
%        \includegraphics[width=\linewidth]{./figures_experiments/reports/mlp_analysis/mlp_8b_one_hop/facts_region/mlp_facts_region_abs_dpd_scores.png}
%        \caption{Facts (One-hop)}
%    \end{subfigure}
%    \hfill
%    \begin{subfigure}[b]{0.285\textwidth}
%        \centering
%        \includegraphics[width=\linewidth]{./figures_experiments/reports/mlp_analysis/mlp_8b_one_hop/expression_region/mlp_expression_region_abs_dpd_scores.png}
%        \caption{Expression (One-hop)}
%    \end{subfigure}
%    \hfill
%    \begin{subfigure}[b]{0.285\textwidth}
%        \centering
%        \includegraphics[width=\linewidth]{./figures_experiments/reports/mlp_analysis/mlp_8b_one_hop/query_region/mlp_query_region_abs_dpd_scores.png}
%        \caption{Query (One-hop)}
%    \end{subfigure}
%
%    \vspace{0.2em}
%
%    \begin{subfigure}[b]{0.285\textwidth}
%        \centering
%        \includegraphics[width=\linewidth]{./figures_experiments/reports/mlp_analysis/mlp_8b_two_hop/facts_region/mlp_facts_region_abs_dpd_scores.png}
%        \caption{Facts (Two-hop)}
%    \end{subfigure}
%    \hfill
%    \begin{subfigure}[b]{0.285\textwidth}
%        \centering
%        \includegraphics[width=\linewidth]{./figures_experiments/reports/mlp_analysis/mlp_8b_two_hop/expression_region/mlp_expression_region_abs_dpd_scores.png}
%        \caption{Expression (Two-hop)}
%    \end{subfigure}
%    \hfill
%    \begin{subfigure}[b]{0.285\textwidth}
%        \centering
%        \includegraphics[width=\linewidth]{./figures_experiments/reports/mlp_analysis/mlp_8b_two_hop/query_region/mlp_query_region_abs_dpd_scores.png}
%        \caption{Query (Two-hop)}
%    \end{subfigure}
%    \caption{\textbf{Qwen3-8B facts-first absolute-dPD MLP gallery.} Using $|\mathrm{dPD}|$ instead of signed dPD shows that the same regions dominate in pure magnitude. The two-hop curves are less sharp than one-hop, but the largest-magnitude intervention sites still follow the same facts-to-expression-to-query ordering.}
%    \label{fig:qwen8_mlp_abs_facts_gallery}
%\end{figure*}

%\begin{figure*}[p]
%    \centering
%    \begin{subfigure}[b]{0.285\textwidth}
%        \centering
%        \includegraphics[width=\linewidth]{./figures_experiments/reports/mlp_analysis/mlp_14b_one_hop/facts_region/mlp_facts_region_abs_dpd_scores.png}
%        \caption{Facts (One-hop)}
%    \end{subfigure}
%    \hfill
%    \begin{subfigure}[b]{0.285\textwidth}
%        \centering
%        \includegraphics[width=\linewidth]{./figures_experiments/reports/mlp_analysis/mlp_14b_one_hop/expression_region/mlp_expression_region_abs_dpd_scores.png}
%        \caption{Expression (One-hop)}
%    \end{subfigure}
%    \hfill
%    \begin{subfigure}[b]{0.285\textwidth}
%        \centering
%        \includegraphics[width=\linewidth]{./figures_experiments/reports/mlp_analysis/mlp_14b_one_hop/query_region/mlp_query_region_abs_dpd_scores.png}
%        \caption{Query (One-hop)}
%    \end{subfigure}
%
%    \vspace{0.2em}
%
%    \begin{subfigure}[b]{0.285\textwidth}
%        \centering
%        \includegraphics[width=\linewidth]{./figures_experiments/reports/mlp_analysis/mlp_14b_two_hop/facts_region/mlp_facts_region_abs_dpd_scores.png}
%        \caption{Facts (Two-hop)}
%    \end{subfigure}
%    \hfill
%    \begin{subfigure}[b]{0.285\textwidth}
%        \centering
%        \includegraphics[width=\linewidth]{./figures_experiments/reports/mlp_analysis/mlp_14b_two_hop/expression_region/mlp_expression_region_abs_dpd_scores.png}
%        \caption{Expression (Two-hop)}
%    \end{subfigure}
%    \hfill
%    \begin{subfigure}[b]{0.285\textwidth}
%        \centering
%        \includegraphics[width=\linewidth]{./figures_experiments/reports/mlp_analysis/mlp_14b_two_hop/query_region/mlp_query_region_abs_dpd_scores.png}
%        \caption{Query (Two-hop)}
%    \end{subfigure}
%    \caption{\textbf{Qwen3-14B facts-first absolute-dPD MLP gallery.} The larger model preserves the same magnitude profile as 8B while smoothing some of the late-layer noise. Facts remain the largest early-layer target, Expression retains the clearest middle-layer peak, and Query broadens into the middle-to-late integration stage.}
%    \label{fig:qwen14_mlp_abs_facts_gallery}
%\end{figure*}

%\begin{figure*}[p]
%    \centering
%    \begin{subfigure}[b]{0.285\textwidth}
%        \centering
%        \includegraphics[width=\linewidth]{./figures_experiments/reports/mlp_analysis_expr_first/mlp_8b_one_hop_expr_first/facts_region/mlp_facts_region_abs_dpd_scores.png}
%        \caption{Facts (One-hop)}
%    \end{subfigure}
%    \hfill
%    \begin{subfigure}[b]{0.285\textwidth}
%        \centering
%        \includegraphics[width=\linewidth]{./figures_experiments/reports/mlp_analysis_expr_first/mlp_8b_one_hop_expr_first/expression_region/mlp_expression_region_abs_dpd_scores.png}
%        \caption{Expression (One-hop)}
%    \end{subfigure}
%    \hfill
%    \begin{subfigure}[b]{0.285\textwidth}
%        \centering
%        \includegraphics[width=\linewidth]{./figures_experiments/reports/mlp_analysis_expr_first/mlp_8b_one_hop_expr_first/query_region/mlp_query_region_abs_dpd_scores.png}
%        \caption{Query (One-hop)}
%    \end{subfigure}
%
%    \vspace{0.2em}
%
%    \begin{subfigure}[b]{0.285\textwidth}
%        \centering
%        \includegraphics[width=\linewidth]{./figures_experiments/reports/mlp_analysis_expr_first/mlp_8b_two_hop_expr_first/facts_region/mlp_facts_region_abs_dpd_scores.png}
%        \caption{Facts (Two-hop)}
%    \end{subfigure}
%    \hfill
%    \begin{subfigure}[b]{0.285\textwidth}
%        \centering
%        \includegraphics[width=\linewidth]{./figures_experiments/reports/mlp_analysis_expr_first/mlp_8b_two_hop_expr_first/expression_region/mlp_expression_region_abs_dpd_scores.png}
%        \caption{Expression (Two-hop)}
%    \end{subfigure}
%    \hfill
%    \begin{subfigure}[b]{0.285\textwidth}
%        \centering
%        \includegraphics[width=\linewidth]{./figures_experiments/reports/mlp_analysis_expr_first/mlp_8b_two_hop_expr_first/query_region/mlp_query_region_abs_dpd_scores.png}
%        \caption{Query (Two-hop)}
%    \end{subfigure}
%    \caption{\textbf{Qwen3-8B expr-first absolute-dPD MLP gallery.} In magnitude space, prompt reordering still pulls Expression earlier and broadens Query. The absolute plots show that the prompt-order shift is not only a signed-direction effect: the same front-loaded expression bias is visible even after removing sign.}
%    \label{fig:qwen8_mlp_abs_expr_gallery}
%\end{figure*}

%\begin{figure*}[p]
%    \centering
%    \begin{subfigure}[b]{0.285\textwidth}
%        \centering
%        \includegraphics[width=\linewidth]{./figures_experiments/reports/mlp_analysis_expr_first/mlp_14b_one_hop_expr_first/facts_region/mlp_facts_region_abs_dpd_scores.png}
%        \caption{Facts (One-hop)}
%    \end{subfigure}
%    \hfill
%    \begin{subfigure}[b]{0.285\textwidth}
%        \centering
%        \includegraphics[width=\linewidth]{./figures_experiments/reports/mlp_analysis_expr_first/mlp_14b_one_hop_expr_first/expression_region/mlp_expression_region_abs_dpd_scores.png}
%        \caption{Expression (One-hop)}
%    \end{subfigure}
%    \hfill
%    \begin{subfigure}[b]{0.285\textwidth}
%        \centering
%        \includegraphics[width=\linewidth]{./figures_experiments/reports/mlp_analysis_expr_first/mlp_14b_one_hop_expr_first/query_region/mlp_query_region_abs_dpd_scores.png}
%        \caption{Query (One-hop)}
%    \end{subfigure}
%
%    \vspace{0.2em}
%
%    \begin{subfigure}[b]{0.285\textwidth}
%        \centering
%        \includegraphics[width=\linewidth]{./figures_experiments/reports/mlp_analysis_expr_first/mlp_14b_two_hop_expr_first/facts_region/mlp_facts_region_abs_dpd_scores.png}
%        \caption{Facts (Two-hop)}
%    \end{subfigure}
%    \hfill
%    \begin{subfigure}[b]{0.285\textwidth}
%        \centering
%        \includegraphics[width=\linewidth]{./figures_experiments/reports/mlp_analysis_expr_first/mlp_14b_two_hop_expr_first/expression_region/mlp_expression_region_abs_dpd_scores.png}
%        \caption{Expression (Two-hop)}
%    \end{subfigure}
%    \hfill
%    \begin{subfigure}[b]{0.285\textwidth}
%        \centering
%        \includegraphics[width=\linewidth]{./figures_experiments/reports/mlp_analysis_expr_first/mlp_14b_two_hop_expr_first/query_region/mlp_query_region_abs_dpd_scores.png}
%        \caption{Query (Two-hop)}
%    \end{subfigure}
%    \caption{\textbf{Qwen3-14B expr-first absolute-dPD MLP gallery.} The 14B absolute plots are the cleanest version of the prompt-order intervention at region level. Expression is the dominant early magnitude peak, Facts broaden into the early-to-middle transition, and Query remains a distributed middle-to-late sink.}
%    \label{fig:qwen14_mlp_abs_expr_gallery}
%\end{figure*}

%\begin{figure*}[p]
%    \centering
%    \begin{subfigure}[b]{0.44\textwidth}
%        \centering
%        \includegraphics[width=\linewidth]{./figures_experiments/reports/token_analysis/8b_facts_first_all_hop_raw/category_comparison_simple.png}
%        \caption{\textit{Qwen3-8B}, \texttt{facts\_first}}
%    \end{subfigure}
%    \hfill
%    \begin{subfigure}[b]{0.44\textwidth}
%        \centering
%        \includegraphics[width=\linewidth]{./figures_experiments/reports/token_analysis/14b_facts_first_all_hop_raw/category_comparison_simple.png}
%        \caption{\textit{Qwen3-14B}, \texttt{facts\_first}}
%    \end{subfigure}
%
%    \vspace{0.2em}
%
%    \begin{subfigure}[b]{0.44\textwidth}
%        \centering
%        \includegraphics[width=\linewidth]{./figures_experiments/reports/token_analysis/8b_expr_first_all_hop_raw/category_comparison_simple.png}
%        \caption{\textit{Qwen3-8B}, \texttt{expr\_first}}
%    \end{subfigure}
%    \hfill
%    \begin{subfigure}[b]{0.44\textwidth}
%        \centering
%        \includegraphics[width=\linewidth]{./figures_experiments/reports/token_analysis/14b_expr_first_all_hop_raw/category_comparison_simple.png}
%        \caption{\textit{Qwen3-14B}, \texttt{expr\_first}}
%    \end{subfigure}
%    \caption{\textbf{Qwen3 simple token category comparisons.} These pre-refinement summaries already expose the same core ordering as the refined plots. \texttt{facts\_first} is more fact-heavy, whereas \texttt{expr\_first} shifts visible mass toward query-facing and derived-assignment behavior.}
%    \label{fig:qwen_token_simple_category}
%\end{figure*}

%\begin{figure*}[p]
%    \centering
%    \begin{subfigure}[b]{0.44\textwidth}
%        \centering
%        \includegraphics[width=\linewidth]{./figures_experiments/reports/token_analysis/8b_facts_first_all_hop_raw/layer_stage_simple.png}
%        \caption{\textit{Qwen3-8B}, \texttt{facts\_first}}
%    \end{subfigure}
%    \hfill
%    \begin{subfigure}[b]{0.44\textwidth}
%        \centering
%        \includegraphics[width=\linewidth]{./figures_experiments/reports/token_analysis/14b_facts_first_all_hop_raw/layer_stage_simple.png}
%        \caption{\textit{Qwen3-14B}, \texttt{facts\_first}}
%    \end{subfigure}
%
%    \vspace{0.2em}
%
%    \begin{subfigure}[b]{0.44\textwidth}
%        \centering
%        \includegraphics[width=\linewidth]{./figures_experiments/reports/token_analysis/8b_expr_first_all_hop_raw/layer_stage_simple.png}
%        \caption{\textit{Qwen3-8B}, \texttt{expr\_first}}
%    \end{subfigure}
%    \hfill
%    \begin{subfigure}[b]{0.44\textwidth}
%        \centering
%        \includegraphics[width=\linewidth]{./figures_experiments/reports/token_analysis/14b_expr_first_all_hop_raw/layer_stage_simple.png}
%        \caption{\textit{Qwen3-14B}, \texttt{expr\_first}}
%    \end{subfigure}
%    \caption{\textbf{Qwen3 simple token stage summaries.} The stage-pooled simple plots are a lighter-weight version of the raw heatmaps. They show that the earliest stages remain fact-heavy, while the later stages move toward query-facing categories, with prompt reordering mostly changing the balance rather than the coarse stage layout.}
%    \label{fig:qwen_token_simple_stage}
%\end{figure*}

\paragraph{Cross-model region atlases.}
The cross-model region figures are primarily full attention and joint-region galleries for Llama and Mistral. The summary panels already show that both families place more weight on the Query Region than Qwen3; the layer-wise curves show how that difference is realized in depth. In Llama, query-facing mass is broad and diffuse. In Mistral, the stage ordering is easier to recognize, but query-side routing still occupies much more of the back half of the network than it does in Qwen3. Figures~\ref{fig:llama_attn_facts_gallery}, \ref{fig:llama_attn_expr_gallery}, \ref{fig:llama_attn_mlp_facts_gallery}, and \ref{fig:llama_attn_mlp_expr_gallery} collect the retained Llama attention and joint views. Figures~\ref{fig:mistral_attn_facts_gallery}, \ref{fig:mistral_attn_expr_gallery}, \ref{fig:mistral_attn_mlp_facts_gallery}, and \ref{fig:mistral_attn_mlp_expr_gallery} provide the corresponding Mistral views.

\begin{figure*}[t]
    \centering
    \begin{subfigure}[b]{0.285\textwidth}
        \centering
        \includegraphics[width=\linewidth]{./figures_experiments/reports_llama3/attn_analysis_facts_first/attn_one_hop_facts_first/facts_region/attn_facts_region_scores.png}
        \caption{Facts (One-hop)}
    \end{subfigure}
    \hfill
    \begin{subfigure}[b]{0.285\textwidth}
        \centering
        \includegraphics[width=\linewidth]{./figures_experiments/reports_llama3/attn_analysis_facts_first/attn_one_hop_facts_first/expression_region/attn_expression_region_scores.png}
        \caption{Expression (One-hop)}
    \end{subfigure}
    \hfill
    \begin{subfigure}[b]{0.285\textwidth}
        \centering
        \includegraphics[width=\linewidth]{./figures_experiments/reports_llama3/attn_analysis_facts_first/attn_one_hop_facts_first/query_region/attn_query_region_scores.png}
        \caption{Query (One-hop)}
    \end{subfigure}

    \vspace{0.2em}

    \begin{subfigure}[b]{0.285\textwidth}
        \centering
        \includegraphics[width=\linewidth]{./figures_experiments/reports_llama3/attn_analysis_facts_first/attn_two_hop_facts_first/facts_region/attn_facts_region_scores.png}
        \caption{Facts (Two-hop)}
    \end{subfigure}
    \hfill
    \begin{subfigure}[b]{0.285\textwidth}
        \centering
        \includegraphics[width=\linewidth]{./figures_experiments/reports_llama3/attn_analysis_facts_first/attn_two_hop_facts_first/expression_region/attn_expression_region_scores.png}
        \caption{Expression (Two-hop)}
    \end{subfigure}
    \hfill
    \begin{subfigure}[b]{0.285\textwidth}
        \centering
        \includegraphics[width=\linewidth]{./figures_experiments/reports_llama3/attn_analysis_facts_first/attn_two_hop_facts_first/query_region/attn_query_region_scores.png}
        \caption{Query (Two-hop)}
    \end{subfigure}
    \caption{\textbf{Llama attention-region gallery under \texttt{facts\_first}.} Llama's attention branch is strongly query-heavy even without prompt reordering. Facts and expression remain early dominant, but the Query Region keeps unusually large middle and late mass in both one-hop and two-hop reasoning.}
    \label{fig:llama_attn_facts_gallery}
\end{figure*}

\begin{figure*}[t]
    \centering
    \begin{subfigure}[b]{0.31\textwidth}
        \centering
        \includegraphics[width=\linewidth]{./figures_experiments/reports_llama3/attn_analysis_expr_first/attn_one_hop_expr_first/facts_region/attn_facts_region_scores.png}
        \caption{Facts (One-hop)}
    \end{subfigure}
    \hfill
    \begin{subfigure}[b]{0.31\textwidth}
        \centering
        \includegraphics[width=\linewidth]{./figures_experiments/reports_llama3/attn_analysis_expr_first/attn_one_hop_expr_first/expression_region/attn_expression_region_scores.png}
        \caption{Expression (One-hop)}
    \end{subfigure}
    \hfill
    \begin{subfigure}[b]{0.31\textwidth}
        \centering
        \includegraphics[width=\linewidth]{./figures_experiments/reports_llama3/attn_analysis_expr_first/attn_one_hop_expr_first/query_region/attn_query_region_scores.png}
        \caption{Query (One-hop)}
    \end{subfigure}

    \vspace{0.15em}

    \begin{subfigure}[b]{0.31\textwidth}
        \centering
        \includegraphics[width=\linewidth]{./figures_experiments/reports_llama3/attn_analysis_expr_first/attn_two_hop_expr_first/facts_region/attn_facts_region_scores.png}
        \caption{Facts (Two-hop)}
    \end{subfigure}
    \hfill
    \begin{subfigure}[b]{0.31\textwidth}
        \centering
        \includegraphics[width=\linewidth]{./figures_experiments/reports_llama3/attn_analysis_expr_first/attn_two_hop_expr_first/expression_region/attn_expression_region_scores.png}
        \caption{Expression (Two-hop)}
    \end{subfigure}
    \hfill
    \begin{subfigure}[b]{0.31\textwidth}
        \centering
        \includegraphics[width=\linewidth]{./figures_experiments/reports_llama3/attn_analysis_expr_first/attn_two_hop_expr_first/query_region/attn_query_region_scores.png}
        \caption{Query (Two-hop)}
    \end{subfigure}
    \caption{\textbf{Llama attention-region gallery under \texttt{expr\_first}.} Front-loading the expression pushes more mass into early Expression processing, but the Query Region remains broad in both one-hop and two-hop reasoning. The two-hop row further widens the query-facing stage, which makes Llama's diffuse, query-centered routing especially clear.}
    \label{fig:llama_attn_expr_gallery}
\end{figure*}

\begin{figure*}[t]
    \centering
    \begin{subfigure}[b]{0.31\textwidth}
        \centering
        \includegraphics[width=\linewidth]{./figures_experiments/reports_llama3/attn_mlp_analysis_facts_first/attn_mlp_one_hop_facts_first/facts_region/attn_mlp_facts_region_scores.png}
        \caption{Facts (One-hop)}
    \end{subfigure}
    \hfill
    \begin{subfigure}[b]{0.31\textwidth}
        \centering
        \includegraphics[width=\linewidth]{./figures_experiments/reports_llama3/attn_mlp_analysis_facts_first/attn_mlp_one_hop_facts_first/expression_region/attn_mlp_expression_region_scores.png}
        \caption{Expression (One-hop)}
    \end{subfigure}
    \hfill
    \begin{subfigure}[b]{0.31\textwidth}
        \centering
        \includegraphics[width=\linewidth]{./figures_experiments/reports_llama3/attn_mlp_analysis_facts_first/attn_mlp_one_hop_facts_first/query_region/attn_mlp_query_region_scores.png}
        \caption{Query (One-hop)}
    \end{subfigure}

    \vspace{0.15em}

    \begin{subfigure}[b]{0.31\textwidth}
        \centering
        \includegraphics[width=\linewidth]{./figures_experiments/reports_llama3/attn_mlp_analysis_facts_first/attn_mlp_two_hop_facts_first/facts_region/attn_mlp_facts_region_scores.png}
        \caption{Facts (Two-hop)}
    \end{subfigure}
    \hfill
    \begin{subfigure}[b]{0.31\textwidth}
        \centering
        \includegraphics[width=\linewidth]{./figures_experiments/reports_llama3/attn_mlp_analysis_facts_first/attn_mlp_two_hop_facts_first/expression_region/attn_mlp_expression_region_scores.png}
        \caption{Expression (Two-hop)}
    \end{subfigure}
    \hfill
    \begin{subfigure}[b]{0.31\textwidth}
        \centering
        \includegraphics[width=\linewidth]{./figures_experiments/reports_llama3/attn_mlp_analysis_facts_first/attn_mlp_two_hop_facts_first/query_region/attn_mlp_query_region_scores.png}
        \caption{Query (Two-hop)}
    \end{subfigure}
    \caption{\textbf{Llama joint attention+MLP gallery under \texttt{facts\_first}.} The joint branch is somewhat more stage-like than pure attention, but still much more query-heavy than Qwen3. Even in one-hop reasoning the Query Region remains broad across the back half of the network, and the two-hop row broadens this profile further.}
    \label{fig:llama_attn_mlp_facts_gallery}
\end{figure*}

\begin{figure*}[t]
    \centering
    \begin{subfigure}[b]{0.31\textwidth}
        \centering
        \includegraphics[width=\linewidth]{./figures_experiments/reports_llama3/attn_mlp_analysis_expr_first/attn_mlp_one_hop_expr_first/facts_region/attn_mlp_facts_region_scores.png}
        \caption{Facts (One-hop)}
    \end{subfigure}
    \hfill
    \begin{subfigure}[b]{0.31\textwidth}
        \centering
        \includegraphics[width=\linewidth]{./figures_experiments/reports_llama3/attn_mlp_analysis_expr_first/attn_mlp_one_hop_expr_first/expression_region/attn_mlp_expression_region_scores.png}
        \caption{Expression (One-hop)}
    \end{subfigure}
    \hfill
    \begin{subfigure}[b]{0.31\textwidth}
        \centering
        \includegraphics[width=\linewidth]{./figures_experiments/reports_llama3/attn_mlp_analysis_expr_first/attn_mlp_one_hop_expr_first/query_region/attn_mlp_query_region_scores.png}
        \caption{Query (One-hop)}
    \end{subfigure}

    \vspace{0.15em}

    \begin{subfigure}[b]{0.31\textwidth}
        \centering
        \includegraphics[width=\linewidth]{./figures_experiments/reports_llama3/attn_mlp_analysis_expr_first/attn_mlp_two_hop_expr_first/facts_region/attn_mlp_facts_region_scores.png}
        \caption{Facts (Two-hop)}
    \end{subfigure}
    \hfill
    \begin{subfigure}[b]{0.31\textwidth}
        \centering
        \includegraphics[width=\linewidth]{./figures_experiments/reports_llama3/attn_mlp_analysis_expr_first/attn_mlp_two_hop_expr_first/expression_region/attn_mlp_expression_region_scores.png}
        \caption{Expression (Two-hop)}
    \end{subfigure}
    \hfill
    \begin{subfigure}[b]{0.31\textwidth}
        \centering
        \includegraphics[width=\linewidth]{./figures_experiments/reports_llama3/attn_mlp_analysis_expr_first/attn_mlp_two_hop_expr_first/query_region/attn_mlp_query_region_scores.png}
        \caption{Query (Two-hop)}
    \end{subfigure}
    \caption{\textbf{Llama joint attention+MLP gallery under \texttt{expr\_first}.} The expression-first intervention sharpens the early Expression peak, but the Query Region remains broad in both rows. Two-hop reasoning further stretches the query-facing stage across the back half of the network, which keeps the joint branch visibly more diffuse than its Qwen3 counterpart.}
    \label{fig:llama_attn_mlp_expr_gallery}
\end{figure*}

\begin{figure*}[t]
    \centering
    \begin{subfigure}[b]{0.31\textwidth}
        \centering
        \includegraphics[width=\linewidth]{./figures_experiments/reports_mistral/attn_analysis_facts_first/attn_one_hop_facts_first/facts_region/attn_facts_region_scores.png}
        \caption{Facts (One-hop)}
    \end{subfigure}
    \hfill
    \begin{subfigure}[b]{0.31\textwidth}
        \centering
        \includegraphics[width=\linewidth]{./figures_experiments/reports_mistral/attn_analysis_facts_first/attn_one_hop_facts_first/expression_region/attn_expression_region_scores.png}
        \caption{Expression (One-hop)}
    \end{subfigure}
    \hfill
    \begin{subfigure}[b]{0.31\textwidth}
        \centering
        \includegraphics[width=\linewidth]{./figures_experiments/reports_mistral/attn_analysis_facts_first/attn_one_hop_facts_first/query_region/attn_query_region_scores.png}
        \caption{Query (One-hop)}
    \end{subfigure}

    \vspace{0.15em}

    \begin{subfigure}[b]{0.31\textwidth}
        \centering
        \includegraphics[width=\linewidth]{./figures_experiments/reports_mistral/attn_analysis_facts_first/attn_two_hop_facts_first/facts_region/attn_facts_region_scores.png}
        \caption{Facts (Two-hop)}
    \end{subfigure}
    \hfill
    \begin{subfigure}[b]{0.31\textwidth}
        \centering
        \includegraphics[width=\linewidth]{./figures_experiments/reports_mistral/attn_analysis_facts_first/attn_two_hop_facts_first/expression_region/attn_expression_region_scores.png}
        \caption{Expression (Two-hop)}
    \end{subfigure}
    \hfill
    \begin{subfigure}[b]{0.31\textwidth}
        \centering
        \includegraphics[width=\linewidth]{./figures_experiments/reports_mistral/attn_analysis_facts_first/attn_two_hop_facts_first/query_region/attn_query_region_scores.png}
        \caption{Query (Two-hop)}
    \end{subfigure}
    \caption{\textbf{Mistral attention-region gallery under \texttt{facts\_first}.} Mistral retains recognizable early Facts and Expression processing in both rows, but keeps considerably more query mass than Qwen3. The two-hop row broadens this query-facing profile further across the middle and late layers.}
    \label{fig:mistral_attn_facts_gallery}
\end{figure*}

\begin{figure*}[t]
    \centering
    \begin{subfigure}[b]{0.285\textwidth}
        \centering
        \includegraphics[width=\linewidth]{./figures_experiments/reports_mistral/attn_analysis_expr_first/attn_one_hop_expr_first/facts_region/attn_facts_region_scores.png}
        \caption{Facts (One-hop)}
    \end{subfigure}
    \hfill
    \begin{subfigure}[b]{0.285\textwidth}
        \centering
        \includegraphics[width=\linewidth]{./figures_experiments/reports_mistral/attn_analysis_expr_first/attn_one_hop_expr_first/expression_region/attn_expression_region_scores.png}
        \caption{Expression (One-hop)}
    \end{subfigure}
    \hfill
    \begin{subfigure}[b]{0.285\textwidth}
        \centering
        \includegraphics[width=\linewidth]{./figures_experiments/reports_mistral/attn_analysis_expr_first/attn_one_hop_expr_first/query_region/attn_query_region_scores.png}
        \caption{Query (One-hop)}
    \end{subfigure}

    \vspace{0.2em}

    \begin{subfigure}[b]{0.285\textwidth}
        \centering
        \includegraphics[width=\linewidth]{./figures_experiments/reports_mistral/attn_analysis_expr_first/attn_two_hop_expr_first/facts_region/attn_facts_region_scores.png}
        \caption{Facts (Two-hop)}
    \end{subfigure}
    \hfill
    \begin{subfigure}[b]{0.285\textwidth}
        \centering
        \includegraphics[width=\linewidth]{./figures_experiments/reports_mistral/attn_analysis_expr_first/attn_two_hop_expr_first/expression_region/attn_expression_region_scores.png}
        \caption{Expression (Two-hop)}
    \end{subfigure}
    \hfill
    \begin{subfigure}[b]{0.285\textwidth}
        \centering
        \includegraphics[width=\linewidth]{./figures_experiments/reports_mistral/attn_analysis_expr_first/attn_two_hop_expr_first/query_region/attn_query_region_scores.png}
        \caption{Query (Two-hop)}
    \end{subfigure}
    \caption{\textbf{Mistral attention-region gallery under \texttt{expr\_first}.} Prompt reordering shifts more weight into early Expression processing but leaves the Query Region broad and comparatively heavy. The accompanying band summary below confirms that Mistral's attention branch remains highly query-facing even under \texttt{expr\_first}.}
    \label{fig:mistral_attn_expr_gallery}
\end{figure*}

%\begin{figure*}[p]
%    \centering
%    \begin{subfigure}[b]{0.44\textwidth}
%        \centering
%        \includegraphics[width=\linewidth]{./figures_experiments/reports_mistral/attn_analysis_expr_first/comparison/band_metrics_panel.png}
%        \caption{Attention branch}
%    \end{subfigure}
%    \hfill
%    \begin{subfigure}[b]{0.44\textwidth}
%        \centering
%        \includegraphics[width=\linewidth]{./figures_experiments/reports_mistral/attn_mlp_analysis_expr_first/comparison/band_metrics_panel.png}
%        \caption{Joint attention+MLP branch}
%    \end{subfigure}
%    \caption{\textbf{Remaining Mistral expr-first branch summaries.} These two band panels complete the Mistral prompt-order atlas for the non-MLP branches. The attention panel is more middle- and query-heavy; the joint panel is somewhat more stage-like but still broader than its Qwen3 counterpart.}
%    \label{fig:mistral_expr_remaining_band_panels}
%\end{figure*}

%\begin{figure*}[p]
%    \centering
%    \includegraphics[width=\textwidth]{./figures_experiments/reports_mistral/attn_mlp_analysis_facts_first/comparison/band_metrics_panel.png}
%    \caption{\textbf{Mistral facts-first joint attention+MLP band summary.} This final omitted branch-level summary completes the Mistral band-panel atlas. The joint intervention preserves the same broad staged template as the MLP branch, but still places comparatively strong mass on the Query Region throughout the back half of the network.}
%    \label{fig:mistral_attn_mlp_facts_band_panel}
%\end{figure*}

\begin{figure*}[t]
    \centering
    \begin{subfigure}[b]{0.285\textwidth}
        \centering
        \includegraphics[width=\linewidth]{./figures_experiments/reports_mistral/attn_mlp_analysis_facts_first/attn_mlp_one_hop_facts_first/facts_region/attn_mlp_facts_region_scores.png}
        \caption{Facts (One-hop)}
    \end{subfigure}
    \hfill
    \begin{subfigure}[b]{0.285\textwidth}
        \centering
        \includegraphics[width=\linewidth]{./figures_experiments/reports_mistral/attn_mlp_analysis_facts_first/attn_mlp_one_hop_facts_first/expression_region/attn_mlp_expression_region_scores.png}
        \caption{Expression (One-hop)}
    \end{subfigure}
    \hfill
    \begin{subfigure}[b]{0.285\textwidth}
        \centering
        \includegraphics[width=\linewidth]{./figures_experiments/reports_mistral/attn_mlp_analysis_facts_first/attn_mlp_one_hop_facts_first/query_region/attn_mlp_query_region_scores.png}
        \caption{Query (One-hop)}
    \end{subfigure}

    \vspace{0.2em}

    \begin{subfigure}[b]{0.285\textwidth}
        \centering
        \includegraphics[width=\linewidth]{./figures_experiments/reports_mistral/attn_mlp_analysis_facts_first/attn_mlp_two_hop_facts_first/facts_region/attn_mlp_facts_region_scores.png}
        \caption{Facts (Two-hop)}
    \end{subfigure}
    \hfill
    \begin{subfigure}[b]{0.285\textwidth}
        \centering
        \includegraphics[width=\linewidth]{./figures_experiments/reports_mistral/attn_mlp_analysis_facts_first/attn_mlp_two_hop_facts_first/expression_region/attn_mlp_expression_region_scores.png}
        \caption{Expression (Two-hop)}
    \end{subfigure}
    \hfill
    \begin{subfigure}[b]{0.285\textwidth}
        \centering
        \includegraphics[width=\linewidth]{./figures_experiments/reports_mistral/attn_mlp_analysis_facts_first/attn_mlp_two_hop_facts_first/query_region/attn_mlp_query_region_scores.png}
        \caption{Query (Two-hop)}
    \end{subfigure}
    \caption{\textbf{Mistral joint attention+MLP gallery under \texttt{facts\_first}.} The joint branch again looks more stage-like than pure attention, but still preserves a heavy Query Region relative to Qwen3. The full curves show that Mistral's joint intervention sharpens early Facts while leaving a substantial amount of middle-to-late query-facing mass.}
    \label{fig:mistral_attn_mlp_facts_gallery}
\end{figure*}

\begin{figure*}[t]
    \centering
    \begin{subfigure}[b]{0.31\textwidth}
        \centering
        \includegraphics[width=\linewidth]{./figures_experiments/reports_mistral/attn_mlp_analysis_expr_first/attn_mlp_one_hop_expr_first/facts_region/attn_mlp_facts_region_scores.png}
        \caption{Facts (One-hop)}
    \end{subfigure}
    \hfill
    \begin{subfigure}[b]{0.31\textwidth}
        \centering
        \includegraphics[width=\linewidth]{./figures_experiments/reports_mistral/attn_mlp_analysis_expr_first/attn_mlp_one_hop_expr_first/expression_region/attn_mlp_expression_region_scores.png}
        \caption{Expression (One-hop)}
    \end{subfigure}
    \hfill
    \begin{subfigure}[b]{0.31\textwidth}
        \centering
        \includegraphics[width=\linewidth]{./figures_experiments/reports_mistral/attn_mlp_analysis_expr_first/attn_mlp_one_hop_expr_first/query_region/attn_mlp_query_region_scores.png}
        \caption{Query (One-hop)}
    \end{subfigure}

    \vspace{0.15em}

    \begin{subfigure}[b]{0.31\textwidth}
        \centering
        \includegraphics[width=\linewidth]{./figures_experiments/reports_mistral/attn_mlp_analysis_expr_first/attn_mlp_two_hop_expr_first/facts_region/attn_mlp_facts_region_scores.png}
        \caption{Facts (Two-hop)}
    \end{subfigure}
    \hfill
    \begin{subfigure}[b]{0.31\textwidth}
        \centering
        \includegraphics[width=\linewidth]{./figures_experiments/reports_mistral/attn_mlp_analysis_expr_first/attn_mlp_two_hop_expr_first/expression_region/attn_mlp_expression_region_scores.png}
        \caption{Expression (Two-hop)}
    \end{subfigure}
    \hfill
    \begin{subfigure}[b]{0.31\textwidth}
        \centering
        \includegraphics[width=\linewidth]{./figures_experiments/reports_mistral/attn_mlp_analysis_expr_first/attn_mlp_two_hop_expr_first/query_region/attn_mlp_query_region_scores.png}
        \caption{Query (Two-hop)}
    \end{subfigure}
    \caption{\textbf{Mistral joint attention+MLP gallery under \texttt{expr\_first}.} In one-hop reasoning, front-loading the expression amplifies the early Expression Region while keeping the Query Region broad. Two-hop reasoning preserves the same reordered stage structure but broadens the Query Region further, so the joint branch remains visibly less cleanly factorized and more query-facing than its Qwen3 counterpart.}
    \label{fig:mistral_attn_mlp_expr_gallery}
\end{figure*}

%\begin{figure*}[p]
%    \centering
%    \begin{subfigure}[b]{0.285\textwidth}
%        \centering
%        \includegraphics[width=\linewidth]{./figures_experiments/reports_llama3/mlp_analysis_facts_first/mlp_one_hop_facts_first/facts_region/mlp_facts_region_abs_dpd_scores.png}
%        \caption{Llama, facts-first (one-hop)}
%    \end{subfigure}
%    \hfill
%    \begin{subfigure}[b]{0.285\textwidth}
%        \centering
%        \includegraphics[width=\linewidth]{./figures_experiments/reports_llama3/mlp_analysis_facts_first/mlp_two_hop_facts_first/facts_region/mlp_facts_region_abs_dpd_scores.png}
%        \caption{Llama, facts-first (two-hop)}
%    \end{subfigure}
%    \hfill
%    \begin{subfigure}[b]{0.285\textwidth}
%        \centering
%        \includegraphics[width=\linewidth]{./figures_experiments/reports_llama3/mlp_analysis_expr_first/mlp_one_hop_expr_first/facts_region/mlp_facts_region_abs_dpd_scores.png}
%        \caption{Llama, expr-first (one-hop)}
%    \end{subfigure}
%
%    \vspace{0.2em}
%
%    \begin{subfigure}[b]{0.285\textwidth}
%        \centering
%        \includegraphics[width=\linewidth]{./figures_experiments/reports_llama3/mlp_analysis_expr_first/mlp_two_hop_expr_first/facts_region/mlp_facts_region_abs_dpd_scores.png}
%        \caption{Llama, expr-first (two-hop)}
%    \end{subfigure}
%    \hfill
%    \begin{subfigure}[b]{0.285\textwidth}
%        \centering
%        \includegraphics[width=\linewidth]{./figures_experiments/reports_mistral/mlp_analysis_facts_first/mlp_one_hop_facts_first/facts_region/mlp_facts_region_abs_dpd_scores.png}
%        \caption{Mistral, facts-first (one-hop)}
%    \end{subfigure}
%    \hfill
%    \begin{subfigure}[b]{0.285\textwidth}
%        \centering
%        \includegraphics[width=\linewidth]{./figures_experiments/reports_mistral/mlp_analysis_facts_first/mlp_two_hop_facts_first/facts_region/mlp_facts_region_abs_dpd_scores.png}
%        \caption{Mistral, facts-first (two-hop)}
%    \end{subfigure}
%    \caption{\textbf{Cross-model absolute-dPD facts-region diagnostics.} The absolute facts-region curves make it clear that both non-Qwen families still retain a meaningful facts-processing stage in magnitude space. The main difference from Qwen3 is therefore not the absence of facts processing, but its weaker separation from the query-facing stages.}
%    \label{fig:cross_model_abs_facts_region}
%\end{figure*}

%\begin{figure*}[p]
%    \centering
%    \begin{subfigure}[b]{0.285\textwidth}
%        \centering
%        \includegraphics[width=\linewidth]{./figures_experiments/reports_mistral/mlp_analysis_expr_first/mlp_one_hop_expr_first/facts_region/mlp_facts_region_abs_dpd_scores.png}
%        \caption{Mistral, expr-first (one-hop)}
%    \end{subfigure}
%    \hfill
%    \begin{subfigure}[b]{0.285\textwidth}
%        \centering
%        \includegraphics[width=\linewidth]{./figures_experiments/reports_mistral/mlp_analysis_expr_first/mlp_two_hop_expr_first/facts_region/mlp_facts_region_abs_dpd_scores.png}
%        \caption{Mistral, expr-first (two-hop)}
%    \end{subfigure}
%    \hfill
%    \begin{subfigure}[b]{0.285\textwidth}
%        \centering
%        \includegraphics[width=\linewidth]{./figures_experiments/reports_llama3/mlp_analysis_facts_first/mlp_one_hop_facts_first/expression_region/mlp_expression_region_abs_dpd_scores.png}
%        \caption{Llama facts-first: Expression}
%    \end{subfigure}
%
%    \vspace{0.2em}
%
%    \begin{subfigure}[b]{0.285\textwidth}
%        \centering
%        \includegraphics[width=\linewidth]{./figures_experiments/reports_llama3/mlp_analysis_expr_first/mlp_one_hop_expr_first/expression_region/mlp_expression_region_abs_dpd_scores.png}
%        \caption{Llama expr-first: Expression}
%    \end{subfigure}
%    \hfill
%    \begin{subfigure}[b]{0.285\textwidth}
%        \centering
%        \includegraphics[width=\linewidth]{./figures_experiments/reports_mistral/mlp_analysis_facts_first/mlp_one_hop_facts_first/expression_region/mlp_expression_region_abs_dpd_scores.png}
%        \caption{Mistral facts-first: Expression}
%    \end{subfigure}
%    \hfill
%    \begin{subfigure}[b]{0.285\textwidth}
%        \centering
%        \includegraphics[width=\linewidth]{./figures_experiments/reports_mistral/mlp_analysis_expr_first/mlp_one_hop_expr_first/expression_region/mlp_expression_region_abs_dpd_scores.png}
%        \caption{Mistral expr-first: Expression}
%    \end{subfigure}
%    \caption{\textbf{Additional cross-model absolute-dPD diagnostics for facts and expression processing.} These remaining magnitude plots reinforce the same conclusion as the signed and BCR summaries: prompt reordering increases early expression mass in both Llama and Mistral, but neither model develops the same clean fact-versus-query separation seen in Qwen3.}
%    \label{fig:cross_model_abs_remaining}
%\end{figure*}
%
%\begin{figure*}[p]
%    \centering
%    \begin{subfigure}[b]{0.285\textwidth}
%        \centering
%        \includegraphics[width=\linewidth]{./figures_experiments/reports_llama3/mlp_analysis_facts_first/mlp_one_hop_facts_first/query_region/mlp_query_region_abs_dpd_scores.png}
%        \caption{Facts-first: Query (One-hop)}
%    \end{subfigure}
%    \hfill
%    \begin{subfigure}[b]{0.285\textwidth}
%        \centering
%        \includegraphics[width=\linewidth]{./figures_experiments/reports_llama3/mlp_analysis_facts_first/mlp_two_hop_facts_first/expression_region/mlp_expression_region_abs_dpd_scores.png}
%        \caption{Facts-first: Expression (Two-hop)}
%    \end{subfigure}
%    \hfill
%    \begin{subfigure}[b]{0.285\textwidth}
%        \centering
%        \includegraphics[width=\linewidth]{./figures_experiments/reports_llama3/mlp_analysis_facts_first/mlp_two_hop_facts_first/query_region/mlp_query_region_abs_dpd_scores.png}
%        \caption{Facts-first: Query (Two-hop)}
%    \end{subfigure}
%
%    \vspace{0.2em}
%
%    \begin{subfigure}[b]{0.285\textwidth}
%        \centering
%        \includegraphics[width=\linewidth]{./figures_experiments/reports_llama3/mlp_analysis_expr_first/mlp_one_hop_expr_first/query_region/mlp_query_region_abs_dpd_scores.png}
%        \caption{Expr-first: Query (One-hop)}
%    \end{subfigure}
%    \hfill
%    \begin{subfigure}[b]{0.285\textwidth}
%        \centering
%        \includegraphics[width=\linewidth]{./figures_experiments/reports_llama3/mlp_analysis_expr_first/mlp_two_hop_expr_first/expression_region/mlp_expression_region_abs_dpd_scores.png}
%        \caption{Expr-first: Expression (Two-hop)}
%    \end{subfigure}
%    \hfill
%    \begin{subfigure}[b]{0.285\textwidth}
%        \centering
%        \includegraphics[width=\linewidth]{./figures_experiments/reports_llama3/mlp_analysis_expr_first/mlp_two_hop_expr_first/query_region/mlp_query_region_abs_dpd_scores.png}
%        \caption{Expr-first: Query (Two-hop)}
%    \end{subfigure}
%    \caption{\textbf{Residual Llama absolute-dPD diagnostics.} These six remaining magnitude plots complete the Llama MLP atlas. They show that the most persistent residual magnitude sits in the Query Region, especially once the query is moved to the front or when reasoning depth increases to two-hop.}
%    \label{fig:llama_abs_residual_gallery}
%\end{figure*}

%\begin{figure*}[p]
%    \centering
%    \begin{subfigure}[b]{0.285\textwidth}
%        \centering
%        \includegraphics[width=\linewidth]{./figures_experiments/reports_mistral/mlp_analysis_facts_first/mlp_one_hop_facts_first/query_region/mlp_query_region_abs_dpd_scores.png}
%        \caption{Facts-first: Query (One-hop)}
%    \end{subfigure}
%    \hfill
%    \begin{subfigure}[b]{0.285\textwidth}
%        \centering
%        \includegraphics[width=\linewidth]{./figures_experiments/reports_mistral/mlp_analysis_facts_first/mlp_two_hop_facts_first/expression_region/mlp_expression_region_abs_dpd_scores.png}
%        \caption{Facts-first: Expression (Two-hop)}
%    \end{subfigure}
%    \hfill
%    \begin{subfigure}[b]{0.285\textwidth}
%        \centering
%        \includegraphics[width=\linewidth]{./figures_experiments/reports_mistral/mlp_analysis_facts_first/mlp_two_hop_facts_first/query_region/mlp_query_region_abs_dpd_scores.png}
%        \caption{Facts-first: Query (Two-hop)}
%    \end{subfigure}
%
%    \vspace{0.2em}
%
%    \begin{subfigure}[b]{0.285\textwidth}
%        \centering
%        \includegraphics[width=\linewidth]{./figures_experiments/reports_mistral/mlp_analysis_expr_first/mlp_one_hop_expr_first/query_region/mlp_query_region_abs_dpd_scores.png}
%        \caption{Expr-first: Query (One-hop)}
%    \end{subfigure}
%    \hfill
%    \begin{subfigure}[b]{0.285\textwidth}
%        \centering
%        \includegraphics[width=\linewidth]{./figures_experiments/reports_mistral/mlp_analysis_expr_first/mlp_two_hop_expr_first/expression_region/mlp_expression_region_abs_dpd_scores.png}
%        \caption{Expr-first: Expression (Two-hop)}
%    \end{subfigure}
%    \hfill
%    \begin{subfigure}[b]{0.285\textwidth}
%        \centering
%        \includegraphics[width=\linewidth]{./figures_experiments/reports_mistral/mlp_analysis_expr_first/mlp_two_hop_expr_first/query_region/mlp_query_region_abs_dpd_scores.png}
%        \caption{Expr-first: Query (Two-hop)}
%    \end{subfigure}
%    \caption{\textbf{Residual Mistral absolute-dPD diagnostics.} These final magnitude plots complete the Mistral MLP atlas. They reinforce the same picture seen throughout the appendix: Mistral preserves a recognizable expression-processing stage, but its strongest residual magnitude is often query-facing once the task becomes two-hop or the prompt is reordered.}
%    \label{fig:mistral_abs_residual_gallery}
%\end{figure*}

%\begin{figure*}[p]
%    \centering
%    \begin{subfigure}[b]{0.44\textwidth}
%        \centering
%        \includegraphics[width=\linewidth]{./figures_experiments/reports_llama3/token_analysis_facts_first_raw/category_comparison_simple.png}
%        \caption{Llama, \texttt{facts\_first}}
%    \end{subfigure}
%    \hfill
%    \begin{subfigure}[b]{0.44\textwidth}
%        \centering
%        \includegraphics[width=\linewidth]{./figures_experiments/reports_llama3/token_analysis_expr_first_raw/category_comparison_simple.png}
%        \caption{Llama, \texttt{expr\_first}}
%    \end{subfigure}
%
%    \vspace{0.2em}
%
%    \begin{subfigure}[b]{0.44\textwidth}
%        \centering
%        \includegraphics[width=\linewidth]{./figures_experiments/reports_mistral/token_analysis_facts_first_raw/category_comparison_simple.png}
%        \caption{Mistral, \texttt{facts\_first}}
%    \end{subfigure}
%    \hfill
%    \begin{subfigure}[b]{0.44\textwidth}
%        \centering
%        \includegraphics[width=\linewidth]{./figures_experiments/reports_mistral/token_analysis_expr_first_raw/category_comparison_simple.png}
%        \caption{Mistral, \texttt{expr\_first}}
%    \end{subfigure}
%    \caption{\textbf{Cross-model simple token category comparisons.} These non-refined token summaries already show that Llama and Mistral allocate more visible mass to query-facing categories than Qwen3. The refined analyses sharpen the same picture rather than changing it.}
%    \label{fig:cross_model_token_simple_category}
%\end{figure*}

%\begin{figure*}[p]
%    \centering
%    \begin{subfigure}[b]{0.44\textwidth}
%        \centering
%        \includegraphics[width=\linewidth]{./figures_experiments/reports_llama3/token_analysis_facts_first_raw/layer_stage_simple.png}
%        \caption{Llama, \texttt{facts\_first}}
%    \end{subfigure}
%    \hfill
%    \begin{subfigure}[b]{0.44\textwidth}
%        \centering
%        \includegraphics[width=\linewidth]{./figures_experiments/reports_llama3/token_analysis_expr_first_raw/layer_stage_simple.png}
%        \caption{Llama, \texttt{expr\_first}}
%    \end{subfigure}
%
%    \vspace{0.2em}
%
%    \begin{subfigure}[b]{0.44\textwidth}
%        \centering
%        \includegraphics[width=\linewidth]{./figures_experiments/reports_mistral/token_analysis_facts_first_raw/layer_stage_simple.png}
%        \caption{Mistral, \texttt{facts\_first}}
%    \end{subfigure}
%    \hfill
%    \begin{subfigure}[b]{0.44\textwidth}
%        \centering
%        \includegraphics[width=\linewidth]{./figures_experiments/reports_mistral/token_analysis_expr_first_raw/layer_stage_simple.png}
%        \caption{Mistral, \texttt{expr\_first}}
%    \end{subfigure}
%    \caption{\textbf{Cross-model simple token stage summaries.} Stage-pooled raw token plots again show that both non-Qwen models are more query-facing. Llama is the most diffuse, while Mistral preserves a somewhat cleaner early-to-late progression but still places more late-stage weight on query-region categories than Qwen3.}
%    \label{fig:cross_model_token_simple_stage}
%\end{figure*}

\clearpage
\twocolumn

\subsection{Failure Case Analysis}
\label{sec:failure_analysis}

We analyze cases where the model produces incorrect outputs despite exhibiting the same functional head profiles observed in successful reasoning. This analysis reveals that failures are not due to a breakdown in head specialization, but rather to limitations in downstream integration.

\begin{figure}[t]
  \includegraphics[width=\columnwidth]{./figures/corporate_heads.pdf}
	\caption{Schematic illustration of Specialized Attention Heads.
The input sequence is segmented into semantic regions (Clauses and Logical Expression).
Color coding indicates functional roles: Splitting Heads (blue) detect semantic boundaries at delimiters;
Entity-Binding and Transmission Heads (orange) associate variables with values and aggregate information within regions;
Fact-Retrieval Heads (green) access truth values from earlier contexts.}
    \vspace{-0.5em}
	\label{fig:corporate_heads}
\end{figure}

\paragraph{Persistent Head Specialization.}
Even when the model fails on logical reasoning tasks involving excessive reasoning depth or element overload, attention heads continue to exhibit their characteristic functional profiles: Splitting Heads still concentrate on delimiters, Transmission Heads still aggregate within regions, and Fact-Retrieval Heads still attend to truth values. This persistence suggests that the identified mechanisms represent robust architectural properties rather than task-contingent behaviors.

\paragraph{Potential Failure Modes.}
We hypothesize three potential sources of failure:
\begin{itemize}
    \item \textbf{Interaction Deficiency}: Insufficient interaction between heads with complementary roles, preventing effective information handoff.
    \item \textbf{Integration Misalignment}: Misalignment between head outputs and subsequent MLP computations, leading to incorrect logical combinations.
    \item \textbf{Dynamic Rigidity}: The model's inability to dynamically reconfigure attention patterns when faced with novel or compositional reasoning demands.
\end{itemize}

\paragraph{Chain-of-Thought Recovery.}
Notably, in certain failure cases, allowing the model to generate a sufficiently long CoT enables it to ultimately arrive at the correct answer. This suggests that some failures in implicit (non-CoT) reasoning stem from premature convergence or incomplete information aggregation at the query position. Extended sequential processing may compensate by providing additional computational steps for information integration. This observation invites further investigation into the relationship between reasoning depth, head collaboration, and the emergence of correct solutions.



% \paragraph{Other Attention Head Patterns.}
% \label{sec:other_heads}
% Beyond the three primary head types identified in our main analysis, we observe additional attention patterns that merit discussion.
% \textbf{Idle Heads}: As network depth increases, certain heads exhibit minimal functional contribution, consistently attending to the first token regardless of input content. This pattern persists even when the first token is semantically neutral (e.g., a period), suggesting these heads may represent redundant capacity or task-irrelevant components.
% \textbf{Information Binding Heads}: In early layers (L0--10), we observe heads that associate entities with their attributes, binding ``A'' with ``True'' and negation operators with their operands. Given that tokenization separates these semantically coupled elements, such binding appears necessary for downstream logical operations.
% \textbf{Self-Processing Heads}: Some heads exhibit strong diagonal attention patterns, where each token primarily attends to itself. This self-focused processing may facilitate residual-like refinement of token representations without cross-position information mixing.
% \textbf{Expression Processing Heads}: In layers 18--22, immediately preceding attention convergence, we identify heads that concentrate attention within the expression region, potentially executing the core logical computations. The layer positioning (pre-convergence) aligns with our Staged Computation findings, suggesting these heads perform critical reasoning operations before final answer aggregation.
% These auxiliary patterns, while not causally validated through patching experiments, provide complementary evidence for the functional diversity of attention heads in logical reasoning.

\section{Discussion}
\label{sec:appendix_discussion}


\subsection{Other Attention Head Patterns}
\label{sec:appendix_other_heads}

Beyond the three primary head types (\textit{i.e.}, Splitting, Transmission, and Fact-Retrieval) identified in our main analysis, we observe additional attention patterns that provide a more complete picture of the model's internal organization.

\paragraph{Idle Heads.}
As network depth increases, certain heads exhibit minimal functional contribution, consistently attending to the first token regardless of input content. Figure~\ref{fig:other_heads}(a) shows Layer 7 Head 4, where attention concentrates heavily on the first column---all tokens attend primarily to position 0. To verify this is not an artifact of specific token semantics, we tested prompts where the first token is a semantically neutral period (``.'') rather than a meaningful character. The pattern persists. Furthermore, we observed that L7H4 exhibits identical idle behavior in multi-hop reasoning tasks and even in unrelated prompts (e.g., ``The weather is sunny, I am''), suggesting that idle heads may be an intrinsic property of the model architecture rather than task-specific behavior. This raises two possibilities: (1) certain heads are permanently idle across all inputs, representing genuinely redundant capacity; or (2) some heads may be conditionally idle, remaining inactive for certain task types while activating for others. A comprehensive analysis of all idle heads across diverse prompts would be needed to distinguish these hypotheses.

\paragraph{Information Binding Heads.}
In early layers (L0--10), we observe heads that associate entities with their truth values. Specifically, these heads bind ``A'' with ``True'' and ``B'' with ``False'' in the Facts region, and bind negation operators (``$\neg$'') with their operands in the Expression region. Figure~\ref{fig:other_heads}(b-c) illustrates two examples from Layer 0: Head 1 and Head 24 both exhibit diagonal patterns with localized binding, where adjacent tokens within semantic units show elevated mutual attention. Since tokenization separates these semantically coupled elements (e.g., ``$\neg$'' and ``A'' are distinct tokens), such binding appears necessary for downstream logical operations. This binding mechanism complements the Fact-Retrieval heads identified in middle-to-late layers.

\paragraph{Self-Processing Heads.}
Some heads exhibit strong diagonal attention patterns, where each token primarily attends to itself with minimal attention to other positions. Figure~\ref{fig:other_heads}(d-e) shows two examples: Layer 12 Head 20 and Layer 22 Head 12. Both display sharp diagonal lines with near-zero off-diagonal attention, indicating that each token's representation is refined independently. This self-focused processing may facilitate residual-like refinement of token representations without cross-position information mixing. Such heads appear across various layers, suggesting a distributed mechanism for token-level representation enhancement.

\paragraph{Expression Processing Heads.}
In layers 18--22, immediately preceding the attention convergence observed in our main analysis, we identify heads that concentrate attention within the expression region tokens. Figure~\ref{fig:other_heads}(f-h) presents three examples from Layers 18, 19, and 21. These heads show elevated attention in the lower-right region of the attention matrix, corresponding to the expression tokens (``($\neg$A or $\neg$B) is''). Notably, Layer 19 Head 31 additionally attends to the first token, exhibiting functional coupling between expression processing and positional anchoring. These heads appear to process the logical structure of expressions, potentially executing core logical computations. The layer positioning aligns with our Staged Computation findings: expression processing occurs after fact encoding (early layers) but before final answer aggregation (late layers).

\begin{figure}[t]
  \centering
  \includegraphics[width=0.24\columnwidth]{./figures/idle-head-l7h4.jpg}
  \includegraphics[width=0.24\columnwidth]{./figures/info-bind-l0h1.jpg}
  \includegraphics[width=0.24\columnwidth]{./figures/info-bind-l0h24.jpg}
  \includegraphics[width=0.24\columnwidth]{./figures/self-process-l12h20.jpg}
  \\
  \includegraphics[width=0.24\columnwidth]{./figures/self-process-l22h12.jpg}
  \includegraphics[width=0.24\columnwidth]{./figures/exp-process-l18h2.jpg}
  \includegraphics[width=0.24\columnwidth]{./figures/exp-process-l19h31.jpg}
  \includegraphics[width=0.24\columnwidth]{./figures/exp-process-l21h13.jpg}
  \caption{Attention patterns of other head types in \textit{Qwen3-8B}. Top row: (a) Idle Head L7H4---attention concentrates on position 0; (b-c) Information Binding Heads L0H1, L0H24---diagonal patterns with local binding; (d) Self-Processing Head L12H20---sharp diagonal. Bottom row: (e) Self-Processing Head L22H12; (f-h) Expression Processing Heads L18H2, L19H31, L21H13---attention concentrated in expression region (lower-right).}
  \label{fig:other_heads}
\end{figure}

\paragraph{Relationship to Primary Mechanisms.}
These additional patterns complement rather than contradict our main findings. Idle Heads may explain why not all attention capacity is utilized for reasoning. Information Binding Heads provide the early-layer foundation that enables later Fact-Retrieval. Self-Processing Heads may contribute to the residual stream's role in information maintenance. Expression Processing Heads bridge the gap between Information Transmission and final prediction, aligning with the Staged Computation framework. While these patterns have not been causally validated through patching experiments, they suggest functional diversity beyond the three primary head types and warrant further investigation.


\subsection{Functional Coupling in Specialized Heads}
\label{sec:appendix_func_couple}


We provide visual evidence for the functional coupling phenomenon described in the main text. Figure~\ref{fig:func_couple} shows three examples from early layers, each exhibiting multiple functional characteristics simultaneously.

\begin{figure}[t]
  \centering
  \includegraphics[width=0.32\columnwidth]{./figures/func-couple-l0h19.jpg}
  \includegraphics[width=0.32\columnwidth]{./figures/func-couple-l0h31.jpg}
  \includegraphics[width=0.32\columnwidth]{./figures/func-couple-l1h22.jpg}
  \caption{Examples of functional coupling in early-layer heads. Left: L0H19 exhibits staircase-like attention blocks indicating both splitting (segment boundaries) and information binding (within-segment attention). Middle: L0H31 shows diagonal self-processing combined with localized binding blocks. Right: L1H22 displays diagonal attention with scattered cross-position connections, suggesting simultaneous self-processing and selective information retrieval. All three heads demonstrate that individual attention heads participate in multiple functional roles.}
  \label{fig:func_couple}
\end{figure}

These attention patterns reveal that early-layer heads commonly exhibit multi-functional profiles:
\begin{itemize}
    \item \textbf{L0H19}: Staircase-like attention blocks at segment boundaries combined with elevated within-segment attention, performing both splitting and information binding.
    \item \textbf{L0H31}: Diagonal emphasis (self-processing) with localized attention blocks within ``A is True'' and ``B is False'' segments (information binding).
    \item \textbf{L1H22}: Strong diagonal pattern with scattered high-attention points at specific cross-positions, suggesting self-processing combined with selective fact retrieval.
\end{itemize}
This multi-functional profile supports the superposition hypothesis: individual heads participate in multiple overlapping computational circuits, achieving parameter efficiency through functional reuse.


% \subsection{Lazy Reasoning}
% \label{sec:lazy_reasoning}

% While the main analyses characterize information flow patterns, a distinct question concerns where the actual logical computation occurs. We hypothesize a \textit{lazy reasoning} strategy: facts are encoded early and routed through middle layers, but the final logical combination is deferred to late layers. We investigate this hypothesis by analyzing dLD (rather than $|\text{dLD}|$) to preserve the direction of causal effects.

% \paragraph{Experimental Evidence.}
% Figure~\ref{fig:lazy_reasoning} presents mean dLD values (averaged across tokens within each category) for three corruption scenarios, grouped by layer stage. Unlike $|\text{dLD}|$ which measures magnitude only, dLD preserves direction: positive values indicate restoration toward the clean prediction, while negative values indicate interference. The results reveal a differential pattern:

% \begin{figure}[t]
%   \centering
%   \includegraphics[width=\columnwidth]{./figures/info-onehop-multi-no-abs.png}
%   \caption{Preliminary evidence for lazy reasoning hypothesis. Each panel shows mean dLD per token category across layer stages. When predictions differ on corrupted vs. clean inputs (Panel 1, 3), patching the query token in late layers produces positive effects (restoration). When predictions match (Panel 2), patching shows opposite signs: positive in middle layers but negative in late layers (interference). This sign reversal suggests that middle layers transmit information while late layers perform logical computation.}
%   \label{fig:lazy_reasoning}
% \end{figure}

% \paragraph{Restoration Effect.} In Panel 1 (B changes) and Panel 3 (both A and B change), the model predicts differently on corrupted vs. clean inputs. Patching the query token in late layers produces strong positive mean dLD values (+4.8 and +2.4, respectively), indicating successful restoration toward the clean prediction.

% \paragraph{Interference Effect.} Panel 2 presents a control case: the corrupted prompt ``A is False, B is False'' yields the same answer (True) as the clean prompt, and predictions already match (baseline LD = +3.0). Importantly, patching the query token shows opposite signs across layers: positive mean dLD in middle layers but negative mean dLD ($-$0.75) in late layers. This sign reversal suggests that middle layers are still transmitting information, while late layers perform the actual computation where interference occurs.

% \paragraph{Implications.}
% This differential behavior is consistent with the lazy reasoning hypothesis: the model may defer actual logical computation to late layers at the query position, rather than computing incrementally throughout the network. Under this hypothesis, facts are encoded in early layers, routed through middle layers via Information Transmission, with the final logical combination (e.g., evaluating $\neg A \lor \neg B$) occurring only at the last moment.

% The sign reversal in Panel 2---positive dLD in middle layers shifting to negative in late layers---provides preliminary evidence supporting this hypothesis, suggesting that middle layers primarily transmit information while late layers perform the actual computation. Further investigation with additional logical rules and model architectures is needed to validate this finding.


\subsection{Complete Accuracy Tables}
\label{sec:appendix_full_accuracy}
For completeness, we collect the full Step 3 accuracy tables at the end of the appendix. This keeps the main analysis sections readable while preserving the full strategy-by-$k$ results for \textit{Qwen3}, \textit{Llama-3.1-8B-Instruct}, and \textit{Mistral-7B-Instruct-v0.1}.

\begin{table}[p]
\centering
\scriptsize
\renewcommand{\arraystretch}{0.95}
\setlength{\tabcolsep}{3pt}
\resizebox{\columnwidth}{!}{%
\begin{tabular}{@{} l r l l @{}}
\toprule
\textbf{Strategy} & \textbf{k} & \textbf{Accuracy} & \textbf{Drop} \\
\midrule
Fact-Retrieval & 1 & 92.6 & 4.0 \\
Fact-Retrieval & 2 & 89.0 & 7.6 \\
Fact-Retrieval & 4 & 83.0 & 13.6 \\
Fact-Retrieval & 8 & 74.6 & 22.0 \\
Fact-Retrieval & 16 & 67.4 & 29.2 \\
Fact-Retrieval & 32 & 57.0 & 39.6 \\
Fact-Retrieval & 64 & 50.6 & 46.0 \\
Splitting & 1 & 93.8 & 2.8 \\
Splitting & 2 & 92.2 & 4.4 \\
Splitting & 4 & 91.4 & 5.2 \\
Splitting & 8 & 84.6 & 12.0 \\
Splitting & 16 & 83.4 & 13.2 \\
Splitting & 32 & 80.4 & 16.2 \\
Splitting & 64 & 70.4 & 26.2 \\
Transmission & 1 & 91.4 & 5.2 \\
Transmission & 2 & 89.6 & 7.0 \\
Transmission & 4 & 85.8 & 10.8 \\
Transmission & 8 & 80.4 & 16.2 \\
Transmission & 16 & 69.4 & 27.2 \\
Transmission & 32 & 69.2 & 27.4 \\
Transmission & 64 & 62.6 & 34.0 \\
Fact-Retrieval + Splitting & 1 & 92.6 & 4.0 \\
Fact-Retrieval + Splitting & 2 & 90.4 & 6.2 \\
Fact-Retrieval + Splitting & 4 & 82.0 & 14.6 \\
Fact-Retrieval + Splitting & 8 & 73.6 & 23.0 \\
Fact-Retrieval + Splitting & 16 & 63.2 & 33.4 \\
Fact-Retrieval + Splitting & 32 & 57.6 & 39.0 \\
Fact-Retrieval + Splitting & 64 & 53.2 & 43.4 \\
Fact-Retrieval + Transmission & 1 & 91.4 & 5.2 \\
Fact-Retrieval + Transmission & 2 & 89.0 & 7.6 \\
Fact-Retrieval + Transmission & 4 & 83.6 & 13.0 \\
Fact-Retrieval + Transmission & 8 & 74.2 & 22.4 \\
Fact-Retrieval + Transmission & 16 & 63.6 & 33.0 \\
Fact-Retrieval + Transmission & 32 & 57.4 & 39.2 \\
Fact-Retrieval + Transmission & 64 & 50.6 & 46.0 \\
Splitting + Transmission & 1 & 91.4 & 5.2 \\
Splitting + Transmission & 2 & 88.4 & 8.2 \\
Splitting + Transmission & 4 & 85.0 & 11.6 \\
Splitting + Transmission & 8 & 80.4 & 16.2 \\
Splitting + Transmission & 16 & 76.2 & 20.4 \\
Splitting + Transmission & 32 & 69.4 & 27.2 \\
Splitting + Transmission & 64 & 63.4 & 33.2 \\
Mixed-All & 1 & 91.4 & 5.2 \\
Mixed-All & 2 & 89.0 & 7.6 \\
Mixed-All & 4 & 81.8 & 14.8 \\
Mixed-All & 8 & 74.6 & 22.0 \\
Mixed-All & 16 & 68.4 & 28.2 \\
Mixed-All & 32 & 62.0 & 34.6 \\
Mixed-All & 64 & 52.2 & 44.4 \\
Random-Outside & 1 & 96.5 $\pm$ 0.1 & 0.1 $\pm$ 0.1 \\
Random-Outside & 2 & 96.7 $\pm$ 0.1 & -0.1 $\pm$ 0.1 \\
Random-Outside & 4 & 96.8 $\pm$ 0.1 & -0.2 $\pm$ 0.1 \\
Random-Outside & 8 & 97.1 $\pm$ 0.2 & -0.5 $\pm$ 0.2 \\
Random-Outside & 16 & 97.4 $\pm$ 0.2 & -0.8 $\pm$ 0.2 \\
Random-Outside & 32 & 97.2 $\pm$ 0.4 & -0.6 $\pm$ 0.4 \\
Random-Outside & 64 & 91.1 $\pm$ 0.9 & 5.5 $\pm$ 0.9 \\
\bottomrule
\end{tabular}%
}
\caption{\textbf{Full Step 3 accuracy table for \textit{Qwen3-8B} with \texttt{facts\_first} prompts.} All values are post-ablation accuracies and accuracy drops in percentage points. The table reports the complete Qwen3 Step 3 results across all strategies and all $k \in \{1,2,4,8,16,32,64\}$, extending the compact summary in Table~\ref{tab:qwen3_head_accuracy_main}.}
\label{tab:qwen3_8b_ff_full_accuracy}
\end{table}

\begin{table}[p]
\centering
\scriptsize
\renewcommand{\arraystretch}{0.95}
\setlength{\tabcolsep}{3pt}
\resizebox{\columnwidth}{!}{%
\begin{tabular}{@{} l r l l @{}}
\toprule
\textbf{Strategy} & \textbf{k} & \textbf{Accuracy} & \textbf{Drop} \\
\midrule
Fact-Retrieval & 1 & 90.0 & 2.8 \\
Fact-Retrieval & 2 & 87.4 & 5.4 \\
Fact-Retrieval & 4 & 82.6 & 10.2 \\
Fact-Retrieval & 8 & 77.4 & 15.4 \\
Fact-Retrieval & 16 & 70.0 & 22.8 \\
Fact-Retrieval & 32 & 63.2 & 29.6 \\
Fact-Retrieval & 64 & 58.4 & 34.4 \\
Splitting & 1 & 89.2 & 3.6 \\
Splitting & 2 & 86.6 & 6.2 \\
Splitting & 4 & 86.2 & 6.6 \\
Splitting & 8 & 81.2 & 11.6 \\
Splitting & 16 & 78.2 & 14.6 \\
Splitting & 32 & 73.2 & 19.6 \\
Splitting & 64 & 69.2 & 23.6 \\
Transmission & 1 & 89.0 & 3.8 \\
Transmission & 2 & 88.0 & 4.8 \\
Transmission & 4 & 84.0 & 8.8 \\
Transmission & 8 & 78.0 & 14.8 \\
Transmission & 16 & 75.0 & 17.8 \\
Transmission & 32 & 73.2 & 19.6 \\
Transmission & 64 & 60.0 & 32.8 \\
Fact-Retrieval + Splitting & 1 & 90.0 & 2.8 \\
Fact-Retrieval + Splitting & 2 & 86.4 & 6.4 \\
Fact-Retrieval + Splitting & 4 & 80.6 & 12.2 \\
Fact-Retrieval + Splitting & 8 & 77.6 & 15.2 \\
Fact-Retrieval + Splitting & 16 & 72.0 & 20.8 \\
Fact-Retrieval + Splitting & 32 & 64.6 & 28.2 \\
Fact-Retrieval + Splitting & 64 & 57.4 & 35.4 \\
Fact-Retrieval + Transmission & 1 & 90.0 & 2.8 \\
Fact-Retrieval + Transmission & 2 & 86.6 & 6.2 \\
Fact-Retrieval + Transmission & 4 & 82.0 & 10.8 \\
Fact-Retrieval + Transmission & 8 & 77.0 & 15.8 \\
Fact-Retrieval + Transmission & 16 & 71.6 & 21.2 \\
Fact-Retrieval + Transmission & 32 & 61.2 & 31.6 \\
Fact-Retrieval + Transmission & 64 & 57.4 & 35.4 \\
Splitting + Transmission & 1 & 89.0 & 3.8 \\
Splitting + Transmission & 2 & 85.4 & 7.4 \\
Splitting + Transmission & 4 & 83.6 & 9.2 \\
Splitting + Transmission & 8 & 78.4 & 14.4 \\
Splitting + Transmission & 16 & 72.2 & 20.6 \\
Splitting + Transmission & 32 & 71.4 & 21.4 \\
Splitting + Transmission & 64 & 60.4 & 32.4 \\
Mixed-All & 1 & 90.0 & 2.8 \\
Mixed-All & 2 & 86.6 & 6.2 \\
Mixed-All & 4 & 81.2 & 11.6 \\
Mixed-All & 8 & 77.6 & 15.2 \\
Mixed-All & 16 & 71.2 & 21.6 \\
Mixed-All & 32 & 67.4 & 25.4 \\
Mixed-All & 64 & 57.4 & 35.4 \\
Random-Outside & 1 & 93.5 $\pm$ 0.1 & -0.7 $\pm$ 0.1 \\
Random-Outside & 2 & 93.4 $\pm$ 0.2 & -0.6 $\pm$ 0.2 \\
Random-Outside & 4 & 94.1 $\pm$ 0.1 & -1.3 $\pm$ 0.1 \\
Random-Outside & 8 & 94.7 $\pm$ 0.2 & -1.9 $\pm$ 0.2 \\
Random-Outside & 16 & 95.0 $\pm$ 0.3 & -2.2 $\pm$ 0.3 \\
Random-Outside & 32 & 92.5 $\pm$ 0.8 & 0.3 $\pm$ 0.8 \\
Random-Outside & 64 & 86.7 $\pm$ 1.1 & 6.1 $\pm$ 1.1 \\
\bottomrule
\end{tabular}%
}
\caption{\textbf{Full Step 3 accuracy table for \textit{Qwen3-8B} with \texttt{expr\_first} prompts.} All values are post-ablation accuracies and accuracy drops in percentage points. The table reports the complete Qwen3 Step 3 results across all strategies and all $k \in \{1,2,4,8,16,32,64\}$, extending the compact summary in Table~\ref{tab:qwen3_head_accuracy_main}.}
\label{tab:qwen3_8b_ef_full_accuracy}
\end{table}

\begin{table}[p]
\centering
\scriptsize
\renewcommand{\arraystretch}{0.95}
\setlength{\tabcolsep}{3pt}
\resizebox{\columnwidth}{!}{%
\begin{tabular}{@{} l r l l @{}}
\toprule
\textbf{Strategy} & \textbf{k} & \textbf{Accuracy} & \textbf{Drop} \\
\midrule
Fact-Retrieval & 1 & 91.0 & 3.8 \\
Fact-Retrieval & 2 & 89.0 & 5.8 \\
Fact-Retrieval & 4 & 85.0 & 9.8 \\
Fact-Retrieval & 8 & 83.2 & 11.6 \\
Fact-Retrieval & 16 & 76.8 & 18.0 \\
Fact-Retrieval & 32 & 63.8 & 31.0 \\
Fact-Retrieval & 64 & 57.2 & 37.6 \\
Splitting & 1 & 93.4 & 1.4 \\
Splitting & 2 & 92.6 & 2.2 \\
Splitting & 4 & 91.6 & 3.2 \\
Splitting & 8 & 90.4 & 4.4 \\
Splitting & 16 & 88.2 & 6.6 \\
Splitting & 32 & 84.8 & 10.0 \\
Splitting & 64 & 75.2 & 19.6 \\
Transmission & 1 & 92.0 & 2.8 \\
Transmission & 2 & 87.8 & 7.0 \\
Transmission & 4 & 84.4 & 10.4 \\
Transmission & 8 & 76.8 & 18.0 \\
Transmission & 16 & 72.2 & 22.6 \\
Transmission & 32 & 69.2 & 25.6 \\
Transmission & 64 & 57.4 & 37.4 \\
Fact-Retrieval + Splitting & 1 & 91.0 & 3.8 \\
Fact-Retrieval + Splitting & 2 & 88.4 & 6.4 \\
Fact-Retrieval + Splitting & 4 & 84.6 & 10.2 \\
Fact-Retrieval + Splitting & 8 & 82.4 & 12.4 \\
Fact-Retrieval + Splitting & 16 & 77.6 & 17.2 \\
Fact-Retrieval + Splitting & 32 & 66.2 & 28.6 \\
Fact-Retrieval + Splitting & 64 & 57.6 & 37.2 \\
Fact-Retrieval + Transmission & 1 & 91.0 & 3.8 \\
Fact-Retrieval + Transmission & 2 & 88.8 & 6.0 \\
Fact-Retrieval + Transmission & 4 & 85.6 & 9.2 \\
Fact-Retrieval + Transmission & 8 & 79.2 & 15.6 \\
Fact-Retrieval + Transmission & 16 & 67.6 & 27.2 \\
Fact-Retrieval + Transmission & 32 & 60.6 & 34.2 \\
Fact-Retrieval + Transmission & 64 & 53.6 & 41.2 \\
Splitting + Transmission & 1 & 92.0 & 2.8 \\
Splitting + Transmission & 2 & 90.6 & 4.2 \\
Splitting + Transmission & 4 & 83.6 & 11.2 \\
Splitting + Transmission & 8 & 81.2 & 13.6 \\
Splitting + Transmission & 16 & 72.4 & 22.4 \\
Splitting + Transmission & 32 & 68.8 & 26.0 \\
Splitting + Transmission & 64 & 62.4 & 32.4 \\
Mixed-All & 1 & 91.0 & 3.8 \\
Mixed-All & 2 & 88.8 & 6.0 \\
Mixed-All & 4 & 83.8 & 11.0 \\
Mixed-All & 8 & 74.2 & 20.6 \\
Mixed-All & 16 & 65.8 & 29.0 \\
Mixed-All & 32 & 61.4 & 33.4 \\
Mixed-All & 64 & 53.0 & 41.8 \\
Random-Outside & 1 & 95.0 $\pm$ 0.1 & -0.2 $\pm$ 0.1 \\
Random-Outside & 2 & 95.0 $\pm$ 0.1 & -0.2 $\pm$ 0.1 \\
Random-Outside & 4 & 95.5 $\pm$ 0.2 & -0.7 $\pm$ 0.2 \\
Random-Outside & 8 & 96.5 $\pm$ 0.3 & -1.7 $\pm$ 0.3 \\
Random-Outside & 16 & 96.9 $\pm$ 0.2 & -2.1 $\pm$ 0.2 \\
Random-Outside & 32 & 98.3 $\pm$ 0.2 & -3.5 $\pm$ 0.2 \\
Random-Outside & 64 & 97.1 $\pm$ 0.4 & -2.3 $\pm$ 0.4 \\
\bottomrule
\end{tabular}%
}
\caption{\textbf{Full Step 3 accuracy table for \textit{Qwen3-14B} with \texttt{facts\_first} prompts.} All values are post-ablation accuracies and accuracy drops in percentage points. The table reports the complete Qwen3 Step 3 results across all strategies and all $k \in \{1,2,4,8,16,32,64\}$, extending the compact summary in Table~\ref{tab:qwen3_head_accuracy_main}.}
\label{tab:qwen3_14b_ff_full_accuracy}
\end{table}

\begin{table}[p]
\centering
\scriptsize
\renewcommand{\arraystretch}{0.95}
\setlength{\tabcolsep}{3pt}
\resizebox{\columnwidth}{!}{%
\begin{tabular}{@{} l r l l @{}}
\toprule
\textbf{Strategy} & \textbf{k} & \textbf{Accuracy} & \textbf{Drop} \\
\midrule
Fact-Retrieval & 1 & 96.0 & 1.6 \\
Fact-Retrieval & 2 & 93.4 & 4.2 \\
Fact-Retrieval & 4 & 85.4 & 12.2 \\
Fact-Retrieval & 8 & 86.0 & 11.6 \\
Fact-Retrieval & 16 & 81.4 & 16.2 \\
Fact-Retrieval & 32 & 77.2 & 20.4 \\
Fact-Retrieval & 64 & 78.0 & 19.6 \\
Splitting & 1 & 98.0 & -0.4 \\
Splitting & 2 & 98.2 & -0.6 \\
Splitting & 4 & 94.8 & 2.8 \\
Splitting & 8 & 93.4 & 4.2 \\
Splitting & 16 & 90.6 & 7.0 \\
Splitting & 32 & 86.8 & 10.8 \\
Splitting & 64 & 64.0 & 33.6 \\
Transmission & 1 & 95.8 & 1.8 \\
Transmission & 2 & 93.4 & 4.2 \\
Transmission & 4 & 87.0 & 10.6 \\
Transmission & 8 & 79.8 & 17.8 \\
Transmission & 16 & 76.0 & 21.6 \\
Transmission & 32 & 76.6 & 21.0 \\
Transmission & 64 & 69.4 & 28.2 \\
Fact-Retrieval + Splitting & 1 & 96.0 & 1.6 \\
Fact-Retrieval + Splitting & 2 & 95.4 & 2.2 \\
Fact-Retrieval + Splitting & 4 & 89.4 & 8.2 \\
Fact-Retrieval + Splitting & 8 & 89.2 & 8.4 \\
Fact-Retrieval + Splitting & 16 & 87.8 & 9.8 \\
Fact-Retrieval + Splitting & 32 & 75.6 & 22.0 \\
Fact-Retrieval + Splitting & 64 & 63.0 & 34.6 \\
Fact-Retrieval + Transmission & 1 & 95.8 & 1.8 \\
Fact-Retrieval + Transmission & 2 & 92.4 & 5.2 \\
Fact-Retrieval + Transmission & 4 & 90.4 & 7.2 \\
Fact-Retrieval + Transmission & 8 & 84.6 & 13.0 \\
Fact-Retrieval + Transmission & 16 & 80.6 & 17.0 \\
Fact-Retrieval + Transmission & 32 & 67.6 & 30.0 \\
Fact-Retrieval + Transmission & 64 & 60.8 & 36.8 \\
Splitting + Transmission & 1 & 95.8 & 1.8 \\
Splitting + Transmission & 2 & 95.8 & 1.8 \\
Splitting + Transmission & 4 & 91.4 & 6.2 \\
Splitting + Transmission & 8 & 83.2 & 14.4 \\
Splitting + Transmission & 16 & 77.8 & 19.8 \\
Splitting + Transmission & 32 & 71.6 & 26.0 \\
Splitting + Transmission & 64 & 55.6 & 42.0 \\
Mixed-All & 1 & 95.8 & 1.8 \\
Mixed-All & 2 & 92.4 & 5.2 \\
Mixed-All & 4 & 91.0 & 6.6 \\
Mixed-All & 8 & 81.6 & 16.0 \\
Mixed-All & 16 & 76.2 & 21.4 \\
Mixed-All & 32 & 62.0 & 35.6 \\
Mixed-All & 64 & 56.8 & 40.8 \\
Random-Outside & 1 & 97.7 $\pm$ 0.0 & -0.1 $\pm$ 0.0 \\
Random-Outside & 2 & 97.7 $\pm$ 0.1 & -0.1 $\pm$ 0.1 \\
Random-Outside & 4 & 98.0 $\pm$ 0.1 & -0.4 $\pm$ 0.1 \\
Random-Outside & 8 & 98.0 $\pm$ 0.1 & -0.4 $\pm$ 0.1 \\
Random-Outside & 16 & 97.9 $\pm$ 0.1 & -0.3 $\pm$ 0.1 \\
Random-Outside & 32 & 98.2 $\pm$ 0.2 & -0.5 $\pm$ 0.2 \\
Random-Outside & 64 & 98.2 $\pm$ 0.1 & -0.7 $\pm$ 0.1 \\
\bottomrule
\end{tabular}%
}
\caption{\textbf{Full Step 3 accuracy table for \textit{Qwen3-14B} with \texttt{expr\_first} prompts.} All values are post-ablation accuracies and accuracy drops in percentage points. The table reports the complete Qwen3 Step 3 results across all strategies and all $k \in \{1,2,4,8,16,32,64\}$, extending the compact summary in Table~\ref{tab:qwen3_head_accuracy_main}.}
\label{tab:qwen3_14b_ef_full_accuracy}
\end{table}

\begin{table}[p]
\centering
\scriptsize
\renewcommand{\arraystretch}{0.95}
\setlength{\tabcolsep}{3pt}
\resizebox{\columnwidth}{!}{%
\begin{tabular}{@{} l r l l @{}}
\toprule
\textbf{Strategy} & \textbf{k} & \textbf{Accuracy} & \textbf{Drop} \\
\midrule
Fact-Retrieval & 1 & 56.2 & 4.0 \\
Fact-Retrieval & 2 & 55.6 & 4.6 \\
Fact-Retrieval & 4 & 55.2 & 5.0 \\
Fact-Retrieval & 8 & 55.2 & 5.0 \\
Fact-Retrieval & 16 & 55.2 & 5.0 \\
Splitting & 1 & 56.0 & 4.2 \\
Splitting & 2 & 55.0 & 5.2 \\
Splitting & 4 & 55.2 & 5.0 \\
Splitting & 8 & 55.2 & 5.0 \\
Splitting & 16 & 55.4 & 4.8 \\
Transmission & 1 & 56.0 & 4.2 \\
Transmission & 2 & 55.4 & 4.8 \\
Transmission & 4 & 55.2 & 5.0 \\
Transmission & 8 & 55.2 & 5.0 \\
Transmission & 16 & 55.2 & 5.0 \\
Fact-Retrieval + Splitting & 1 & 56.2 & 4.0 \\
Fact-Retrieval + Splitting & 2 & 55.0 & 5.2 \\
Fact-Retrieval + Splitting & 4 & 55.0 & 5.2 \\
Fact-Retrieval + Splitting & 8 & 55.2 & 5.0 \\
Fact-Retrieval + Splitting & 16 & 55.2 & 5.0 \\
Fact-Retrieval + Splitting & 32 & 55.2 & 5.0 \\
Fact-Retrieval + Transmission & 1 & 56.2 & 4.0 \\
Fact-Retrieval + Transmission & 2 & 55.4 & 4.8 \\
Fact-Retrieval + Transmission & 4 & 55.2 & 5.0 \\
Fact-Retrieval + Transmission & 8 & 55.2 & 5.0 \\
Fact-Retrieval + Transmission & 16 & 55.2 & 5.0 \\
Fact-Retrieval + Transmission & 32 & 55.2 & 5.0 \\
Splitting + Transmission & 1 & 56.0 & 4.2 \\
Splitting + Transmission & 2 & 55.4 & 4.8 \\
Splitting + Transmission & 4 & 55.0 & 5.2 \\
Splitting + Transmission & 8 & 55.2 & 5.0 \\
Splitting + Transmission & 16 & 55.2 & 5.0 \\
Splitting + Transmission & 32 & 55.2 & 5.0 \\
Mixed-All & 1 & 56.2 & 4.0 \\
Mixed-All & 2 & 55.0 & 5.2 \\
Mixed-All & 4 & 55.0 & 5.2 \\
Mixed-All & 8 & 55.2 & 5.0 \\
Mixed-All & 16 & 55.2 & 5.0 \\
Mixed-All & 32 & 55.2 & 5.0 \\
Mixed-All & 64 & 55.2 & 5.0 \\
Random-Outside & 1 & 60.2 $\pm$ 0.1 & -0.1 $\pm$ 0.1 \\
Random-Outside & 2 & 60.6 $\pm$ 0.2 & -0.4 $\pm$ 0.2 \\
Random-Outside & 4 & 61.3 $\pm$ 0.2 & -1.1 $\pm$ 0.2 \\
Random-Outside & 8 & 61.4 $\pm$ 0.4 & -1.2 $\pm$ 0.4 \\
Random-Outside & 16 & 62.4 $\pm$ 0.5 & -2.2 $\pm$ 0.5 \\
Random-Outside & 32 & 64.3 $\pm$ 0.7 & -4.1 $\pm$ 0.7 \\
Random-Outside & 64 & 67.4 $\pm$ 0.9 & -7.2 $\pm$ 0.9 \\
\bottomrule
\end{tabular}%
}
\caption{\textbf{Full Step 3 accuracy table for \textit{Llama-3.1-8B-Instruct} with \texttt{facts\_first} prompts.} Base held-out accuracy is $60.2\%$. All values are post-ablation accuracies and accuracy drops in percentage points. Compared with the Qwen3 tables, Llama shows much earlier saturation: many role-targeted ablations collapse to nearly the same accuracy floor, while the Random-Outside control often improves performance rather than preserving baseline.}
\label{tab:llama_ff_full_accuracy}
\end{table}

\begin{table}[p]
\centering
\scriptsize
\renewcommand{\arraystretch}{0.95}
\setlength{\tabcolsep}{3pt}
\resizebox{\columnwidth}{!}{%
\begin{tabular}{@{} l r l l @{}}
\toprule
\textbf{Strategy} & \textbf{k} & \textbf{Accuracy} & \textbf{Drop} \\
\midrule
Fact-Retrieval & 1 & 64.0 & 5.2 \\
Fact-Retrieval & 2 & 61.0 & 8.2 \\
Fact-Retrieval & 4 & 58.6 & 10.6 \\
Fact-Retrieval & 8 & 58.4 & 10.8 \\
Fact-Retrieval & 16 & 58.4 & 10.8 \\
Splitting & 1 & 60.6 & 8.6 \\
Splitting & 2 & 58.4 & 10.8 \\
Splitting & 4 & 58.4 & 10.8 \\
Splitting & 8 & 58.4 & 10.8 \\
Splitting & 16 & 58.4 & 10.8 \\
Splitting & 32 & 58.4 & 10.8 \\
Transmission & 1 & 62.0 & 7.2 \\
Transmission & 2 & 61.6 & 7.6 \\
Transmission & 4 & 59.2 & 10.0 \\
Transmission & 8 & 58.6 & 10.6 \\
Transmission & 16 & 58.4 & 10.8 \\
Fact-Retrieval + Splitting & 1 & 60.6 & 8.6 \\
Fact-Retrieval + Splitting & 2 & 58.8 & 10.4 \\
Fact-Retrieval + Splitting & 4 & 58.4 & 10.8 \\
Fact-Retrieval + Splitting & 8 & 58.4 & 10.8 \\
Fact-Retrieval + Splitting & 16 & 58.4 & 10.8 \\
Fact-Retrieval + Splitting & 32 & 58.4 & 10.8 \\
Fact-Retrieval + Splitting & 64 & 58.4 & 10.8 \\
Fact-Retrieval + Transmission & 1 & 62.0 & 7.2 \\
Fact-Retrieval + Transmission & 2 & 59.4 & 9.8 \\
Fact-Retrieval + Transmission & 4 & 58.6 & 10.6 \\
Fact-Retrieval + Transmission & 8 & 58.4 & 10.8 \\
Fact-Retrieval + Transmission & 16 & 58.4 & 10.8 \\
Fact-Retrieval + Transmission & 32 & 58.4 & 10.8 \\
Splitting + Transmission & 1 & 60.6 & 8.6 \\
Splitting + Transmission & 2 & 58.6 & 10.6 \\
Splitting + Transmission & 4 & 58.4 & 10.8 \\
Splitting + Transmission & 8 & 58.4 & 10.8 \\
Splitting + Transmission & 16 & 58.4 & 10.8 \\
Splitting + Transmission & 32 & 58.4 & 10.8 \\
Splitting + Transmission & 64 & 58.4 & 10.8 \\
Mixed-All & 1 & 60.6 & 8.6 \\
Mixed-All & 2 & 58.6 & 10.6 \\
Mixed-All & 4 & 58.4 & 10.8 \\
Mixed-All & 8 & 58.4 & 10.8 \\
Mixed-All & 16 & 58.4 & 10.8 \\
Mixed-All & 32 & 58.4 & 10.8 \\
Mixed-All & 64 & 58.4 & 10.8 \\
Random-Outside & 1 & 69.3 $\pm$ 0.3 & -0.1 $\pm$ 0.3 \\
Random-Outside & 2 & 69.8 $\pm$ 0.4 & -0.6 $\pm$ 0.4 \\
Random-Outside & 4 & 69.1 $\pm$ 0.3 & 0.1 $\pm$ 0.3 \\
Random-Outside & 8 & 69.8 $\pm$ 0.5 & -0.6 $\pm$ 0.5 \\
Random-Outside & 16 & 70.9 $\pm$ 0.5 & -1.7 $\pm$ 0.5 \\
Random-Outside & 32 & 70.5 $\pm$ 0.9 & -1.3 $\pm$ 0.9 \\
Random-Outside & 64 & 70.9 $\pm$ 0.8 & -1.7 $\pm$ 0.8 \\
\bottomrule
\end{tabular}%
}
\caption{\textbf{Full Step 3 accuracy table for \textit{Llama-3.1-8B-Instruct} with \texttt{expr\_first} prompts.} Base held-out accuracy is $69.2\%$. All values are post-ablation accuracies and accuracy drops in percentage points. The same saturation pattern appears here as in the \texttt{facts\_first} table: once a small number of selected heads are ablated, many strategy variants collapse to the same performance floor, while the Random-Outside control remains near or slightly above baseline.}
\label{tab:llama_ef_full_accuracy}
\end{table}

\begin{table}[p]
\centering
\scriptsize
\renewcommand{\arraystretch}{0.95}
\setlength{\tabcolsep}{3pt}
\resizebox{\columnwidth}{!}{%
\begin{tabular}{@{} l r l l @{}}
\toprule
\textbf{Strategy} & \textbf{k} & \textbf{Accuracy} & \textbf{Drop} \\
\midrule
Fact-Retrieval & 1 & 79.2 & 20.2 \\
Fact-Retrieval & 2 & 53.4 & 46.0 \\
Fact-Retrieval & 4 & 42.6 & 56.8 \\
Fact-Retrieval & 8 & 41.6 & 57.8 \\
Fact-Retrieval & 16 & 41.2 & 58.2 \\
Fact-Retrieval & 32 & 41.2 & 58.2 \\
Fact-Retrieval & 64 & 41.2 & 58.2 \\
Splitting & 1 & 75.0 & 24.4 \\
Splitting & 2 & 67.4 & 32.0 \\
Splitting & 4 & 58.8 & 40.6 \\
Splitting & 8 & 47.0 & 52.4 \\
Splitting & 16 & 42.2 & 57.2 \\
Splitting & 32 & 41.4 & 58.0 \\
Splitting & 64 & 41.2 & 58.2 \\
Transmission & 1 & 88.8 & 10.6 \\
Transmission & 2 & 77.0 & 22.4 \\
Transmission & 4 & 63.0 & 36.4 \\
Transmission & 8 & 45.2 & 54.2 \\
Transmission & 16 & 41.6 & 57.8 \\
Transmission & 32 & 41.2 & 58.2 \\
Transmission & 64 & 41.2 & 58.2 \\
Fact-Retrieval + Splitting & 1 & 79.2 & 20.2 \\
Fact-Retrieval + Splitting & 2 & 53.4 & 46.0 \\
Fact-Retrieval + Splitting & 4 & 42.0 & 57.4 \\
Fact-Retrieval + Splitting & 8 & 41.4 & 58.0 \\
Fact-Retrieval + Splitting & 16 & 41.2 & 58.2 \\
Fact-Retrieval + Splitting & 32 & 41.2 & 58.2 \\
Fact-Retrieval + Splitting & 64 & 41.2 & 58.2 \\
Fact-Retrieval + Transmission & 1 & 79.2 & 20.2 \\
Fact-Retrieval + Transmission & 2 & 64.8 & 34.6 \\
Fact-Retrieval + Transmission & 4 & 43.0 & 56.4 \\
Fact-Retrieval + Transmission & 8 & 41.6 & 57.8 \\
Fact-Retrieval + Transmission & 16 & 41.2 & 58.2 \\
Fact-Retrieval + Transmission & 32 & 41.2 & 58.2 \\
Fact-Retrieval + Transmission & 64 & 41.2 & 58.2 \\
Splitting + Transmission & 1 & 75.0 & 24.4 \\
Splitting + Transmission & 2 & 64.8 & 34.6 \\
Splitting + Transmission & 4 & 51.6 & 47.8 \\
Splitting + Transmission & 8 & 42.8 & 56.6 \\
Splitting + Transmission & 16 & 41.2 & 58.2 \\
Splitting + Transmission & 32 & 41.2 & 58.2 \\
Splitting + Transmission & 64 & 41.2 & 58.2 \\
Mixed-All & 1 & 79.2 & 20.2 \\
Mixed-All & 2 & 53.4 & 46.0 \\
Mixed-All & 4 & 41.8 & 57.6 \\
Mixed-All & 8 & 41.6 & 57.8 \\
Mixed-All & 16 & 41.4 & 58.0 \\
Mixed-All & 32 & 41.2 & 58.2 \\
Mixed-All & 64 & 41.2 & 58.2 \\
Random-Outside & 1 & 98.7 $\pm$ 0.2 & 0.7 $\pm$ 0.2 \\
Random-Outside & 2 & 98.1 $\pm$ 0.3 & 1.3 $\pm$ 0.3 \\
Random-Outside & 4 & 96.5 $\pm$ 0.4 & 2.9 $\pm$ 0.4 \\
Random-Outside & 8 & 95.9 $\pm$ 0.4 & 3.5 $\pm$ 0.4 \\
Random-Outside & 16 & 92.8 $\pm$ 0.7 & 6.6 $\pm$ 0.7 \\
Random-Outside & 32 & 88.4 $\pm$ 1.2 & 11.0 $\pm$ 1.2 \\
Random-Outside & 64 & 79.5 $\pm$ 1.5 & 19.9 $\pm$ 1.5 \\
\bottomrule
\end{tabular}%
}
\caption{\textbf{Full Step 3 accuracy table for \textit{Mistral-7B-Instruct-v0.1} with \texttt{facts\_first} prompts.} Base held-out accuracy is $99.4\%$. All values are post-ablation accuracies and accuracy drops in percentage points. Compared with the Qwen3 tables, Mistral exhibits much stronger overall sensitivity: role-targeted ablations collapse quickly toward the low-40\% range, and even the Random-Outside control causes a substantial degradation.}
\label{tab:mistral_ff_full_accuracy}
\end{table}

\begin{table}[p]
\centering
\scriptsize
\renewcommand{\arraystretch}{0.95}
\setlength{\tabcolsep}{3pt}
\resizebox{\columnwidth}{!}{%
\begin{tabular}{@{} l r l l @{}}
\toprule
\textbf{Strategy} & \textbf{k} & \textbf{Accuracy} & \textbf{Drop} \\
\midrule
Fact-Retrieval & 1 & 93.0 & 6.8 \\
Fact-Retrieval & 2 & 82.2 & 17.6 \\
Fact-Retrieval & 4 & 77.4 & 22.4 \\
Fact-Retrieval & 8 & 65.4 & 34.4 \\
Fact-Retrieval & 16 & 44.6 & 55.2 \\
Fact-Retrieval & 32 & 43.8 & 56.0 \\
Fact-Retrieval & 64 & 43.2 & 56.6 \\
Splitting & 1 & 80.8 & 19.0 \\
Splitting & 2 & 67.2 & 32.6 \\
Splitting & 4 & 62.4 & 37.4 \\
Splitting & 8 & 50.4 & 49.4 \\
Splitting & 16 & 43.6 & 56.2 \\
Splitting & 32 & 43.2 & 56.6 \\
Splitting & 64 & 43.2 & 56.6 \\
Transmission & 1 & 92.0 & 7.8 \\
Transmission & 2 & 85.6 & 14.2 \\
Transmission & 4 & 81.4 & 18.4 \\
Transmission & 8 & 77.6 & 22.2 \\
Transmission & 16 & 56.4 & 43.4 \\
Transmission & 32 & 44.8 & 55.0 \\
Transmission & 64 & 43.2 & 56.6 \\
Fact-Retrieval + Splitting & 1 & 80.8 & 19.0 \\
Fact-Retrieval + Splitting & 2 & 74.8 & 25.0 \\
Fact-Retrieval + Splitting & 4 & 51.4 & 48.4 \\
Fact-Retrieval + Splitting & 8 & 46.4 & 53.4 \\
Fact-Retrieval + Splitting & 16 & 43.2 & 56.6 \\
Fact-Retrieval + Splitting & 32 & 43.2 & 56.6 \\
Fact-Retrieval + Splitting & 64 & 43.2 & 56.6 \\
Fact-Retrieval + Transmission & 1 & 93.0 & 6.8 \\
Fact-Retrieval + Transmission & 2 & 88.0 & 11.8 \\
Fact-Retrieval + Transmission & 4 & 75.0 & 24.8 \\
Fact-Retrieval + Transmission & 8 & 63.0 & 36.8 \\
Fact-Retrieval + Transmission & 16 & 46.4 & 53.4 \\
Fact-Retrieval + Transmission & 32 & 43.4 & 56.4 \\
Fact-Retrieval + Transmission & 64 & 43.2 & 56.6 \\
Splitting + Transmission & 1 & 80.8 & 19.0 \\
Splitting + Transmission & 2 & 79.2 & 20.6 \\
Splitting + Transmission & 4 & 65.0 & 34.8 \\
Splitting + Transmission & 8 & 48.2 & 51.6 \\
Splitting + Transmission & 16 & 44.4 & 55.4 \\
Splitting + Transmission & 32 & 43.2 & 56.6 \\
Splitting + Transmission & 64 & 43.2 & 56.6 \\
Mixed-All & 1 & 80.8 & 19.0 \\
Mixed-All & 2 & 74.8 & 25.0 \\
Mixed-All & 4 & 59.4 & 40.4 \\
Mixed-All & 8 & 47.4 & 52.4 \\
Mixed-All & 16 & 43.4 & 56.4 \\
Mixed-All & 32 & 43.2 & 56.6 \\
Mixed-All & 64 & 43.2 & 56.6 \\
Random-Outside & 1 & 98.7 $\pm$ 0.3 & 1.1 $\pm$ 0.3 \\
Random-Outside & 2 & 97.6 $\pm$ 0.5 & 2.2 $\pm$ 0.5 \\
Random-Outside & 4 & 96.9 $\pm$ 0.5 & 2.9 $\pm$ 0.5 \\
Random-Outside & 8 & 94.5 $\pm$ 0.8 & 5.3 $\pm$ 0.8 \\
Random-Outside & 16 & 90.7 $\pm$ 1.3 & 9.1 $\pm$ 1.3 \\
Random-Outside & 32 & 85.2 $\pm$ 1.3 & 14.5 $\pm$ 1.3 \\
Random-Outside & 64 & 77.6 $\pm$ 1.1 & 22.2 $\pm$ 1.1 \\
\bottomrule
\end{tabular}%
}
\caption{\textbf{Full Step 3 accuracy table for \textit{Mistral-7B-Instruct-v0.1} with \texttt{expr\_first} prompts.} Base held-out accuracy is $99.8\%$. All values are post-ablation accuracies and accuracy drops in percentage points. The same qualitative pattern remains visible as in the \texttt{facts\_first} table: the discovered roles are highly consequential, but the random-outside control is also substantially destructive, indicating a sensitive yet less cleanly isolated head set.}
\label{tab:mistral_ef_full_accuracy}
\end{table}

\end{document}

% Appendix: 1.并行结果讲不完（不同的domain，相同的方法）；2. Detailed（unimportant details，具体实验的参数设置等等，大体的Setting）；

% 语言：1. 尽量用简单的表达（句式 & 词汇）。 表达清楚意思即可，不要增加阅读障碍。
